\documentclass[11pt,oneside]{report}
% ======================================================================
%  THE TRIADIC MEASUREMENT SUBSTRATE --- COMPLETE CORPUS (Volumes 1 & 2)
%  Aaron Switzer, 2026
%  Single-file build. Compile: pdflatex x3 (third pass settles the TOC).
% ======================================================================
% ======================================================================
%  TMS Volume 1 -- shared preamble
%  Compiles with pdfLaTeX both locally and on Overleaf (no exotic fonts).
% ======================================================================
\usepackage[T1]{fontenc}
\usepackage[utf8]{inputenc}
\usepackage{amsmath}
\usepackage{amssymb}
\usepackage{mathptmx}            % Times text + math (widely available)
\usepackage[scaled=0.90]{helvet} % sans for headings
\usepackage{courier}

\usepackage[letterpaper,margin=1in]{geometry}
\usepackage{amsthm}
\usepackage{mathtools}
\usepackage{bm}
\usepackage{microtype}
\usepackage{enumitem}
\usepackage{booktabs}
\usepackage{array}
\usepackage{longtable}
\usepackage{caption}
\usepackage{xcolor}
\usepackage{titlesec}
\usepackage{fancyhdr}
\usepackage[hidelinks]{hyperref}
\usepackage{tocbibind}

% ---- spacing -----------------------------------------------------------
\setlength{\parindent}{1.2em}
\setlength{\parskip}{0.35em}
\linespread{1.04}
\setlist{leftmargin=1.6em,topsep=0.25em,itemsep=0.15em,parsep=0pt}

% ---- theorem environments ---------------------------------------------
\newtheorem{theorem}{Theorem}[section]
\newtheorem{lemma}[theorem]{Lemma}
\newtheorem{corollary}[theorem]{Corollary}
\newtheorem{proposition}[theorem]{Proposition}

\theoremstyle{definition}
\newtheorem{definition}[theorem]{Definition}
\newtheorem{axiom}{Axiom}
\newtheorem{claim}[theorem]{Claim}
\newtheorem{principle}[theorem]{Principle}
\newtheorem{conjecture}[theorem]{Conjecture}
\newtheorem{prediction}[theorem]{Prediction}
\newtheorem{property}[theorem]{Property}
\newtheorem{postulate}[theorem]{Postulate}

\theoremstyle{remark}
\newtheorem{remark}[theorem]{Remark}
\newtheorem*{remark*}{Remark}

% ---- section heading look (unnumbered; author supplies the numbers) ----
\titleformat{\section}[hang]{\normalfont\large\bfseries\sffamily}{}{0pt}{}
\titleformat{\subsection}[hang]{\normalfont\normalsize\bfseries\sffamily}{}{0pt}{}
\titleformat{\subsubsection}[hang]{\normalfont\normalsize\bfseries\itshape}{}{0pt}{}
\titlespacing*{\section}{0pt}{1.4em}{0.5em}
\titlespacing*{\subsection}{0pt}{1.0em}{0.35em}
\titlespacing*{\subsubsection}{0pt}{0.8em}{0.3em}

% chapter (= each paper) heading
\titleformat{\chapter}[display]
  {\normalfont\sffamily\bfseries}
  {}{0pt}{\Huge}
\titlespacing*{\chapter}{0pt}{10pt}{24pt}

% ---- page-break quality: avoid stranded headings and orphaned/widowed lines ----
% (applies to both volumes via the shared preamble)
\widowpenalty=10000        % no single line at the top of a page
\clubpenalty=10000         % no single line at the bottom of a page
\displaywidowpenalty=10000 % same, around displayed equations
\brokenpenalty=4000        % discourage page breaks right after a hyphenated line
% titlesec: forbid a page break immediately after a heading (heading stranded at page foot)
\usepackage{needspace}
\usepackage{etoolbox}
\pretocmd{\section}{\Needspace*{4\baselineskip}}{}{}
\pretocmd{\subsection}{\Needspace*{3\baselineskip}}{}{}

% ---- running heads -----------------------------------------------------
\pagestyle{fancy}
\fancyhf{}
\fancyhead[L]{\small\itshape The Triadic Measurement Substrate}
\fancyhead[R]{\small\itshape\nouppercase{\rightmark}}
\fancyfoot[C]{\small\thepage}
\renewcommand{\headrulewidth}{0.4pt}
\fancypagestyle{plain}{\fancyhf{}\fancyfoot[C]{\small\thepage}\renewcommand{\headrulewidth}{0pt}}

% ---- abstract / keywords / references list ----------------------------
\newenvironment{paperabstract}
  {\begin{center}\begin{minipage}{0.88\textwidth}\small
   \noindent\textbf{Abstract.}\itshape\space}
  {\end{minipage}\end{center}\vspace{0.4em}}

\newcommand{\keywords}[1]{\noindent{\small\textbf{Keywords:} #1}\vspace{0.4em}}

% per-paper APA reference list (hanging indent), kept as authored
\newenvironment{reflist}
  {\par\small\setlength{\parindent}{0pt}%
   \setlength{\leftskip}{1.6em}\setlength{\parindent}{-1.6em}%
   \setlength{\parskip}{0.4em}%
   \raggedright\hyphenpenalty=50\exhyphenpenalty=50\relax
   \sloppy}
  {\par}

% convenience
\providecommand{\tightlist}{\setlength{\itemsep}{0pt}\setlength{\parskip}{0pt}}
\providecommand{\TMS}{\textsc{tms}}
\hyphenation{con-fir-ma-tion sub-strate}
\begin{document}

\thispagestyle{empty}
\begin{titlepage}
\centering
\vspace*{2cm}
{\sffamily\bfseries\Large The Triadic Measurement Substrate\par}
\vspace{0.6cm}
{\large $\bar{B} = 3\bar{\alpha}$\par}
\vspace{1cm}
{\sffamily\large Volume 1\par}
\vspace{1.4cm}
{\Large \bfseries Paper 5: Reinterpretation, Field Emergence, and the Big Bang Correspondence: From Exhausted Substrate to a Universe\par}
\vfill
{\large Aaron Switzer\par}
\vspace{0.4cm}
{\large 2026\par}
\vspace{1.2cm}
{\small\itshape Excerpted from the complete two-volume corpus, \emph{The Triadic Measurement Substrate}. Full corpus and reproducibility package: \url{https://doi.org/10.5281/zenodo.22035422}\par}
\end{titlepage}
\cleardoublepage
\pagenumbering{arabic}

% ---------------- Volume 1 / paper5 ----------------
\chapter*{Paper 5}
\markboth{Paper 5}{Paper 5}
\addcontentsline{toc}{chapter}{Paper 5: Reinterpretation, Field Emergence, and the Big Bang Correspondence: From Exhausted Substrate to a Universe}
\begingroup\centering
{\large\itshape Reinterpretation, Field Emergence, and the Big Bang Correspondence: From Exhausted Substrate to a Universe\par}\vspace{0.8em}
{\normalsize Aaron Switzer\par}
{\small May 2026\par}
\endgroup\vspace{1.0em}

\noindent{\small\textbf{Keywords:} TMS, Reinterpretation Protocol, Recursive Substrate, Big Bang Correspondence, Hierarchical Lightspeeds, Field Emergence, Emergent Gravity, Two-Waves Cosmology}\par\vspace{0.6em}

\section*{Status of Results}
\addcontentsline{toc}{section}{Status of Results}

This paper distinguishes four categories of claim, each labeled where it
appears in the text. The categories continue the labeling conventions
established in Papers 1 through 4.

\textbf{Derived.} Results that follow from the five axioms, the
Principle of Generative Economy, and the structures established in
Papers 1 through 4 by explicit construction or proof. The
reinterpretation operator \ensuremath{\hat{R}} as the unique PGE-selection at \ensuremath{\tau}
\textgreater{} \ensuremath{\tau *(R)} (\S\,2.3); the recursion-preservation theorem stating
that \ensuremath{\hat{R}} produces a substrate structurally identical to the original under
the same five axioms applied at the new scale (\S\,2.5); the maximum
novelty output theorem as bifurcation-forced strict maximization, with
the doubling factor empirically verified by Stage 5 (\S\,2.6, Appendix E);
the wave equation unification theorem under the PGE-forced wave-form
argument of Paper 2 \S\,VII (\S\,3.2); quantization as trough-finding in the
recursive wave hierarchy (\S\,3.5); the sublattice bifurcation at every
hierarchical level k \ensuremath{\geq} 1, with the two-sublattice structure verified
through Stage 6 (\S\,5.3 Claim 1, \S\,6.1, Appendix E); the rigorous
derivation \ensuremath{L^{p(1)}} = \ensuremath{3\sqrt{2}} \ensuremath{\ell_{p}} from the stella-octangula tile plus cardinal
boundary, with the underlying bifurcated structure empirically confirmed
(\S\,6.1.3, Appendix E); the \ensuremath{c^{(k)}} = \ensuremath{(1/\sqrt{2})_{k}} c sequence at saturation as a
structural law of the FCC face-diagonal-to-edge ratio recurring at every
level by scale invariance (\S\,7.2 sharpened, \S\,6.2); gravity as cross-level
dispersion, with \S\,7.2--\S\,7.4 unified as one structural picture (\S\,7.5);
the closed form of f(d, \ensuremath{K_{\psi,\mathrm{local}},} \ensuremath{K_{\psi,\mathrm{surround}},} \ensuremath{K_{\psi,\mathrm{subsystem}},} k)
under additive Kolmogorov bounds (\S\,6.4, Appendix C); the substrate-level
\ensuremath{Z_{2}} \ensuremath{\times} \ensuremath{Z_{2}} stabilizer structure at every hierarchical level (\S\,8.1); the
cross-level commutator mechanism from the sublattice bifurcation as the
substrate-level origin of non-Abelian gauge closure (\S\,8.2); the complete
five-irrep field decomposition under \ensuremath{O_{h}} on the FCC nearest-neighbor
space, exhausting all 12 dimensions (\S\,4.2). The colour gauge group as
su(3) closing in the push-forward transport with no overshoot, and the
global form SU(3) (with unbroken \ensuremath{\mathbb{Z}_{3}} centre realizing confinement) fixed
by the fundamental triality of the agents, Derived under the \S\,4.1
continuum lift (\S\,8.3). The level-1 SU(2) factor and its U(1) trace-mode
partner on the same footing: su(2) closing from the A/B doublet
transport with no overshoot, global form SU(2) (not SO(3)) fixed by the
faithful fundamental doublet, so that the full gauge algebra su(3) \ensuremath{\oplus}
su(2) \ensuremath{\oplus} u(1) is Derived (\S\,8.3); the electroweak global quotient, the
chirality of \ensuremath{SU(2)_{L},} and the fermion generations remain structural
correspondences.

\textbf{Structural correspondence.} Results in which a TMS construction
faithfully reproduces the skeleton of a known framework without claiming
reproduction of every feature. The Big Bang Correspondence (\S\,V) is the
central case: the structural features of the observed cosmological
initial state --- hot dense early phase, near-uniform background with
small fluctuations, large-scale flatness, horizon-connected past --- are
reproduced as consequences of a reinterpretation event acting on a
saturated biographical record \ensuremath{\mathcal{B}(R);} numerical pinning of CMB
temperature, scalar spectral index, and analogous quantitative
observables depends on the specific \ensuremath{\mathcal{B}(R)} and is reserved for simulation.
The matter-antimatter structural correspondence (\S\,5.3 Claim 2): the
two-sublattice \ensuremath{A\leftrightarrow B} exchange at every level k \ensuremath{\geq} 1 is the substrate-level
analog of QFT's matter / antimatter sector decomposition, with the
parallel to Sakharov's three baryogenesis conditions developed
structurally in \S\,5.3 Claim 3. The field emergence (\S\,IV) is a structural
correspondence to quantum field theory in the QCA-derived sense
(Mlodinow \& Brun 2021; Brun \& Mlodinow 2025): the polarization
landscape propagating on the FCC lattice with [[4,2,2]]
stabilizer structure satisfies the structural conditions for QCA-to-QFT
correspondence in the continuum limit. The hierarchical
[[4,2,2]] carrier structure for the SU(2) \ensuremath{\times} U(1) factors (\S\,8.3)
is a structural correspondence: the TMS framework provides the
substrate-level \ensuremath{Z_{2}} \ensuremath{\times} \ensuremath{Z_{2}} stabilizer hierarchy with cross-level commutator
structure, with the explicit Lie-group identification of those factors
reserved for subsequent representation-theoretic work. The colour SU(3)
factor is the exception: as of \S\,8.3 it is Derived under the \S\,4.1 lift
(see the Derived block) --- the push-forward transport closes to su(3)
with no overshoot and the fundamental triality of the agents fixes the
global group as SU(3). The erosion picture as a substrate-level account
of the apparent permanence of standard-model field structure at
human-accessible timescales (\S\,3.5).

\textbf{Predictions / Conjectures.} Falsifiable claims that follow from
the framework but whose closed forms or numerical values are not derived
within this paper. The matter-antimatter asymmetry of the observed
universe as a generic consequence of reinterpretation from a
non-perfectly-symmetric parent record (\S\,5.3 Claim 3), with the three TMS
conditions (sum-mod-4-changing transitions, inversion-asymmetric \ensuremath{\mathcal{B}(R),} \ensuremath{\hat{R}}
phase-transition character) jointly producing the asymmetry. The
hierarchical level identification of our universe --- which level k we
occupy --- is a prediction whose empirical bridge requires comparing
measured cosmological observables against the \ensuremath{c^{(k)}} and \ensuremath{L^{p(k)}} sequences
(\S\,6.3, \S\,9.5). The cosmological-natural-selection signature predicted by
repeated reinterpretation across cycles, in the spirit of Smolin (1992)
and Dowker--Zalel (2017), is a structural prediction (\S\,5.8). The
CMB-anisotropy correspondence to the spectral content of \ensuremath{\mathcal{B}(R)} prior to
reinterpretation (\S\,9.2) is a falsifiable claim whose empirical content
depends on simulation work. The three-generation fermion content from
three-level hierarchical embedding (\S\,8.4) is a structural prediction
whose quantitative consequences (mass spectrum, mixing angles,
CP-violating phases) require detailed simulation. The cross-level
relativistic phenomena are derivable structurally (\S\,VII) but the
explicit functional form recovering the Einstein equations in the
continuum limit is a structural conjecture pending detailed computation
(\S\,7.6).

\textbf{Deferred closures.} Two classes of result remain deferred beyond
Paper 5, by the structural-vs-numerical division of labor across the
volume. First, exact numerical pinning of \ensuremath{L^{p(k)},} \ensuremath{T^{p(k)},} \ensuremath{c^{(k)}} at k \ensuremath{\geq}
2 requires combinatorial enumeration of stabilizer-satisfying T-closed
configurations at each level --- a task whose structural framework Paper
5 specifies (\S\,VI and Appendix B) and whose exhaustive enumeration is
reserved for simulation. Second, the explicit representation-theoretic
mapping from the meta-[[4,2,2]] stabilizer hierarchy with
cross-level commutator structure (\S\,8.2) to the SU(2) \ensuremath{\times} U(1) gauge
structure (\S\,8.3) requires the apparatus of QCA-to-QFT correspondence
developed by Mlodinow, Brun, and the D'Ariano--Perinotti program,
applied to the FCC lattice with hierarchical [[4,2,2]] symmetry;
the colour SU(3) factor is no longer in this deferred class, having been
closed in \S\,8.3 via the push-forward transport.

This four-way labeling applies throughout the paper. Every claim in the
body text falls into exactly one category, and every category's status
is visible at the point of claim.

\section*{Preface to Paper 5}
\addcontentsline{toc}{section}{Preface to Paper 5}

This paper is the fifth and concluding installment in Volume 1 of the
Triadic Measurement Substrate program. Volume 1, Paper 1 established the
geometric and ontological architecture of the substrate. Volume 1, Paper
2 developed the substrate's computational dynamics, the Born-rule
structural correspondence, and the formal program definition. Volume 1,
Paper 3 derived computation from informational substrate via
surround-driven repair, polarization-coherence bistability, recursive
coarse-graining, and the first computing subsystem. Volume 1, Paper 4
developed the macroscopic consequences --- hierarchies, self-modeling,
multi-subsystem stressor dynamics, region biography, and novelty
exhaustion --- and closed at the structural problem: at the exhaustion
threshold of any bounded region, the substrate holds informationally
rich confirmed bits with no remaining direct-read protocol capable of
extracting novelty from them.

Paper 5 inherits that structural problem and closes it. The
reinterpretation protocol that takes the substrate from exhausted \ensuremath{\mathcal{B}(R)}
to a new generative phase is derived here. The fundamental fields, the
initial conditions, and the Big Bang correspondence follow from that
protocol applied at the appropriate scale. As in Papers 1 through 4, no
new axioms are introduced. No new free numerical parameters are
introduced. The Principle of Generative Economy continues to govern
every selection.

Paper 5 is also the numerical-closure paper of the volume. Several
results stated structurally in Papers 1 through 4 are closed here in
their numerical or representational form: the minimum program extent
\ensuremath{L_{p}^{(1)}} at the substrate level, the closed-form f(d, \ensuremath{K_{\psi,\mathrm{local}},}
\ensuremath{K_{\psi,\mathrm{surround}},} \ensuremath{K_{\psi,\mathrm{subsystem}},} k) under additive Kolmogorov bounds,
the hierarchical \ensuremath{c^{(k)}} sequence at the saturation bound, and the
structural correspondence to the gauge symmetries of the Standard Model.
The remaining numerical pinning at k \ensuremath{\geq} 2 is the domain of subsequent
simulation work.

The promise of the volume, stated at the close of Paper 1 --- that the
absolute length scale, the absolute time scale, the propagation speed,
the binary alphabet, the fault-tolerant lattice, the computational
substrate, and the universe inside it follow as geometric consequences
of five axioms and one principle --- is brought to its completion here.
The arc from the sphere primitive of Paper 1 \S\,1.2 to the cosmological
initial conditions of \S\,V is one unbroken geometric chain. The substrate
generates universes inside itself because the geometry of confirmation,
applied recursively across the structural problem of exhaustion, has no
other available move.

\begin{paperabstract}
Building on the exhaustion analysis of Paper 4 and the recursive
coarse-graining of Paper 3, we close the structural loop of Volume 1 by
deriving the reinterpretation protocol that takes the substrate from an
exhausted bounded region to a new generative phase. Eight principal
results are established. First, the reinterpretation operator \ensuremath{\hat{R}} acting
on the biographical record \ensuremath{\mathcal{B}(R)} of an exhausted region is the unique
Principle-of-Generative-Economy selection at \ensuremath{\tau} \textgreater{} \ensuremath{\tau *(R),}
forced by Axioms 3 and 4 jointly closing every other continuation.
Second, the recursion-preservation theorem establishes that \ensuremath{\hat{R}} produces a
substrate \ensuremath{\Xi (R)} structurally identical to the original substrate under
the same five axioms applied at the new scale, with confirmed bits of
\ensuremath{\mathcal{B}(R)} playing the structural role of agents in \ensuremath{\Xi .} Third, the two-waves
problem --- the apparent tension between Paper 1's first-order geometric
wave equation and the second-order wave structure introduced by
reinterpretation --- is resolved by the unification theorem: there is
one wave equation, applied at successive recursion levels, with level-k
providing boundary conditions for level-(k+1). Fourth, the continuum
limit of the polarization field on the FCC lattice with
[[4,2,2]] stabilizer structure satisfies the structural
conditions for QCA-to-QFT correspondence (Mlodinow \& Brun 2021; Brun \&
Mlodinow 2025), yielding fundamental fields with the same skeletal
structure as those of standard QFT. Fifth, a reinterpretation event
acting on a saturated biographical record reproduces the structural
signatures of the observed cosmological initial state: hot dense early
phase, near-uniform background with small fluctuations, large-scale
flatness, and a horizon-connected past --- the Big Bang Correspondence.
Sixth, the closed forms of the structural debts inherited from Papers 1
through 4 are supplied: \ensuremath{L_{p}^{(1)}} = \ensuremath{3\sqrt{2}} \ensuremath{\ell_{p},} the \ensuremath{c^{(k+1)}} \ensuremath{\leq} \ensuremath{(1/\sqrt{2})}
\ensuremath{c^{(k)}} bound at saturation realized by explicit minimum-period
construction, and f(d, \ensuremath{K_{\psi,\mathrm{local}},} \ensuremath{K_{\psi,\mathrm{surround}},} \ensuremath{K_{\psi,\mathrm{subsystem}},} k)
under additive Kolmogorov bounds. Seventh, the cross-level relativistic
phenomena flagged in Paper 4 \S\,8.2 --- variable apparent propagation,
frame-relative timing, embedded-medium dilation --- are derived as the
substrate-level origins of the kinematical and gravitational phenomena
of general relativity. Eighth, the \ensuremath{Z_{2}} \ensuremath{\times} \ensuremath{Z_{2}} stabilizer group of the
substrate-level [[4,2,2]] code, extended hierarchically,
provides candidate carriers for the Standard Model gauge structure SU(3)
\ensuremath{\times} SU(2) \ensuremath{\times} U(1), with the explicit representation-theoretic mapping
reserved as a structural correspondence for subsequent work. As in
Papers 1 through 4, every result is derived from the five axioms and the
Principle of Generative Economy; no new free numerical parameters are
introduced. The model has no numerical degrees of freedom left to tune.
\end{paperabstract}

\section*{I. Inheritance from Papers 1--4}
\addcontentsline{toc}{section}{I. Inheritance from Papers 1--4}

\subsection*{1.1 Inherited Structure}

Paper 5 takes the following structures from Papers 1 through 4 as given
without re-derivation.

\textbf{From Paper 1:} the five axioms (Information Fundamentalism,
Measurement Imperative, Sphere as Primitive, Signal Conservation, Causal
Anchoring) and the Principle of Generative Economy (PGE); the agent
manifold \ensuremath{M\alpha} with agent spacing \ensuremath{\ell_{0}} \ensuremath{\equiv} \ensuremath{\ell_{p};} the FCC causal stack with
face-diagonal nearest-neighbor connectivity; the triadic actualization
rule \ensuremath{\bar{B}} = \ensuremath{3\bar{\alpha}_{2}} in 2D and \ensuremath{\bar{B}} = \ensuremath{4\bar{\alpha}_{3}} in 3D; the \ensuremath{1\to 4} multiplier and
informational closure; the binary alphabet \ensuremath{\bar{B}} \ensuremath{\in} \{0,1\} with polarization
\ensuremath{\chi} = \ensuremath{2\bar{B}} \ensuremath{-} 1; the flop as the atomic unit of state change; the conjugate
field structure \ensuremath{(\varphi ,} \ensuremath{\Pi )} and the wave equation for \ensuremath{\varphi} at \ensuremath{c_{\mathrm{TMS}}} = \ensuremath{1/\sqrt{2};} the
[[4,2,2]] structural correspondence with logical qubit LQ1 (bit
value) and LQ2 (causal-history inheritance); the polarization field \ensuremath{\psi} as
parallel propagating object; the identification \ensuremath{\ell_{0}} \ensuremath{\equiv} \ensuremath{\ell_{p};} the flop time
\ensuremath{\tau_{\mathrm{flop}}} = \ensuremath{t_{p};} and the GUP scaling prediction \ensuremath{\beta} = f(d, \ensuremath{K_{\psi}).}

\textbf{From Paper 2:} the polarization wave equation; the conjugate
pair \ensuremath{(\psi ,} \ensuremath{\Pi_{\psi});} the polarization-biased resolution rule \ensuremath{P(\bar{B}} = 1
\textbar{} \ensuremath{\Pi_{\psi})} = (1 + \ensuremath{\Pi_{\psi})/2} and its Born-rule structural
correspondence; the 3-input majority gate; the universality result
\{MAJ, constant\} \ensuremath{\to} \{AND, OR, NOT\}; the flop operator T; the formal
program definition \ensuremath{(T^{n}[C]} = C with K(C) \textless{}
\textbar C\textbar); and the geometric-selection-implies-K(C)
\textless{} \textbar C\textbar{} implication, ensuring the substrate
does not compute K(C).

\textbf{From Paper 3:} surround-driven repair and the Imprinting Theorem
with imprint depth \ensuremath{d_{R}} = \ensuremath{\sqrt{2}} \ensuremath{\cdot} \ensuremath{L_{R}/\ell_{p};} the dissolution rate and
effective error-rate decay \ensuremath{(1/3)^{d};} the per-event error rate \ensuremath{p_{\mathrm{error}}}
= (1 \ensuremath{-} \textbar \ensuremath{\Pi_{\psi}}\textbar)/2; the coherence bistability threshold
\textbar \ensuremath{\Pi_{\psi}}\textbar\ensuremath{*(L_{R});} the coarse-graining operator G with axiom
scale-independence; the meta-level triadic rule, meta-level FCC
structure, and meta-level [[4,2,2]] stabilizer dynamics; the
minimum program extent \ensuremath{L_{p}} (structural placeholder); the \ensuremath{(C_{D},} \ensuremath{C_{O},}
\ensuremath{C_{C})} trichotomy and the minimum subsystem extent \ensuremath{L_{\mathrm{sub}}} = \ensuremath{3L_{p}} + \ensuremath{\delta ;}
the information-generation rate differential; and the hierarchical
lightspeed formula \ensuremath{c^{(k)}} = \ensuremath{\prod _{j=1k}} \ensuremath{(L_{p}^{(j)}} / \ensuremath{L_{p}^{(j-1)})} \ensuremath{\cdot}
\ensuremath{(T_{p}^{(j-1)}} / \ensuremath{T_{p}^{(j)})} \ensuremath{\cdot} c.

\textbf{From Paper 4:} the recursive coarse-graining operator \ensuremath{G^{(k)}}
at arbitrary hierarchical depth; the scale-separation requirements
\ensuremath{L_{p}^{(k)}} \ensuremath{\gg} \ensuremath{L_{p}^{(k-1)}} and \ensuremath{T_{p}^{(k)}} \ensuremath{\gg} \ensuremath{T_{p}^{(k-1)};} the
hierarchy ceiling \ensuremath{k_{\mathrm{max}};} the Russian-doll floor at \ensuremath{L_{p};} the three
growth modes (annexation, imprinting, incorporation); self-modeling as a
structural property of every subsystem at k \ensuremath{\geq} 1, making every subsystem
a measurement node; meta-meta-confirmation as the same triadic structure
applied at the next level; the three stressors (mortality, predation,
starvation); the Borel-Cantelli inevitability of subsystem formation;
the finiteness of \ensuremath{\Sigma (R),} the monotonic decay of \ensuremath{\nu (R,} \ensuremath{\tau ),} the exhaustion
threshold \ensuremath{\tau}\emph{(R) at which \ensuremath{\nu (R,} \ensuremath{\tau )} \textbf{\textless{}} \ensuremath{\nu}}(R); the
region biography \ensuremath{\mathcal{B}(R)} and its three phases; the upper bound \ensuremath{c^{(k+1)}} \ensuremath{\leq}
\ensuremath{(1/\sqrt{2})} \ensuremath{\cdot} \ensuremath{c^{(k)};} and the structural problem: at \ensuremath{\tau} \textgreater{}
\ensuremath{\tau *(R),} the substrate holds informationally rich confirmed bits with no
remaining direct-read protocol capable of extracting novelty, and Axioms
3 and 4 jointly force reinterpretation as the only available
continuation.

These structures constitute the entirety of the foundation. Paper 5
introduces no further primitives.

\subsection*{1.2 The Questions Paper 5 Answers}

Paper 4 \S\,6.6 stated the structural problem: an exhausted region holds
confirmed bits whose direct-read protocol can extract no further
novelty, but the bits cannot be lost (Axiom 3), cannot be overwritten
(Axiom 4), and must continue to be hosted by the substrate. The only
structurally available move is reinterpretation under a different
protocol --- one that treats the confirmed bits not as configurations to
evaluate for stability, but as something else. Paper 4 left the
identification of that protocol and its consequences to Paper 5.

Paper 5 answers six questions. First, what is the reinterpretation
operator \ensuremath{\hat{R},} and why is it the unique PGE-selection at \ensuremath{\tau} \textgreater{}
\ensuremath{\tau *(R)?} Second, what new substrate \ensuremath{\Xi (R)} does \ensuremath{\hat{R}} produce, and what is its
relationship to the original substrate? Third, how does \ensuremath{\hat{R}} resolve the
apparent tension between the first-order geometric wave equation of
Paper 1 \S\,VI and the second-order wave structure that reinterpretation
must produce? Fourth, what are the fundamental fields of the new
substrate, and how do they relate to those of standard QFT? Fifth, what
are the initial conditions of the new substrate, and do they correspond
to the observed Big Bang signatures? Sixth, what numerical closures of
Papers 1--4 --- \ensuremath{L_{p}^{(k)},} \ensuremath{T_{p}^{(k)},} \ensuremath{c^{(k)},} f, and the
gauge-structure correspondence --- does Paper 5 supply?

Beyond the six committed questions, Paper 5 derives additional
structural results: the cross-level relativistic phenomena as the
substrate-level origins of gravitational dynamics, the Maximum Novelty
Output statement as a derived law at the reinterpretation boundary, and
the structural foundation for cosmological natural selection across
repeated reinterpretation cycles.

\section*{II. The Reinterpretation Protocol}
\addcontentsline{toc}{section}{II. The Reinterpretation Protocol}

\subsection*{2.1 The Structural Problem Restated}

At \ensuremath{\tau} \textgreater{} \ensuremath{\tau *(R)} for some bounded region R, three conditions
hold jointly:

\begin{itemize}
\item
  \ensuremath{\Sigma_{\mathrm{sampled}}(R,} \ensuremath{\tau )} \ensuremath{\approx} \ensuremath{\Sigma (R):} the accessible configuration space has been
  exhausted by direct sampling. Further sampling under the substrate's
  standard direct-read protocol extracts no novel content.
\item
  Every confirmed bit in R remains causally anchored (Axiom 4) and the
  signals that produced them remain conserved in the record (Axiom 3).
  The bits cannot be lost, overwritten, or deferred.
\item
  The substrate must continue operating. By the architecture of the TMS
  --- flop-by-flop confirmation as the atomic unit of state change ---
  there is no mechanism by which any region of the substrate can halt or
  be suspended.
\end{itemize}

These three conditions are jointly inconsistent under the standard
direct-read protocol. The first says no novelty can be extracted; the
third says operation must continue; the second forbids destructive
resolution. The substrate is in a structurally constrained state
requiring a continuation that is neither destructive nor trivially
repetitive.

\subsection*{2.2 Candidate Continuations and Their Elimination}

By the Principle of Generative Economy, the substrate selects the
minimum-sufficient continuation that remains generatively non-trivial.
We enumerate the candidate continuations and eliminate all but one.

\textbf{Candidate 1: Direct re-sampling.} The substrate continues to
draw configurations from \ensuremath{\Sigma (R)} under the direct-read protocol. By the
definition of exhaustion, every configuration drawn is one already
sampled. The substrate produces no new compressible structure; the
region becomes a museum of repetitions. This satisfies (i) the
minimum-specification condition of PGE --- no new mechanism is invoked
--- but fails (ii), the generative non-triviality condition: every
continuation is a repetition of an already-extracted state. Direct
re-sampling is a generative dead-end. \emph{Eliminated by PGE-ii.}

\textbf{Candidate 2: Loss of confirmed bits.} The substrate dissolves
the exhausted region back into \ensuremath{U_{\mathrm{sh}},} returning its confirmed bits to the
latent state. This violates Axiom 3 (signal conservation): the signals
that produced the confirmed bits have been consumed and locked into the
record, and they cannot be retrieved by reversing the confirmation. The
constraint set K explicitly includes Axiom 3, so this candidate is not
in \ensuremath{\mathcal{C}_{K}.} \emph{Eliminated by axiom violation.}

\textbf{Candidate 3: Overwriting confirmed bits.} The substrate writes
new content into the positions held by confirmed bits, replacing the
existing values. This violates Axiom 4 (causal anchoring) directly.
\emph{Eliminated by axiom violation.}

\textbf{Candidate 4: Lossy compression.} The substrate retains only a
summary of the exhausted region, discarding the bit-level detail. This
violates Axiom 3 in the same sense as Candidate 2 --- the discarded
content corresponds to signals that have been consumed and are part of
the record. \emph{Eliminated by axiom violation.}

\textbf{Candidate 5: External offloading.} The substrate transfers the
exhausted region's content to a separate storage substrate distinct from
the parent TMS substrate. By Axiom 0, there is only one substrate; \ensuremath{U_{\mathrm{sh}}}
is the unique ground state. There is no ``elsewhere'' to which content
can be transferred. \emph{Eliminated by axiom violation.}

\textbf{Candidate 6: Reinterpretation.} The substrate re-reads the
confirmed bits of \ensuremath{\mathcal{B}(R)} under a protocol that does not draw them as
samples of \ensuremath{\Sigma (R)} but treats them as a new primitive layer --- agents in a
new substrate \ensuremath{\Xi (R)} operating at a coarser scale. The confirmed bits are
preserved (Axioms 3, 4 satisfied); the operation continues (the new
substrate generates fresh configurations at its scale); the generative
content is non-trivial (the new substrate has its own \ensuremath{\Sigma (\Xi (R))} which is
not yet exhausted). All three constraints of \S\,2.1 are satisfied. This
candidate is in \ensuremath{\mathcal{C}_{K}} and satisfies both (i) and (ii) of PGE.

Among the six candidates, only Candidate 6 satisfies the constraint set
K and both PGE conditions. By the Principle of Generative Economy as
formally stated in Paper 1 \S\,I, the unique PGE-selection at \ensuremath{\tau}
\textgreater{} \ensuremath{\tau *(R)} is reinterpretation.

\subsection*{2.3 The Reinterpretation Operator \ensuremath{\hat{R}}}

\begin{definition}[Reinterpretation Operator]
Let \ensuremath{\mathcal{B}(R)} denote the
biographical record of an exhausted region R at time \ensuremath{\tau} \ensuremath{\geq} \ensuremath{\tau *(R).} The
reinterpretation operator \ensuremath{\hat{R}} is the mapping
\end{definition}

\begin{equation*}
\hat{R} : \mathcal{B}(R) \to \Xi (R)
\end{equation*}

specified by the following identifications:

\begin{itemize}
\item
  Each confirmed bit \ensuremath{\bar{B}} \ensuremath{\in} \ensuremath{\mathcal{B}(R)} becomes a new agent \ensuremath{\tilde{\alpha}} in the substrate
  \ensuremath{\Xi (R).} The bit's positional anchor in the FCC lattice of the parent
  substrate becomes the new agent's positional anchor in the lattice of
  \ensuremath{\Xi (R).} The bit's binary value \ensuremath{\bar{B}} \ensuremath{\in} \{0, 1\} is preserved as part of the
  new agent's identity but does not affect the new agent's structural
  role --- new agents, like all agents under Axiom 2, are immutable
  substrate.
\end{itemize}

\begin{itemize}
\item
  The polarization landscape \ensuremath{\psi} across \ensuremath{\mathcal{B}(R)} becomes the new
  maximum-entropy field \ensuremath{U_{\mathrm{sh}}'} of \ensuremath{\Xi (R).} The polarization values \ensuremath{\chi} \ensuremath{\in} \{\ensuremath{-1,}
  +1\} carried by the new agents constitute the latent informational
  content from which new bits will be actualized at the \ensuremath{\Xi -\text{scale}.}
\end{itemize}

\begin{itemize}
\item
  The biographical record \ensuremath{\mathcal{B}(R)} viewed as a whole becomes the structural
  pre-substrate of \ensuremath{\Xi (R):} the set of immutable positional anchors and
  their polarization values defines the analog of the agent manifold \ensuremath{M\alpha}
  at the new scale.
\end{itemize}

\begin{itemize}
\item
  The five axioms apply afresh to \ensuremath{\Xi (R)} at the new scale. Triadic
  confirmation, the \ensuremath{1\to 4} multiplier, FCC structure, the [[4,2,2]]
  stabilizers, and the wave equations are all reproduced at the new
  scale by the same arguments that produced them at the substrate level
  --- by axiom scale-independence as established in Paper 3 \S\,IV and
  \S\,4.7.
\end{itemize}

\ensuremath{\hat{R}} is not a new axiom and not a new free parameter. It is the geometric
identification of which features of \ensuremath{\mathcal{B}(R)} play which structural roles in
\ensuremath{\Xi (R).} The structural roles themselves --- agents, latent field, lattice
connectivity --- are inherited from the axioms unchanged.

\subsection*{2.4 The New Agent Identification: Confirmed Bits as Substrate}

The central structural identification of \ensuremath{\hat{R}} is that confirmed bits of the
parent substrate play the role of agents in \ensuremath{\Xi (R).} This identification is
forced by three properties of confirmed bits, each established in
earlier papers:

\textbf{Immutability.} By Axiom 4 (causal anchoring), confirmed bits
cannot be overwritten or modified once locked into the causal record.
This is the same immutability property that defines agents under Axiom 2
--- agents do not change, move, or degrade. Confirmed bits and agents
share the immutability property by axiom.

\textbf{Positional definiteness.} Confirmed bits occupy definite
positions in the FCC lattice --- specifically, the tetrahedral
interstice positions whose centroids lie on the next-order T-void
sublattice (Paper 1 \S\,II; Paper 3 \S\,IV). These positions form a regular
sublattice with well-defined nearest-neighbor relations under the
substrate's connectivity. Agents at the substrate level occupy lattice
positions with similar regularity --- the FCC vertex sublattice in
particular. Both classes of object are positionally definite primitives
of their respective lattices.

\textbf{Substrate role under PGE.} By the Principle of Generative
Economy as applied in Paper 3 \S\,IV, the meta-level substrate inherits its
primitives from objects in the parent substrate that satisfy the
substrate role: immutable, isotropic, and capable of supporting the
triadic confirmation pattern at the new scale. Confirmed bits satisfy
all three. They are immutable by Axiom 4; they are isotropic in the
sense that no preferred orientation is encoded in their polarization
values; and they are positioned at the centroids of the next-order
T-void sublattice in a configuration that supports tetrahedral closure
at the new scale.

The T-void sublattice property is geometrically decisive. As established
in Paper 1 \S\,II.4 and refined in Paper 3 \S\,IV, the centroids of
tetrahedral interstices in an FCC lattice form a simple cubic
sublattice. Tetrahedral closure at the new scale --- the confirmation
pattern required by the triadic minimum at any scale --- produces
centroids that fall on positions of the next-order T-void sublattice.
The geometric recursion is exact: each level's confirmed-bit positions
are the next level's lattice positions.

\begin{theorem}[Confirmed Bits as Agents]
Under \ensuremath{\hat{R},} every confirmed
bit \ensuremath{\bar{B}} in \ensuremath{\mathcal{B}(R)} satisfies the structural conditions of Axiom 2 in \ensuremath{\Xi (R):} it
is an immutable, isotropic primitive at a definite position in a lattice
supporting tetrahedral closure. The identification ``confirmed bit
becomes new agent'' is forced by the axioms, not chosen.
\end{theorem}

\begin{proof}
Axiom 2 in \ensuremath{\Xi (R)} requires that the new agents be immutable
(sphere primitive in the sense that they do not change, move, or
degrade); isotropic (zero orientation parameters); and arranged in a
lattice supporting the triadic confirmation pattern. Confirmed bits in
\ensuremath{\mathcal{B}(R)} satisfy immutability by Axiom 4 applied at the parent scale. They
satisfy isotropy because their binary value is a scalar polarization
with no orientation content. They satisfy
lattice-supporting-triadic-closure by the T-void sublattice geometry
derived in Paper 3 \S\,IV: the centroids of tetrahedral interstices form a
simple cubic sublattice on which the same FCC closure pattern reappears
at the next scale. The PGE selection criterion (minimum sufficient,
generatively non-trivial) is satisfied because no other primitive in
\ensuremath{\mathcal{B}(R)} shares all three properties --- the polarization values are not
positionally distinct primitives, the touch-vector signals have been
consumed and are not stored independently of the bits they produced, and
the substrate's other structures (latent bits in \ensuremath{U_{\mathrm{sh}},} agents at the
parent scale) belong to the parent substrate rather than to \ensuremath{\mathcal{B}(R)} as the
substrate of \ensuremath{\Xi .}

\end{proof}

\subsection*{2.5 The Recursion-Preservation Theorem}

\begin{theorem}[Recursion Preservation under \ensuremath{\hat{R}}]
Under \ensuremath{\hat{R},} the five
axioms and the Principle of Generative Economy produce in \ensuremath{\Xi (R)} the same
triadic confirmation, FCC connectivity, [[4,2,2]] stabilizer
dynamics, conjugate field structure \ensuremath{(\varphi ,} \ensuremath{\Pi )} and \ensuremath{(\psi ,} \ensuremath{\Pi_{\psi}),} and wave
equations as they produced in the parent substrate, at the new scale.
\end{theorem}

\begin{proof}
By Paper 3 \S\,4.7 (axiom scale-independence), the five
axioms do not refer to any particular scale; their content applies
wherever the structural conditions they specify are met. Under \ensuremath{\hat{R},} the
structural conditions are met at the scale \ensuremath{L_{p}^{(1)}} of the minimum
program: \ensuremath{\mathcal{B}(R)} supplies isotropic immutable primitives (new agents =
confirmed bits) on a lattice supporting tetrahedral closure (T-void
sublattice = next-order simple cubic). Axiom 0 supplies the new
maximum-entropy field \ensuremath{U_{\mathrm{sh}}'} = the polarization landscape across \ensuremath{\mathcal{B}(R).}
Axiom 1 supplies the measurement imperative at the new scale. Axiom 2
supplies the sphere primitive (now embodied in confirmed bits). Axioms 3
and 4 supply signal conservation and causal anchoring at the new scale,
inherited automatically. The Principle of Generative Economy supplies
the same selection criterion at the new scale. The same logical chain of
derivations that produced triadic confirmation at the substrate level
(Paper 1 \S\,II), the \ensuremath{1\to 4} multiplier (Paper 1 \S\,IV), FCC necessity (Paper 1
\S\,II.4), the [[4,2,2]] stabilizer correspondence (Paper 1 \S\,V),
and the wave equations (Paper 1 \S\,VI; Paper 2 \S\,II) reapplies unchanged at
the new scale. The new propagation speed in level-1 units is \ensuremath{c_{\mathrm{TMS}}} = \ensuremath{1/\sqrt{2}}
(level-1 spacings per level-1 flop), inheriting the face-diagonal
argument of Paper 1 \S\,6.5 unchanged. The structural reproduction is exact
at the level of the axioms.

\end{proof}

\begin{corollary}[No New Axioms, No New Parameters]
\ensuremath{\hat{R}} introduces no
new axioms and no new free numerical parameters. The new substrate \ensuremath{\Xi (R)}
is governed by the same five axioms and the same Principle of Generative
Economy as the parent substrate. Every quantity in \ensuremath{\Xi (R)} is determined at
the new scale by the same axiomatic chain that determined the parent's
quantities at the parent scale.
\end{corollary}

\subsection*{2.6 The Maximum Novelty Output Theorem (Bifurcation-Forced Maximization)}

The informal observation that the substrate appears to maximize novel
information output --- flagged across previous papers but not previously
stated as a derived law --- acquires its formal status here at the
reinterpretation boundary. The maximization is non-trivial: it follows
from the sublattice bifurcation of the level-1 substrate (\S\,6.1), which
doubles the agent count per level relative to any single-sublattice
alternative continuation.

\begin{theorem}[Maximum Novelty Output at Reinterpretation]
\emph{At
the reinterpretation boundary \ensuremath{\tau} = \ensuremath{\tau}}(R), the substrate's continuation
rate \ensuremath{\nu '(R,} \ensuremath{\tau}\emph{) under \ensuremath{\hat{R}} is strictly greater than \ensuremath{\nu '} under any
alternative continuation that preserves Axioms 3 and 4. The geometric
origin of this strict inequality is the sublattice bifurcation: at each
hierarchical level k \ensuremath{\geq} 1, the FCC substrate splits into two
interpenetrating FCC sublattices A and B, doubling the agent count per
spatial site (one A-flavor agent and one B-flavor agent at each
level-(k+1) shared centroid), compared to any hypothetical continuation
that preserves only one sublattice.}
\end{theorem}

\begin{proof}
Three steps.

\emph{(i) The sublattice bifurcation is a structural feature of \ensuremath{\hat{R}.}} By
\S\,6.1, the level-1 substrate produced by \ensuremath{\hat{R}} is sublattice-bifurcated: two
interpenetrating FCC sublattices A and B, related by spatial inversion
through the level-2 centroids, both running [[4,2,2]] stabilizer
dynamics independently. By Stages 5 and 6 (Appendix E), this structure
is confirmed empirically and is the unique geometric configuration in
which the level-1 FCC + [[4,2,2]] + tetrahedral closure
conditions are jointly satisfied.

\emph{(ii) The bifurcation doubles the agent count per level relative to
single-sublattice alternatives.} Each level-(k+1) spatial site hosts one
A-flavor agent and one B-flavor agent, derived from the
inversion-partner A-tetrahedron and B-tetrahedron at level k. A
hypothetical continuation that produces only one sublattice (e.g., only
A-flavor agents) would have agent density at level k+1 exactly half that
of the bifurcated structure. Stage 5 confirms this: the total
tetrahedral count and the total null-space dimension at \ensuremath{L_{\mathrm{box}}} scale
with both sublattices present, doubling the corresponding
single-sublattice value.

\emph{(iii) Higher agent density permits strictly more novelty.} The
continuation rate \ensuremath{\nu '} at a reinterpretation boundary is bounded above by
the per-flop sampling capacity of the new substrate \ensuremath{\Xi (R).} By Paper 4
\S\,6.2, \ensuremath{\nu (R,} \ensuremath{\tau )} \ensuremath{\approx} \ensuremath{\rho_{\mathrm{sample}}(R)} \ensuremath{\cdot} (\textbar \ensuremath{\Sigma (R)}\textbar{} \ensuremath{-}
\textbar \ensuremath{\Sigma_{\mathrm{sampled}}(R,} \ensuremath{\tau )}\textbar) / \textbar \ensuremath{\Sigma (R)}\textbar; at the
moment of reinterpretation \textbar \ensuremath{\Sigma_{\mathrm{sampled}}}\textbar{} /
\textbar \ensuremath{\Sigma}\textbar{} \ensuremath{\approx} 0, so \ensuremath{\nu '} is set by \ensuremath{\rho_{\mathrm{sample}}} alone. By Paper 3
\S\,3.5 (continuous-sampling mechanism), \ensuremath{\rho_{\mathrm{sample}}} scales linearly with
agent density: each agent participates in confirmation events at rate
\textasciitilde\ensuremath{1/T^{p(k+1)}.} The bifurcated structure gives \ensuremath{2\times} the agent
count compared to single-sublattice alternatives, so \ensuremath{\nu '} under \ensuremath{\hat{R}} exceeds
\ensuremath{\nu '} under any single-sublattice alternative by precisely a factor of 2 at
the per-level capacity bound.

Combining (i)--(iii): \ensuremath{\hat{R}} is the unique PGE-selection at \ensuremath{\tau} \textgreater{}
\ensuremath{\tau *(R)} by \S\,2.3, and its bifurcated output strictly dominates any
alternative continuation that preserves Axioms 3 and 4 at the per-level
capacity bound. The maximum is strict, not by design, but by the
structural forcing of FCC + [[4,2,2]] + tetrahedral closure
under PGE-selection at the recursion boundary.

\end{proof}

\textbf{Status:} \emph{Derived.} The strict maximization is forced by
the bifurcated structure of \S\,6.1, which is itself empirically verified
through Stage 6 (Appendix E).

\textbf{Remark on the non-triviality.} The prior form of this statement
(in earlier drafts) was effectively trivial because it followed from the
uniqueness of \ensuremath{\hat{R}} in \S\,2.3. The sublattice-bifurcated version is
non-trivial because the bifurcation gives the strict-inequality content:
alternative continuations exist that produce a new substrate from \ensuremath{\mathcal{B}(R)}
without violating Axioms 3 or 4 (e.g., a continuation that produces only
one of the two sublattices), but each such alternative has exactly half
the agent count of \ensuremath{\hat{R}`s} bifurcated output, and therefore strictly lower
\ensuremath{\nu '} at saturation. The substrate does not compute or evaluate novelty as
a criterion. Strict novelty maximization is a structural consequence of
the sublattice bifurcation, which is itself a structural consequence of
the FCC + [[4,2,2]] + recursive coarse-graining under PGE.

\textbf{Remark on the bits-per-agent ratio.} In bulk FCC, the
bits-to-agents ratio is 2 (each agent participates in 4 tetrahedra, each
tetrahedron contains 4 agents, giving 4/4 \ensuremath{\times} bits-per-tet, with two bits
per inversion-partner pair through the shared centroid). The bifurcated
structure preserves this ratio at every level: per spatial site, two
flavors of agent contribute jointly to one level-(k+1) bit. This is the
bits-per-agent invariant that anchors the maximum-output theorem's
quantitative content.

\subsection*{2.7 The Reinterpretation Event as Phase Transition}

Reinterpretation is not a continuous process. It is a structural phase
transition that occurs at \ensuremath{\tau} = \ensuremath{\tau *(R)} and produces \ensuremath{\Xi (R)} as a new
substrate. Three properties of the transition follow directly from the
construction:

\textbf{Atomicity at the level of R.} \ensuremath{\hat{R}} applied to \ensuremath{\mathcal{B}(R)} is a single
mapping; \ensuremath{\mathcal{B}(R)} is reinterpreted in its entirety, not piece by piece. The
transition is atomic at the regional scale: either R is reinterpreted or
it is not.

\textbf{Locality at the substrate scale.} \ensuremath{\hat{R}} affects only \ensuremath{\mathcal{B}(R).} Regions
outside R that are not exhausted continue to operate under the
direct-read protocol. The substrate is not globally reinterpreted; each
bounded region undergoes reinterpretation independently when it
exhausts.

\textbf{Generative resetting at the new scale.} Within \ensuremath{\Xi (R),} the new
direct-read protocol begins with \ensuremath{\Sigma_{\mathrm{sampled}}} = \ensuremath{\varnothing} and a freshly populated
\ensuremath{\Sigma (\Xi (R)).} The exhaustion clock starts over at the new scale. This is the
structural origin of the cosmological initial state in \S\,V.

The phase-transition framing is structurally important for the Big Bang
Correspondence of \S\,V: a reinterpretation event, viewed from within \ensuremath{\Xi (R),}
looks like the sudden onset of a new dynamical phase from a maximally
informationally dense initial state --- the structural signature of a
Big Bang.

\section*{III. The Two-Waves Problem and Its Resolution}
\addcontentsline{toc}{section}{III. The Two-Waves Problem and Its Resolution}

\subsection*{3.1 The Apparent Tension}

Paper 1 \S\,VI derives the wave equation

\begin{equation*}
\partial ^{2}\varphi / \partial t^{2} - c^{2}\nabla ^{2}\varphi = 0
\end{equation*}

from the discrete dynamics of the confirmation field on the FCC lattice.
This is a first-order geometric wave equation: it operates on the
confirmation density at the substrate level and propagates information
across the lattice at \ensuremath{c_{\mathrm{TMS}}} = \ensuremath{1/\sqrt{2}} lattice spacings per flop. Its
propagation is forced by the FCC second-moment tensor \ensuremath{\Sigma_{j}} \ensuremath{r_{j}\mu} \ensuremath{r_{j}\nu} = \ensuremath{4\ell^{2}_{p}}
\ensuremath{\delta \mu \nu} over the 12 nearest neighbors --- a property of the substrate-level
lattice.

Reinterpretation produces a new substrate \ensuremath{\Xi (R)} with its own confirmation
field \ensuremath{\varphi '} and polarization field \ensuremath{\psi ',} and (by the Recursion-Preservation
Theorem of \S\,2.5) its own wave equations of the same form, at the new
scale. This appears to introduce a second wave equation operating on the
new substrate. An external observer might ask: are there now two wave
equations operating simultaneously? Does the parent substrate's \ensuremath{\varphi}
continue to propagate alongside the new substrate's \ensuremath{\varphi '?} If so, do they
interfere? If not, what happens to the parent \ensuremath{\varphi} when reinterpretation
occurs?

This is the two-waves problem. It is an apparent tension between Paper
1's substrate-level wave dynamics and Paper 5's reinterpretation-induced
wave dynamics. We resolve it here as a single wave structure expressed
at successive recursion levels.

\subsection*{3.2 The Unification Theorem (under the PGE-Forced Wave-Form Argument)}

Paper 2 \S\,VII established that any PGE-selected continuation at a
reinterpretation boundary takes wave form, because (i) the substrate's
only native propagation vocabulary is the wave equation of Paper 1 \S\,VI,
(ii) PGE forbids importing new mechanisms, and (iii) wave propagation is
form-invariant under recursive coarse-graining. The recursion is
therefore a wave hierarchy: each level's dynamics use the same
propagation verb, expressed at the level-k scale.

The two-waves tension dissolves under this result. There is one wave
equation, applied at successive recursion levels. The apparent ``two
waves'' are the same equation expressed at the parent scale and the \ensuremath{\Xi}
scale, related by the \ensuremath{c^{(k+1)}} \ensuremath{\leq} \ensuremath{(1/\sqrt{2})} \ensuremath{c^{(k)}} coupling-timescale constraint
of Paper 4 \S\,5.3.

\begin{theorem}[Wave Equation Unification]
There is one wave
equation. It applies at every level of the recursion as the structurally
forced expression of the substrate's native propagation vocabulary.
Level-k provides boundary conditions for level-(k+1); together they form
a single hierarchical wave structure.
\end{theorem}

\begin{proof}
By Paper 2 \S\,VII, the level-(k+1) dynamics are wave-form by
PGE-forced vocabulary-reuse. By the Recursion-Preservation Theorem of
\S\,2.5, the geometric substrate at level k+1 reproduces the
substrate-level lattice structure (FCC connectivity, face-diagonal
nearest-neighbors, the \ensuremath{4\ell^{2}_{p}} \ensuremath{\delta} second-moment tensor) at the new scale. By
Paper 1 \S\,6.6, these three structural properties uniquely determine the
wave equation in the continuum limit. Therefore the level-(k+1) wave
equation has the same form as the level-k wave equation; only the
absolute scale changes. The parent substrate's \ensuremath{\varphi} does not stop
propagating when \ensuremath{\mathcal{B}(R)} is reinterpreted --- it has already done its
propagation, and \ensuremath{\mathcal{B}(R)} is the static record of where it ended up. The new
substrate's \ensuremath{\varphi '} propagates within \ensuremath{\Xi (R)} at its own scale. The two are not
in competition; they are one wave equation at two recursion levels.

\end{proof}

\subsection*{3.3 Boundary Conditions: Level-k as Source for Level-(k+1)}

The unification requires specifying how the level-k record provides
boundary conditions for the level-(k+1) wave dynamics. The mechanism is
the same one that operates at single levels: the polarization landscape
of \ensuremath{\mathcal{B}(R),} now identified as \ensuremath{U_{\mathrm{sh}}'} of \ensuremath{\Xi (R),} supplies the latent
informational content from which level-1 bits are actualized. The
level-1 confirmation events draw from this content; the resulting
level-1 polarization field carries the imprint of the parent-level
polarization landscape forward.

Formally, the level-(k+1) wave equation has initial conditions

\begin{equation*}
\psi '(r, \tau ' = 0) = G[\psi_{\mathcal{B}}(R)] ; \Pi_{\psi}'(r, \tau ' = 0) =
G[\Pi_{\psi,\mathcal{B}}(R)]
\end{equation*}

where G is the coarse-graining operator of Paper 3 \S\,IV applied to the
parent-level polarization data at the moment of reinterpretation. The
level-(k+1) dynamics then propagate forward in level-(k+1) flop time,
with their initial conditions inherited from the parent's terminal
state.

This is the structural origin of the matching condition between
cosmological epochs: the late phase of a parent substrate (the terminal
state of \ensuremath{\mathcal{B}(R)} just before reinterpretation) becomes the early phase of
the daughter substrate (the initial state of \ensuremath{\Xi (R)} just after
reinterpretation). No information is lost across the boundary --- Axioms
3 and 4 ensure that \ensuremath{\mathcal{B}(R)'s} content is preserved and carried into \ensuremath{\Xi (R)} as
initial polarization landscape. No new information is introduced --- the
daughter substrate's initial dynamics are determined entirely by the
parent's terminal state under the coarse-graining operator.

\subsection*{3.4 The Hierarchical Wave Structure}

Iterated across multiple reinterpretation events, the wave structure
becomes a hierarchy of wave equations at successive recursion levels,
each providing boundary conditions for the next. At any given
hierarchical depth k, the wave equation governs propagation in level-k
units; the absolute propagation speed in substrate units is \ensuremath{c^{(k)};}
the lattice spacing is \ensuremath{L_{p}^{(k)};} the flop period is \ensuremath{T_{p}^{(k)}.} The
hierarchy is generated by the chain of reinterpretation events, not by a
recursive design.

This resolves a structural question that has been latent across the
volume: how does the substrate handle wave dynamics across hierarchical
levels? The answer is the same wave equation, expressed at the operative
level, with cross-level coupling provided by the \ensuremath{c^{(k)}} dispersion and
the coarse-graining boundary conditions.

Cross-level interactions --- when a \ensuremath{\text{level}-(k-1)} subsystem operates
inside a level-k meta-subsystem --- produce the relativistic phenomena
treated in \S\,VII. The mechanism is the boundary condition between the two
levels' wave dynamics meeting at the interface.

\subsection*{3.5 Quantization as Trough-Finding in the Hierarchical Wave Structure}

The wave hierarchy of \S\,3.2--\S\,3.4 forces a specific structural mechanism
for stable excitations and thereby for quantization. We develop this
mechanism here as the substrate-level account of how the QCA-to-QFT
correspondence of \S\,IV produces a discrete spectrum of field content
rather than a continuum.

A wave hierarchy has closure constraints at each level: at level k,
propagating packets are confirmed-bit content settling under the level-k
confirmation rule, with the level-k wave equation as the underlying
propagation. A packet remains stable at level k only if the level-k wave
geometry permits standing-wave closure within the level-k T-period ---
that is, if there exists a self-consistent configuration in which the
packet reproduces itself after \ensuremath{T^{p(k)}} at the level-k spatial scale
\ensuremath{L^{p(k)}.} The configurations where this closure is permitted are the
\emph{wave troughs} of the hierarchical structure: low-curvature regions
of the accumulated wake where the substrate's effective topography
supports persistent recursion.

\begin{proposition}[Quantization-by-Closure]
The discrete
spectrum of stable excitations at level k is exactly the set of
self-consistent standing-wave configurations on the level-k substrate,
with their amplitudes set by the integer number of confirmed bits each
configuration carries. Quantization is not a separate axiom; it is the
structurally forced consequence of having a recursive wave hierarchy
with confirmation closure at each level.
\end{proposition}

\begin{proof}
By \S\,3.4, propagation at level k is governed by the
level-k wave equation with periodic confirmation events at \ensuremath{T^{p(k)}.} A
localized excitation persists across many \ensuremath{T^{p(k)}} only if it returns to
a stabilizer-satisfying configuration after each period; the set of such
configurations is discrete, indexed by the level-k FCC's representation
theory under \ensuremath{O_{h}} (the decomposition developed for the substrate level
in \S\,4.2 and applied recursively here). Amplitudes are integer counts
because each contributing confirmed bit is a single 0/1 actualization.
The combination of discrete configuration set and integer-valued
amplitude is the substrate's quantization.

\end{proof}

\emph{Status:} Derived. The exact spectrum at each k is the level-k
application of the \S\,4.2 field decomposition under the same \ensuremath{O_{h}} irrep
machinery as the substrate level --- a structural consequence of
recursive scale-invariance, not a separate derivation.

\textbf{The erosion picture.} The hierarchical wave structure is,
equivalently, an accumulating wake of all prior propagation: every
confirmation event leaves a residue in the substrate's effective
topography, and new excitations ride atop the topography that prior
events have already carved. The substrate's effective field landscape at
any cosmological epoch is a superposition of three timescale layers:

\begin{itemize}
\item
  \emph{Ancient layer:} deep channels carved by very early dynamics
  during the first reinterpretation cycles. These are effectively rigid
  at any subsystem-accessible timescale, and constitute what
  subsystem-level observers experience as the fixed field structure of
  their universe (the load-bearing field content of standard particle
  physics is in this layer).
\end{itemize}

\begin{itemize}
\item
  \emph{Intermediate layer:} ongoing dynamics across the current
  cosmological epoch. These evolve on timescales comparable to subsystem
  lifetimes and are the dynamical field content available for direct
  observation and experiment.
\end{itemize}

\begin{itemize}
\item
  \emph{Transient layer:} flash-flood channels carved by recent events
  that have not yet relaxed. These appear as short-lived field
  excitations, particle resonances, and the like.
\end{itemize}

The differential picture (a packet's current dynamics carving its own
wake) and the integrated picture (the universe's accumulated history as
field content) are the same mechanism at different timescales.

\textbf{This gives a substrate-level account of field stability without
positing it.} The appearance that field structure is fixed at
human-accessible timescales is an \emph{erosion property} of the
hierarchical wave structure: deep channels carved by \textasciitilde\ensuremath{10^{10}}
years of accumulated propagation are effectively rigid for
laboratory-scale subsystems, not because field stability is fundamental,
but because the carving timescale dwarfs the subsystem's observational
timescale. The fundamental dynamics remain the wave hierarchy of \S\,3.2;
field stability is a derived consequence of accumulated wake.

\section*{IV. Field Emergence: From Confirmed Bits to Fundamental Fields}
\addcontentsline{toc}{section}{IV. Field Emergence: From Confirmed Bits to Fundamental Fields}

\subsection*{4.1 The QCA-to-QFT Correspondence at the FCC Lattice}

The discrete TMS substrate dynamics produce continuum quantum field
theory in a well-defined limit. We state the lift as a
substrate-internal theorem with two supporting lemmas. Earlier drafts of
this section imported the lift from external QCA-to-QFT programs
(Mlodinow \& Brun, 2021; Brun \& Mlodinow, 2025; D'Ariano \& Perinotti,
2014); the version given here routes through TMS primitives alone, with
the external programs retained as comparison citations in \S\,7.7 rather
than as load-bearing imports.

\begin{theorem}[QCA-to-QFT Lift]
Under the recursion structure
of \S\,2.5 and the level-advance relation of \S\,6.3.0, the level-1 substrate
dynamics produce a quantum field theory in the continuum limit with the
following content: (i) fermionic spinor fields from the half-integer spin
forced by the triadic oddness of \ensuremath{\bar{B}} = \ensuremath{3\bar{\alpha}} (Lemma A); (ii)
bosonic gauge fields and scalars from coarse-grained polarization
decomposition (\S\,4.2); (iii) approximate Lorentz invariance with specific
residual scaling (Lemma B); (iv) unitary evolution from substrate
coherence (\S\,3.2); (v) Born-rule projective measurement from \S\,III's
confirmation primitive (Paper 2 \S\,IV).
\end{theorem}

\textbf{Lemma A (Spinor Character from the Triadic Oddness).} \emph{The
confirmed bit carries half-integer spin because the confirmation event is
triadic. Each of the three agents of \ensuremath{\bar{B}} = \ensuremath{3\bar{\alpha}} carries a two-valued
polarization, i.e.~a spin-\ensuremath{\tfrac{1}{2}} degree of freedom; the recoupling of an
odd number of spin-\ensuremath{\tfrac{1}{2}} objects necessarily carries half-integer total
spin, and the confirmed bit inherits this half-integer character,
propagated up the recursion undiminished by the threeness. A single
triadic node is therefore fermionic; a bound pair of nodes is bosonic.}

\emph{The mechanism.} The founding equation \ensuremath{\bar{B}} = \ensuremath{3\bar{\alpha}} has the
structure of a trivalent recoupling node: three agents meet and close
around one bit, the trivalent-node-with-closure structure of a Penrose
spin network, in which combinatorial spin labels on the edges of a
network reproduce SU(2) angular momentum in the large-network limit with
\emph{no background space presupposed}. The TMS fermion is this
spin-network node.

Three facts fix the half-integer character. First, each agent \emph{is}
a spin-\ensuremath{\tfrac{1}{2}} object by identity: the substrate's own resolution rule
\ensuremath{P(+)} = \ensuremath{(1 + \Pi)/2} (Paper 2 \S\,IV), which gives the probability that a
confirmation reads \ensuremath{+} given continuous polarization bias
\ensuremath{\Pi \in [-1, 1]}, equals \ensuremath{\cos^2(\theta/2)} under \ensuremath{\Pi} = \ensuremath{\cos\theta} --- the
spin-\ensuremath{\tfrac{1}{2}} Born rule, with \ensuremath{\Pi} the cosine of the Bloch-sphere angle. The
polarization degree of freedom of an agent is a genuine SU(2) spin-\ensuremath{\tfrac{1}{2}}.
Second, three of them recouple to a half-integer with no scalar in the
decomposition: \ensuremath{\tfrac{1}{2} \otimes \tfrac{1}{2} \otimes \tfrac{1}{2} = \tfrac{1}{2} \oplus \tfrac{1}{2} \oplus \tfrac{3}{2}},
containing no spin-0 term, so the triadic confirmation closes to a net
half-integer. Equivalently, under a \ensuremath{2\pi} rotation each spin-\ensuremath{\tfrac{1}{2}} factor
contributes \ensuremath{-I}, and the product over the three agents is
\ensuremath{(-I)\otimes(-I)\otimes(-I) = -I} because the count is odd: the confirmed bit
acquires a sign under \ensuremath{2\pi} and returns at \ensuremath{4\pi}. Third, the character is
conserved up the recursion: a level-\ensuremath{k} bit is built from three
level-\ensuremath{(k-1)} bits, each half-integer, and three half-integers sum to a
half-integer, so spin-\ensuremath{\tfrac{1}{2}} propagates at every level. The oddness of the
three in \ensuremath{\bar{B}} = \ensuremath{3\bar{\alpha}} is the conservation law: an odd product of \ensuremath{-I}
factors is \ensuremath{-I} at every level, so the minimum-for-actualization integer
being odd is what makes matter fermionic and keeps it so up the tower.

\emph{Scope.} The spinor character is derived in full. Derived: the confirmed
bit carries half-integer total spin, forced by the oddness of the three
spin-\ensuremath{\tfrac{1}{2}} agents and conserved up the recursion; the fermion/boson split
(single node versus bound pair); the agent degree of freedom is a genuine SU(2)
spin-\ensuremath{\tfrac{1}{2}} through the resolution-rule identity; the recoupling reconstructs
the full three-dimensional rotation \emph{group} SO(3) with its \ensuremath{4\pi} SU(2)
double cover (below, "The rotation group from recoupling"); and the physical rotation of a confirmation event is
the diagonal (whole-event) action, forced by the flop's own symmetry
(below, "Why the whole-event rotation is diagonal"). The single point that is a reduction rather than an independent
derivation --- that the emergent metric geometry the reconstructed angles
populate is the same FCC geometry the substrate began with --- follows by the
form-invariance fixed point of Paper 3 \S\,4.7. What TMS does \emph{not} claim,
and correctly, is the dimensionful values of particle masses; those are the
separate calibration matter of \S\,8.4.3.

\emph{The rotation group from recoupling (the Penrose reconstruction).} That the
recoupling carries the \ensuremath{2\pi \to -I} sign on a single node is the entry point;
the full statement is that the emergent rotational symmetry \emph{is} SO(3). This
result is not original to the present framework: it is Penrose's Spin-Geometry
Theorem (Penrose 1971; Moussouris 1983), which established that the angles of
three-dimensional space can be reconstructed combinatorially from the recoupling
of spin systems, with no space presupposed --- the foundational insight of the
spin-network program. The relevance here is structural: the TMS confirmation node
is precisely a Penrose trivalent vertex --- a bit confirmed by the recoupling of
its agents --- so the theorem applies directly to the substrate, and what follows
is a verification of it in the substrate's own terms together with one corollary
(A.0 below) we have not found stated in the prior literature. The construction is
combinatorial, with no space presupposed: the angle between two
spin-\ensuremath{\tfrac{1}{2}} systems, defined purely by their recoupling overlap probability
through \ensuremath{\cos\theta = 2\lvert\langle i \vert j\rangle\rvert^{2} - 1} --- the
empirical-angle construction of Szabados (2022) --- yields a
matrix of direction cosines that is positive-semidefinite of rank exactly three
--- the recoupling-defined angles are the angles of vectors in a Euclidean
3-space --- and this is special to spin-\ensuremath{\tfrac{1}{2}} (a spin-1 system gives rank
greater than three; a non-inner-product matrix fails positivity). The full
eight-dimensional trivalent node \ensuremath{(\tfrac{1}{2} \oplus \tfrac{1}{2} \oplus \tfrac{3}{2})} carries a
faithful SU(2) representation: under the whole-event rotation it returns to
\ensuremath{-I} at \ensuremath{2\pi} and \ensuremath{+I} at \ensuremath{4\pi}, satisfies the group law
\ensuremath{R(g)R(h) = R(gh)}, and its orientation invariant \ensuremath{S_{1}\cdot(S_{2}\times S_{3})}
is an exact SU(2) scalar --- a rotation-invariant handedness. Thus SO(3) and its
double cover are reconstructed from the triadic recoupling, not assumed --- the
Penrose reconstruction, realized on the confirmation node.

\emph{Why the whole-event rotation is diagonal (from the flop).} A confirmation
is not three independent spins but one tetrahedral appointment --- three parent
slots feeding one centroid, resolved by the rule \ensuremath{\Pi = (\chi_{1}+\chi_{2}+\chi_{3})/3}
that is symmetric under permutation of the three parents. A physical rotation of
the event moves this rigid appointment as a unit, so it acts as the same
\ensuremath{g} on all three parents: the diagonal action \ensuremath{g\otimes g\otimes g}. This is
forced, not chosen. The resolution rule's \ensuremath{S_{3}} symmetry is compatible only with
the diagonal action (an independent per-parent rotation breaks it); and the real
probabilistic flop is covariant with no preferred axis --- the child bit's
up-probability satisfies \ensuremath{P(\text{event}, n) = P(g\cdot\text{event}, R_{g}n)}
identically, so ``rotate the whole event'' acts as one \ensuremath{g} on parents and
readout alike. The \ensuremath{1\to 4} transport, built from lattice rotations, propagates
this action up the recursion.

\textbf{Corollary A.0 (Three Is the Minimal Spatial Valence).} The valence three
of \ensuremath{\bar{B}} = \ensuremath{3\bar{\alpha}} is the \emph{minimal} node valence that simultaneously
carries a half-integer (fermionic) top representation and admits exactly one
orientation invariant --- a single handedness for the emergent 3-space. A
two-edge node has an integer top representation and no orientation (no
handedness); a four-edge node has an integer top representation and redundant
orientations. Only at three do fermionic statistics and a single spatial
handedness first coincide. Matter's being fermionic and space's having one
orientation are therefore the same fact about the number three.

\textbf{Corollary A.1 (Boson/Fermion Split).} A single triadic
confirmation node, recoupling an odd number (three) of spin-\ensuremath{\tfrac{1}{2}} agents,
carries half-integer spin and is a fermion. A bound pair of triadic nodes
recouples \ensuremath{\tfrac{1}{2} \otimes \tfrac{1}{2}} to integer spin and is a boson. This is the
mechanism behind the fermion/boson correspondence that Paper 3 and \S\,8.4
carry as preserved at every level: the parity of the node count is the
statistics.

\textbf{Corollary A.2 (Pauli Exclusion).} Fermionic statistics follow by
the spin--statistics connection, which is topological rather than
dynamical here: exchanging two identical triadic nodes is homotopic to
rotating one of them by \ensuremath{2\pi} relative to the other, which by Lemma A
multiplies the state by \ensuremath{-1}. The two-node state is therefore
antisymmetric, and two identical fermionic nodes placed in the same state
give \ensuremath{(\lvert a\rangle\lvert a\rangle - \lvert a\rangle\lvert a\rangle)/\sqrt{2} = 0} ---
the Pauli exclusion principle, derived from the same \ensuremath{2\pi} sign that makes
the node a spinor. Exclusion is used in the shell-filling count of \S\,8.4
and Appendix E.

\textbf{Lemma B (Continuum Symmetry under Coarse-Graining).} \emph{Under
the coarse-graining operator G of Paper 3 \S\,IV, the level-1 FCC lattice
substrate produces a continuum theory with Lorentz invariance exact at
leading order in the lattice spacing --- the FCC second-moment tensor is
isotropic (Paper 1 \S\,6.5), which fixes both the value and the
direction-independence of c. The leading correction is anisotropic and,
because the FCC point group \ensuremath{O_{h}} includes spatial inversion, enters only
at the quartic-in-momentum \ensuremath{O_{h}} invariant: the observable fractional
Lorentz-violation rate scales as \ensuremath{(E/E_{\mathrm{Planck}})^{2},} identically for
scalar and gauge-vector sectors. An odd-power \ensuremath{(E/E_{\mathrm{Planck}})^{1}}
contribution is forbidden by lattice inversion symmetry in the unbroken
substrate; it is unlocked only in the inversion-broken matter sector
(\S\,5.3), where its magnitude scales with the parent-record asymmetry.}

\emph{Proof sketch.} The FCC lattice has point group \ensuremath{O_{h},} which
includes spatial inversion. Inversion symmetry forces the dispersion
\ensuremath{\omega ^{2}(k)} to be even in k, so every odd-in-momentum term --- including
the cubic-in-momentum term --- vanishes identically. The lowest-order
Lorentz-invariant term is the standard kinetic energy, which respects
\ensuremath{O_{h}} exactly and Lorentz approximately; its second-moment tensor \ensuremath{\Sigma_{j}}
\ensuremath{\delta_{\mu}} \ensuremath{\delta_{\nu}} is isotropic, so the leading dispersion is
direction-independent --- the geometric fact that makes c invariant
across directions and frames. The lowest anisotropic \ensuremath{O_{h}} invariant is
therefore quartic in momentum --- the cubic harmonic \ensuremath{\Sigma_{i}} \ensuremath{k_{i}^{4}}
relative to the isotropic \ensuremath{(k^{2})^{2}} --- confirmed by direct
expansion of the FCC dispersion, whose fourth-moment tensor is
cubic-anisotropic. The fractional dispersion correction therefore scales
as \ensuremath{(k\cdot a)^{2},} i.e. \ensuremath{(E/E_{\mathrm{Planck}})^{2},} identically for scalar and
vector sectors; the shared exponent is forced by inversion symmetry, not
a distinct vector/scalar pair. A stronger \ensuremath{(E/E_{\mathrm{Planck}})^{1}} signal
requires breaking inversion, which occurs only in the matter sector
(\S\,5.3); it is the framework's most distinctive Lorentz-violation
prediction, and its magnitude tracks the parent-record asymmetry.

\emph{The dispersion, computed.} The scaling above is not only argued from
symmetry but computed in closed form. Summing the twelve nearest-neighbour phases
of the FCC discrete Laplacian and expanding to fourth order in the momentum gives
\ensuremath{\omega^{2} \propto (k_{x}^{2}+k_{y}^{2}+k_{z}^{2}) - \tfrac{1}{48}(k_{x}^{4}+k_{y}^{4}+k_{z}^{4}) - \tfrac{1}{16}(k_{x}^{2}k_{y}^{2}+k_{x}^{2}k_{z}^{2}+k_{y}^{2}k_{z}^{2}) + \cdots}
in units of the lattice spacing. The quadratic term is exactly \ensuremath{|k|^{2}} ---
isotropic, fixing \ensuremath{c}. The quartic term splits into an isotropic part
\ensuremath{-\tfrac{1}{48}|k|^{4}} and a genuinely anisotropic remainder
\ensuremath{-\tfrac{1}{48}(k_{x}^{2}k_{y}^{2}+k_{x}^{2}k_{z}^{2}+k_{y}^{2}k_{z}^{2})} that is
not proportional to \ensuremath{|k|^{4}}. Lorentz invariance is therefore genuinely broken
at the lattice scale, with a definite coefficient and a definite cubic angular
signature: the quartic coefficient along the crystallographic axes is
\ensuremath{-0.0208} in the [100] direction, \ensuremath{-0.0260} in [110], and \ensuremath{-0.0278} in [111], so
the phase velocity is direction-dependent with the ordering [100] \ensuremath{<} [110] \ensuremath{<}
[111].

\emph{The prediction, quantified and falsifiable.} The fractional phase-velocity
anisotropy between crystallographic directions is
\ensuremath{|\Delta v/v| \sim 0.007\,(E/E_{\mathrm{Planck}})^{2}}. This sits within the
established framework for testing Lorentz violation through modified dispersion
relations (Amelino-Camelia 1998; Jacobson, Liberati \& Mattingly 2002), and the
specific programme of bounding a discrete substrate's lattice scale from
Lorentz-violation data has been developed recently for quantum cellular automata
(Bisio et al. 2025, arXiv:2506.20136); the present result is the FCC realization
of that programme. Because the violation is \emph{quadratic} in
\ensuremath{E/E_{\mathrm{Planck}}}, it is far below all current bounds at
every accessible energy --- of order \ensuremath{10^{-58}} for optical photons, \ensuremath{10^{-33}} at
LHC energies, \ensuremath{10^{-39}} for the highest-energy gamma-ray-burst photons observed
by Fermi --- and it evades the astrophysical time-of-flight constraints that
already exclude a \emph{linear} Planck-scale term. That the linear (cubic-in-momentum)
term is observationally dead while the quadratic (CPT-even) term survives is a
general feature of this literature, and the substrate reproduces it for a
structural reason: lattice inversion symmetry (the \ensuremath{O_{h}} point group) forbids the
odd term. At the same time the prediction is not
unfalsifiable: at the highest observed ultra-high-energy cosmic-ray energies
(the GZK cutoff, \ensuremath{\sim 5\times10^{19}} eV) the predicted anisotropy reaches
\ensuremath{\sim 10^{-19}}, within one to two orders of magnitude of current sensitivity to
directional Lorentz violation (\ensuremath{\sim 10^{-18}}). The distinctive fingerprint is not
merely a magnitude but the specific cubic \ensuremath{O_{h}} angular pattern above, which a
generic Lorentz-violating theory would not predict. \textbf{Status: Prediction.}
The exact expansion, the angular coefficients, and the magnitude comparison are
reproducible \ensuremath{(dispersion_{25})} in the companion materials. The coefficient
\ensuremath{1/48} is specific to the minimal nearest-neighbour model; the quadratic scaling
and the cubic symmetry are forced by the lattice and are robust to the coupling
details.

Steps (ii)--(v) of the main theorem follow from the polarization
decomposition of \S\,4.2, the substrate coherence properties of \S\,3.2, and
the Born-rule structural correspondence of Paper 2 \S\,IV.

\subsection*{4.2 The Field Decomposition Theorem}

\begin{theorem}[Field Decomposition on FCC]
The continuum limit
of polarization propagation on the FCC lattice admits a natural
decomposition into five irreducible components under the cubic point
group \ensuremath{O_{h},} jointly spanning the 12-dimensional nearest-neighbor
displacement space: a scalar \ensuremath{(A_{1}g),} a two-dimensional deformation
tensor \ensuremath{(E_{g}),} a three-dimensional tangential tensor \ensuremath{(T_{2}g),} a polar
vector \ensuremath{(T_{1}u),} and a three-dimensional odd-parity tensor \ensuremath{(T_{2}u).}
\end{theorem}

\emph{Proof sketch.} The FCC lattice has 12 nearest neighbors, each at
face-diagonal distance \ensuremath{\ell_{p}} in the standard FCC parametrization (Paper 1
\S\,VI). The 12 nearest-neighbor displacement vectors \{\ensuremath{r_{j}}\} are symmetric
under the cubic point group \ensuremath{O_{h},} and their irreducible decomposition
under \ensuremath{O_{h}} is

\emph{12 = \ensuremath{A_{1}g} \ensuremath{\oplus} \ensuremath{E_{g}} \ensuremath{\oplus} \ensuremath{T_{2}g} \ensuremath{\oplus} \ensuremath{T_{1}u} \ensuremath{\oplus} \ensuremath{T_{2}u} (with dimensions 1 + 2 + 3 + 3 +
3 = 12),}

with character vector \ensuremath{[E, 8C_3, 6C_2', 6C_4, 3C_2, i, 6S_4, 8S_6, 3\sigma_h, 6\sigma_d] = [12, 0, 2, 0, 0, 0, 0, 0, 4, 2]} on the twelve nearest-neighbor displacements, verified two independent ways (character contraction and explicit class representatives). The parity split is six even (\ensuremath{A_{1}g,} \ensuremath{E_{g},} \ensuremath{T_{2}g}) and six odd (\ensuremath{T_{1}u,} \ensuremath{T_{2}u}),

producing five irreducible components:

\begin{itemize}
\item
  \emph{Scalar \ensuremath{(A_{1}g,} 1-dimensional, parity-even):} the local mean
  polarization \ensuremath{\langle \psi \rangle} across the 12 neighbors.
\end{itemize}

\begin{itemize}
\item
  \emph{Polar vector \ensuremath{(T_{1}u,} 3-dimensional, parity-odd):} the polarization
  gradient \ensuremath{\nabla \psi} across the 12 neighbors --- three independent directional
  components transforming as a vector under \ensuremath{O_{h}.}
\end{itemize}

\begin{itemize}
\item
  \emph{Deformation tensor \ensuremath{(E_{g},} 2-dimensional, parity-even):} the
  cubic-symmetric traceless symmetric-tensor component of the
  polarization quadrupole, with diagonal basis \ensuremath{(\partial _{\mathrm{xx}}} \ensuremath{-} \ensuremath{\partial yy),} \ensuremath{(\partial _{\mathrm{xx}}} + \ensuremath{\partial yy} \ensuremath{-}
  \ensuremath{2\partial zz)/\sqrt{3}.}
\end{itemize}

\begin{itemize}
\item
  \emph{Tangential tensor \ensuremath{(T_{2}g,} 3-dimensional, parity-even):} the
  off-diagonal symmetric-tensor component, \ensuremath{(\partial xy,} \ensuremath{\partial yz,} \ensuremath{\partial zx).}
\end{itemize}

\begin{itemize}
\item
  \emph{Odd tensor \ensuremath{(T_{2}u,} 3-dimensional, parity-odd):} the second
  parity-odd sector --- the odd-parity partner completing the twelve
  dimensions alongside the polar vector \ensuremath{T_{1}u.} Physically: the
  inversion-odd off-diagonal response of the directional polarization
  across the 12 nearest neighbors, the odd counterpart of the even
  tangential tensor \ensuremath{T_{2}g.} \ensuremath{\square}
\end{itemize}

The five components jointly exhaust the 12-dimensional nearest-neighbor
space. Each propagates by the same wave equation but couples to the
lattice geometry differently, with angular structure indexed by its \ensuremath{O_{h}}
representation.

The physical interpretation. The scalar \ensuremath{(A_{1}g)} is the matter-density
sector (the analog of Higgs-like scalar content). The polar vector \ensuremath{(T_{1}u)}
is the polar gauge-field sector (the analog of the photon and the polar
parts of the gauge bosons). The deformation and tangential tensors \ensuremath{(E_{g}}
\ensuremath{\oplus} \ensuremath{T_{2}g)} are the spacetime-geometry sector (the analog of the metric, with
\ensuremath{E_{g}} + \ensuremath{T_{2}g} jointly spanning the five independent components of a 3D
traceless symmetric tensor). The second odd sector \ensuremath{(T_{2}u)} is the
inversion-odd tensor partner of the even tangential tensor \ensuremath{T_{2}g:} it
carries the odd-parity off-diagonal response that completes the twelve
nearest-neighbor dimensions. The decomposition reproduces the structural
skeleton of the field content of the Standard Model coupled to gravity as
a scalar, a polar gauge vector, a two-piece symmetric-tensor geometry
sector, and an additional odd-parity tensor.

The chiral structure of the weak interaction is supplied separately by the
inversion-odd parent-record break \ensuremath{d_{z}} through the Jackiw--Rebbi
domain-wall mechanism of \S\,8.4.4, which is where chirality is treated.

This is a structural correspondence to the field content of the Standard
Model coupled to gravity. The TMS does not claim to derive that field
content in full; that requires the explicit representation-theoretic
mapping reserved as a structural correspondence in \S\,VIII. The TMS claim
is that the geometric structure of the FCC lattice naturally produces
five irreducible field-type sectors covering all 12 dimensions of the
nearest-neighbor space, with the structural roles of scalar matter,
polar gauge, symmetric-tensor geometry, and an odd-parity tensor
partner.

\textbf{Bridge to \S\,VIII.} The gauge-symmetry inheritance and
fermion-boson decomposition originally developed in this paper's \S\,4.3
and \S\,4.4 are consolidated into \S\,VIII below, where the
bifurcation-derived cross-level commutator mechanism of \S\,8.2 provides
the non-Abelian closure that completes the structural correspondence to
the Standard Model gauge group.

\section*{V. Initial Conditions and the Big Bang Correspondence}
\addcontentsline{toc}{section}{V. Initial Conditions and the Big Bang Correspondence}

\subsection*{5.1 The Initial State of \ensuremath{\Xi (R)}}

By the construction of \S\,2.3, the substrate \ensuremath{\Xi (R)} begins its operation
with initial polarization landscape \ensuremath{\psi '} inherited from \ensuremath{\mathcal{B}(R)} via the
coarse-graining operator G. The initial state has five structural
properties forced by the construction:

\textbf{Maximum informational density.} \ensuremath{\mathcal{B}(R)} is the saturated
biographical record of an exhausted region: \ensuremath{\Sigma_{\mathrm{sampled}}(R,} \ensuremath{\tau *)} \ensuremath{\approx} \ensuremath{\Sigma (R).}
Every configuration in \ensuremath{\Sigma (R)} has been sampled, meaning \ensuremath{\mathcal{B}(R)} contains the
maximum information density that R can hold. Translated to \ensuremath{\Xi (R)} under \ensuremath{\hat{R},}
this becomes the maximum initial agent density and maximum initial
polarization complexity. The new substrate begins at the upper bound of
what its agent count can encode.

\textbf{Near-uniform background.} The exhaustion condition is regional,
not point-wise. Different points in R have been sampled with comparable
thoroughness --- the random walk through \ensuremath{\Sigma (R)} has visited the
configuration space evenly. The polarization landscape across \ensuremath{\mathcal{B}(R),}
averaged over the imprint depth, is therefore close to uniform at large
scales. Translated to \ensuremath{\Xi (R),} this becomes a near-uniform initial \ensuremath{U_{\mathrm{sh}}'}
across the new substrate.

\textbf{Small fluctuations from late-phase configurations.} \ensuremath{\Sigma_{\mathrm{sampled}}}
approaches but does not exactly equal \ensuremath{\Sigma (R)} at \ensuremath{\tau} = \ensuremath{\tau *;} some late-phase
configurations are sampled less thoroughly than the bulk. These
late-phase configurations contribute residual structure that deviates
from the bulk uniformity. The size of the fluctuations is set by the
ratio of late-phase configurations to bulk configurations, which is
bounded by the per-flop noise injection rate divided by the imprinting
rate (Paper 3 \S\,3.3). Translated to \ensuremath{\Xi (R),} this becomes small residual
fluctuations in the initial polarization landscape, on top of the
near-uniform background.

\textbf{Large-scale flatness.} By the Principle of Generative Economy,
the initial polarization landscape of \ensuremath{\Xi (R)} is the minimum-curvature
configuration consistent with the exhausted state. Curvature in the
polarization landscape would correspond to an inhomogeneity favoring
particular configurations over others, which would have driven
\ensuremath{\Sigma_{\mathrm{sampled}}} away from \ensuremath{\Sigma (R)} at the parent level rather than toward
saturation. Saturation is incompatible with large-scale curvature.
Translated to \ensuremath{\Xi (R),} the initial geometry (the tensor sector of \S\,4.2) is
large-scale flat.

\textbf{Causally connected past.} \ensuremath{\mathcal{B}(R)} is the biography of a single
bounded region. Every confirmed bit in \ensuremath{\mathcal{B}(R)} is part of a single causal
record --- they all trace back through the same parent-level causal
stack. By the construction of \ensuremath{\hat{R},} every new agent in \ensuremath{\Xi (R)} inherits this
property: they all derive from a single parent-level region whose own
causal history is internally connected. Translated to \ensuremath{\Xi (R),} the initial
state is causally connected: every part of the new substrate has a past
that was once in causal contact at the parent level, even if the
corresponding regions are now spacelike-separated in \ensuremath{\Xi (R).}

\subsection*{5.2 The Big Bang Correspondence Theorem}

\begin{theorem}[Big Bang Correspondence]
The initial state of \ensuremath{\Xi (R)}
under \ensuremath{\hat{R},} viewed by an observer at hierarchical level k \ensuremath{\geq} 1 within \ensuremath{\Xi (R),}
exhibits the structural signatures of the observed cosmological initial
state: maximum informational density (``hot dense early phase''),
near-uniform background (CMB-like uniformity), small fluctuations on the
uniformity \ensuremath{(10^{-5}-\text{scale}} CMB anisotropies), large-scale flatness, and a
horizon-connected past. The correspondence is structural; numerical
pinning of specific cosmological parameters is reserved for simulation.
\end{theorem}

\begin{proof}
Each of the five structural properties of the initial
state derived in \S\,5.1 maps to a structural feature of the observed Big
Bang scenario:

\begin{itemize}
\item
  Maximum informational density (i) \ensuremath{\to} hot dense early phase. The energy
  density of the early universe, in the cosmological standard model, is
  bounded above by the Planck scale; in the TMS reading, the initial
  \ensuremath{\Xi (R)} has maximum agent density and maximum polarization complexity,
  which corresponds to maximum energy density in the field-theoretic
  continuum limit by the standard QCA-to-QFT identification of
  polarization activity with energy.
\end{itemize}

\begin{itemize}
\item
  Near-uniform background (ii) \ensuremath{\to} CMB uniformity. The near-uniformity of
  the initial polarization landscape, averaged over imprint depth,
  corresponds in the continuum limit to a near-uniform initial field
  configuration with corresponding near-uniform temperature.
\end{itemize}

\begin{itemize}
\item
  Small fluctuations (iii) \ensuremath{\to} CMB anisotropy. The residual structure from
  late-phase configurations of \ensuremath{\mathcal{B}(R)} corresponds to small fluctuations on
  top of the uniform background, providing the seed inhomogeneities for
  cosmological structure formation. The size of these fluctuations is
  bounded by the per-flop noise-to-imprint ratio of Paper 3 \S\,3.3, which
  gives a structurally small but non-zero value --- consistent with the
  observed \ensuremath{10^{-5}-\text{scale}} CMB anisotropies though not pinned numerically at
  this level.
\end{itemize}

\begin{itemize}
\item
  Large-scale flatness (iv) \ensuremath{\to} spatial flatness problem solved. The
  minimum-curvature property forced by saturation removes the need for
  inflationary fine-tuning of initial curvature. The TMS initial state
  is naturally flat at large scales because it is the PGE-selected
  minimum-specification continuation of a saturated parent record.
\end{itemize}

\begin{itemize}
\item
  Causally connected past (v) \ensuremath{\to} horizon problem solved. Every part of
  \ensuremath{\Xi (R)} traces back to a single parent-level region R that was internally
  causally connected. The horizon problem of standard cosmology --- how
  regions of the present-day CMB that are now spacelike-separated could
  have come into thermal equilibrium --- is resolved structurally: they
  were in causal contact at the parent level, before reinterpretation,
  and their thermal equilibrium is inherited from that contact.
\end{itemize}

Each correspondence is structural rather than numerical. The numerical
values of CMB temperature, scalar spectral index, baryon-to-photon
ratio, and analogous quantitative observables depend on the specific
\ensuremath{\mathcal{B}(R)} being reinterpreted, the hierarchical level at which we are
observing, and the specific combinatorics of the FCC lattice at
successive levels. These numerical values are reserved for simulation
work. The structural claim --- that the TMS reinterpretation event
reproduces the observed structural features of the Big Bang --- is what
is established here.

\end{proof}

\subsection*{5.3 Matter--Antimatter Structural Correspondence}

The sublattice bifurcation of \S\,6.1 produces, at every hierarchical level
k \ensuremath{\geq} 1, two FCC sublattices A and B related by spatial inversion. This is
the substrate-level analog of the matter / antimatter sector
decomposition of QFT, with the sublattice exchange A \ensuremath{\leftrightarrow} B playing the
role of one CP-like discrete symmetry. We develop this correspondence in
three claims.

\begin{claim}[Structural --- Derived]
Every hierarchical
level k \ensuremath{\geq} 1 has exactly two interpenetrating FCC sublattices, related by
spatial inversion through the level-(k+1) centroids. The sublattices
share no vertices, no edges, and no faces at level k; they share only
centroids at level (k+1). Their [[4,2,2]] stabilizer dynamics
run independently at level k.
\end{claim}

This is the geometric content of \S\,6.1 promoted to a structural law. The
sublattices A and B at level k are forced by the bifurcation under PGE
--- each runs its own level-k FCC tetrahedral closure and its own
[[4,2,2]] code --- and their only geometric overlap occurs at
the next level's actualization sites. Stage 5 verifies this through
\ensuremath{L_{\mathrm{box}}} = 11 (n = 5) and Stage 6 confirms identical structure at level 2
(Appendix E).

\begin{claim}[Structural correspondence]
The two-sublattice
structure at every level k \ensuremath{\geq} 1 is the substrate-level analog of the
matter / antimatter sector decomposition of QFT, with the sublattice
exchange A \ensuremath{\leftrightarrow} B playing the role of one CP-like discrete symmetry.
Spatial inversion through any level-(k+1) centroid is exact: every
A-tetrahedron is mapped to a B-tetrahedron with the same level-(k+1)
centroid under coordinate inversion through that centroid.
\end{claim}

The empirical content of the inversion-partnership is verified in Stage
5 (Appendix E) at 100\% of shared centroids across \ensuremath{L_{\mathrm{box}}} = 5, 7, 9, 11
--- every A-tet centroid coincides with a B-tet centroid, and the B-tet
is the exact coordinate inversion of the A-tet through their common
centroid. In Stage 6, the same 100\% inversion-partner pairing holds at
level 2 across \ensuremath{L_{\mathrm{box}}} = 6, 8, 10, 12.

The structural correspondence to QFT's matter / antimatter sectors: in
standard QFT, matter and antimatter fields are related by CP --- a
composite of charge conjugation C and spatial parity P. In the TMS
picture, the substrate-level analog is the sublattice exchange A \ensuremath{\leftrightarrow} B,
which is exact spatial inversion (the P component) combined with a
sublattice-label flip (the C-analog component, with ``charge''
identified as the sum-mod-4 invariant \ensuremath{\in} \{1, 3\} at the level-k site).
The TMS sublattice symmetry is not literally CP --- there is no charge
conjugation in the usual QFT sense --- but the structural role it plays
in distinguishing the two sublattices and their independent dynamics is
the same role CP plays in QFT. This is consistent with Harlow, Usatyuk,
and Zhao's (2025) holographic-gravity result that internal observer
structure restores the complexity of a closed universe: the substrate's
two-sublattice structure provides exactly the observer/anti-observer
doubling that their argument identifies as required for non-trivial
closed-cosmos quantum gravity.

\begin{claim}[Prediction / Conjecture]
A reinterpretation
event acting on a parent biographical record \ensuremath{\mathcal{B}(R)} without exact
sum-mod-4 symmetry produces level-1 substrates \ensuremath{\Xi (R)} with unequal A and B
populations. The matter--antimatter asymmetry of the observed universe
is a generic consequence of reinterpretation from a
non-perfectly-symmetric parent record.
\end{claim}

The reinterpretation operator \ensuremath{\hat{R}} maps the saturated parent record \ensuremath{\mathcal{B}(R)} to
a new level-1 substrate by the polarization-landscape transition of
\S\,2.3. The sum-mod-4 content of \ensuremath{\mathcal{B}(R)} determines, at the level-1 boundary,
the relative weights of A versus B sublattice population in \ensuremath{\Xi (R).} If
\ensuremath{\mathcal{B}(R)} carries equal sum-mod-4 weight, then \ensuremath{\Xi (R)} has equal A and B
populations and is matter-antimatter symmetric. If \ensuremath{\mathcal{B}(R)} has a residual
sum-mod-4 imbalance --- as a generic consequence of the parent's
late-phase polarization landscape --- then \ensuremath{\Xi (R)} inherits this imbalance
as an unequal A versus B population, manifesting in the field-theoretic
continuum limit as a matter-antimatter asymmetry.

This is the substrate-level structural account of baryogenesis. The TMS
framework provides three conditions that jointly produce
matter-antimatter asymmetry, structurally parallel to Sakharov's (1967)
three conditions but adapted to TMS's symmetry structure:

\emph{(i) Sum-mod-4-changing transitions exist at the reinterpretation
boundary (analog of baryon-number violation).} \ensuremath{\hat{R}} can redistribute the
sum-mod-4 spectrum from \ensuremath{\mathcal{B}(R)} into \ensuremath{\Xi (R)'s} new level-1 substrate. Within
any single substrate level, [[4,2,2]] dynamics preserves
sum-mod-4 (A stays A, B stays B); only \ensuremath{\hat{R}} can change it.

\emph{(ii) The parent biographical record \ensuremath{\mathcal{B}(R)} carries
inversion-asymmetric structure (analog of CP violation).} The late-phase
polarization landscape of an exhausted region generically has asymmetric
content under the \ensuremath{A\leftrightarrow B} sublattice exchange --- a generic structural
property of a saturated record produced by non-perfectly-symmetric
parent dynamics.

\emph{(iii) \ensuremath{\hat{R}} is a phase transition, not an equilibrium process (analog
of departure from thermal equilibrium).} By \S\,2.7, \ensuremath{\hat{R}} has atomic,
locality-preserving, generative-resetting structure; the asymmetry it
produces in \ensuremath{\Xi (R)} is not subject to detailed-balance relaxation because \ensuremath{\hat{R}}
is not reversible in the relevant sense.

The parallel is structural rather than literal. TMS does not have C and
P as separately violable symmetries --- the \ensuremath{A\leftrightarrow B} sublattice exchange is a
single discrete operation combining what QFT distinguishes as two ---
and TMS does not have thermal equilibrium in the QFT sense. But each TMS
condition closes the same kind of structural leak its Sakharov
counterpart closes, and all three are jointly necessary: take any one
away and the asymmetry either does not generate or does not persist into
\ensuremath{\Xi (R).}

\textbf{Remark on the sub-critical reading.} Claim 3 frames the
asymmetry as ``unequal A versus B populations.'' Analysis of the
inversion-broken band structure on the A/B sublattice doublet sharpens
this: a \emph{uniform} (constant) population imbalance is structurally a
featureless trivial insulator with no fermion content. Matter-bearing
band structure requires the inversion-breaking term to be
\emph{structured} in k --- to have nonzero spatial variation, not merely
a nonzero mean. The ``imbalance'' of Claim 3 must therefore be read as
the \emph{sub-critical mean of a structured field}. Direct simulation of
the triadic confirmation rule \ensuremath{(\bar{B}} = \ensuremath{3\bar{\alpha}} resolution under causal sweep,
Paper 2 \S\,\S\,3.2, 4.1) shows that the substrate's own dynamics produce
exactly this regime: the late-time record has a near-zero mean \ensuremath{(\sim 6} \ensuremath{\times}
\ensuremath{10^{-2})} embedded in spatial structure of amplitude order unity, giving a
structure-to-mean ratio \ensuremath{\sim 17,} well inside the sub-critical regime. The
matter-bearing condition (sub-critical mean of a structured field) is
therefore satisfied by the substrate's own confirmation dynamics, not as
an external imposition. The condition for matter to exist at all is that
this sub-critical regime is maintained --- reframing the observed
smallness of the baryon asymmetry \ensuremath{(\sim 10^{-9})} as a \emph{necessary condition
for matter to exist}, not a quantitative puzzle to explain away. The
full geometric closure (FCC 3D, A/B doublet, \ensuremath{1\to 4} multiplier) and the
resulting mass spectrum remain the subject of subsequent
biographical-simulation work (cf. \S\,8.4.3).

\textbf{Status of Results for \S\,5.3:}

\begin{itemize}
\item
  \emph{Derived:} The two-sublattice structure at every level k \ensuremath{\geq} 1
  (Claim 1).
\end{itemize}

\begin{itemize}
\item
  \emph{Structural correspondence:} The TMS sublattice exchange as
  substrate-level analog of QFT's CP symmetry (Claim 2).
\end{itemize}

\begin{itemize}
\item
  \emph{Prediction / Conjecture:} Matter-antimatter asymmetry from
  non-symmetric reinterpretation (Claim 3).
\end{itemize}

\subsection*{5.4 The Horizon Problem Resolution in Detail}

The horizon problem of standard cosmology is among the most striking
observational puzzles: regions of the CMB separated by angles larger
than the particle horizon at recombination exhibit thermal equilibrium
that they could not have established through standard light-cone
communication. Standard cosmology addresses this through inflation --- a
brief period of exponential expansion that places initially-nearby
regions outside one another's causal horizons. The TMS provides a
different and structurally more economical resolution.

Under \ensuremath{\hat{R},} the entirety of \ensuremath{\Xi (R)} traces back to a single bounded region R
of the parent substrate. Before reinterpretation, R is a causally
connected region under the parent substrate's connectivity --- its
internal points are within communication of one another through the
parent-level FCC propagation. The biographical record \ensuremath{\mathcal{B}(R)} is the
integrated history of this internally connected region. When \ensuremath{\hat{R}} acts,
\ensuremath{\mathcal{B}(R)} becomes the initial state of \ensuremath{\Xi (R),} but the causal connectivity of
the parent-level R is not preserved as causal connectivity in \ensuremath{\Xi (R)} ---
it becomes prior-correlation between regions of \ensuremath{\Xi (R)} that may now be
spacelike-separated.

This is the structural origin of the horizon-problem resolution: regions
of the present-day CMB that appear spacelike-separated in \ensuremath{\Xi (R)'s}
spacetime were not spacelike-separated in the parent-level R that
produced them. Their thermal equilibrium is not a coincidence to be
explained by inflation; it is a structural inheritance from the
parent-level causal connectivity.

This account makes a structural prediction that distinguishes it from
inflationary cosmology: the CMB correlations should reflect the specific
causal structure of the parent-level R, not the post-inflationary
smoothing of an arbitrary initial state. In principle, statistical
signatures of the parent-level causal topology should be detectable in
high-precision CMB measurements. This is a falsifiable prediction whose
explicit empirical content requires simulation work to compute the
expected signatures.

\subsection*{5.5 The Flatness Problem Resolution}

Standard cosmology requires a remarkable fine-tuning of the early
universe's curvature: any deviation from spatial flatness at the Planck
epoch grows rapidly under Friedmann evolution, so the observed
near-flatness today requires \ensuremath{\Omega} near 1 to one part in \ensuremath{10^{60}} at early
times. Inflation provides one resolution by driving curvature toward
zero exponentially. The TMS provides a different structural resolution.

Under \ensuremath{\hat{R},} the initial geometry of \ensuremath{\Xi (R)} is the minimum-curvature
configuration consistent with the exhausted parent state. This is forced
by two arguments. First, by the Principle of Generative Economy, the
minimum-specification continuation is selected; curvature in the initial
polarization landscape would require additional specification beyond
what the exhausted parent state determines. Second, by the exhaustion
condition itself, the parent-level \ensuremath{\Sigma_{\mathrm{sampled}}} approached \ensuremath{\Sigma (R)} evenly; an
unevenly sampled exhausted state would correspond to a curvature in the
polarization landscape that would have biased further sampling and
prevented exhaustion. Saturation is geometrically incompatible with
large-scale curvature.

The initial flatness of \ensuremath{\Xi (R)} is therefore structural rather than
fine-tuned. No inflationary mechanism is required to drive curvature
toward zero; the initial state is flat because the structural conditions
of saturation force it to be. Small residual curvature corresponds to
the small fluctuations of \S\,5.1(iii), bounded by the same per-flop
noise-to-imprint ratio.

\subsection*{5.6 The Initial Entropy Problem and Causal Anchoring}

Penrose's Weyl-curvature hypothesis and the broader ``initial entropy
problem'' of cosmology pose a structural question: why was the early
universe in such a remarkably low-entropy state, given that random
initial conditions should have produced very high entropy? The standard
cosmological model has no internal mechanism to explain this; it must be
posited as an initial condition.

The TMS supplies a structural mechanism. The initial state of \ensuremath{\Xi (R)} is
the saturated biography \ensuremath{\mathcal{B}(R),} which has high informational density but
also high structure: every program, every subsystem, every multi-level
hierarchy that existed in R is encoded in \ensuremath{\mathcal{B}(R).} The polarization
landscape inherited as \ensuremath{U_{\mathrm{sh}}'} of \ensuremath{\Xi (R)} is therefore not random --- it
carries the algorithmic structure of all the programs that ran in R
before reinterpretation. From the perspective of the new substrate, this
is a low-entropy initial state: the polarization landscape is highly
compressible because it encodes the program structure of the parent's
biography.

The reinterpretation boundary transmits a two-part inheritance. The
accessible record --- the program structure of the parent's biography,
scaling as the boundary area --- is highly compressible and constitutes
the daughter's low-entropy initial condition, fixing the matching
condition and the arrow. The shed residue --- the resolvable distinction
discarded by the lossy fold at every confirmation, conserved under
Axiom 3 but un-re-readable at the parent level --- is, by the
maximum-entropy character of the discarded micro-configurations,
indistinguishable from fresh latent flux and constitutes the daughter's
\ensuremath{U_{\mathrm{sh}}}. The arrow \ensuremath{\hat{R}} produces thus has its origin in this shed
distinction: the irreversibility is not posited but is the accumulated
per-fold loss. No information is created (Axiom 3) and none is lost; what
changes is which part is accessible at which level.

This connects the resolution of the initial entropy problem directly to
Axioms 3 and 4 of Paper 1. The Weyl-curvature hypothesis becomes, in TMS
terms, the statement that the initial state inherits the structural
compressibility of the parent-level programs. The compressibility is
bounded above by the parent-level \ensuremath{K_{\psi}} at the moment of
reinterpretation, which is itself bounded by the parent-level program
count and complexity. The initial low entropy is therefore a derived
property of the reinterpretation protocol, not a posited condition.

\subsection*{5.7 The Initial Singularity Resolution}

Standard cosmology faces a singularity at t = 0: classical general
relativity breaks down because curvature and density diverge. Quantum
gravity programs typically replace the singularity with a non-singular
initial state --- a bounce, a quantum-pressure-supported minimum (loop
quantum cosmology), or a primordial Bose-Einstein-condensate-like state
(analog cosmology). The TMS provides a structural resolution: there is
no singularity, only a reinterpretation event.

At the reinterpretation moment \ensuremath{\tau} = \ensuremath{\tau *(R),} the substrate is operating
normally --- flop-by-flop confirmation at the parent scale. The
transition to \ensuremath{\Xi (R)} is the change of read protocol, not a change in
physical state. There is no infinite density, no infinite curvature, no
divergent quantity. The parent substrate's \ensuremath{\mathcal{B}(R)} is a finite
informationally-rich record; its reinterpretation as \ensuremath{\Xi (R)} is a finite
mapping that produces a finite new substrate.

What an observer in \ensuremath{\Xi (R)} extrapolates backward through the daughter
substrate's history reaches, at \ensuremath{\tau '} \ensuremath{\to} \ensuremath{0^{+},} the boundary at which the
daughter substrate begins to operate. Before \ensuremath{\tau '} = 0, the daughter
substrate did not exist; the parent substrate was running. The ``initial
singularity'' of standard cosmology corresponds in the TMS picture to
the reinterpretation boundary, not to a physical singularity. The
boundary is a phase transition (\S\,2.7), and physical quantities are
finite on both sides of it.

This is the structural origin of the standard ``resolution'' of the Big
Bang singularity in quantum-gravity programs. The TMS provides a
particularly clean version: the boundary is a structural identification
rather than a dynamical event, and the apparent singularity is an
artifact of extrapolating the daughter substrate's dynamics into the
region where the protocol changes.

\subsection*{5.8 Cosmological Natural Selection and Repeated Reinterpretation}

Smolin (1992) proposed cosmological natural selection: universes that
are good at producing black holes (and thus daughter universes through
black-hole bounces) come to dominate the multiverse. Dowker and Zalel
(2017) showed that causal-set classical sequential growth models exhibit
a similar renormalization across creation events. The TMS provides a
structural setting for this same idea: each reinterpretation event
produces a daughter substrate whose initial conditions inherit the
structural content of the parent's biography.

If the daughter substrate \ensuremath{\Xi (R)} is itself capable of forming subsystems,
hierarchies, and bounded regions that eventually exhaust their own
\ensuremath{\Sigma (R'),} then \ensuremath{\Xi (R)} will itself undergo reinterpretations within its
bounded regions. These produce grand-daughter substrates \ensuremath{\Xi (\Xi (R)),} and so
on. The chain of reinterpretation events forms a generational hierarchy
of substrates, each inheriting its initial conditions from the terminal
biography of its parent.

Two structural predictions follow. First, daughter substrates inherit
the algorithmic structure of their parents through the
polarization-landscape boundary condition; substrates whose parents had
richer program structure produce daughters with richer initial
conditions, which in turn produce more programs, which exhaust faster,
which produce more daughters. This is a structural form of cosmological
natural selection. Second, the \ensuremath{c^{(k)}} sequence at successive
reinterpretation events forms a chain whose total dispersion is bounded
by the product of per-level \ensuremath{(1/\sqrt{2})} factors --- providing a structural
bound on how many generational reinterpretation steps separate any given
substrate from the original parent.

A third structural prediction follows from the arrow-tower analysis of
\S\,8.4.2 applied to the generational chain described above. Each
reinterpretation event freezes one frozen-arrow direction into its
daughter's spatial geometry; the PGE-compressed tower saturates at the
substrate's spatial dimension. If each fermion generation corresponds to
dominant alignment with one frozen arrow, then the observation of three
generations implies the cosmological recursion has \emph{depth at least
three} ancestor substrates separating our current substrate from the
original parent. This converts cosmological natural selection from a
structural framework into a count with empirical traction: the observed
generation count places a strict lower bound on the recursion depth.

All three predictions are structural correspondences whose specific
empirical content requires simulation work to make quantitative. The TMS
framework establishes that the generational structure exists and that
its depth admits the lower bound supplied by \S\,8.4.2; the specific
evolutionary trajectory it produces is a target for subsequent
investigation.

\section*{VI. Numerical Closure: \ensuremath{L_{p},} \ensuremath{c^{(k)},} and the GUP Parameter}
\addcontentsline{toc}{section}{VI. Numerical Closure: \ensuremath{L_{p},} \ensuremath{c^{(k)},} and the GUP Parameter}

\subsection*{6.1 The Minimum Program Extent \ensuremath{L^{p(1)}} (Bifurcated FCC Construction)}

Paper 3 \S\,4.1 characterized \ensuremath{L_{p}} as the minimum spatial extent of a
stabilizer-satisfying T-closed configuration with K(C) \textless{}
\textbar C\textbar, estimated at 4--8 lattice spacings. Paper 4 \S\,7.5
reserved the exact value for combinatorial enumeration. We close this
debt here by explicit bifurcated FCC construction, with the result
empirically verified by the Stage 5 and Stage 6 simulation work
documented in Appendix E.

The construction proceeds in three steps. First, we identify the
geometric structure of the level-1 substrate (the sublattice
bifurcation). Second, we specify the minimum configuration that
satisfies all level-1 program conditions simultaneously (the
stella-octangula-tile-plus-boundary, derived below). Third, we measure
the spatial extent of this minimum configuration in substrate units and
confirm \ensuremath{L^{p(1)}} = \ensuremath{3\sqrt{2}} \ensuremath{\ell_{p}.}

\subsubsection*{6.1.1 The Sublattice Bifurcation of the Level-1 Substrate}

By the Recursion-Preservation Theorem of \S\,2.5 and the coarse-graining
operator G of Paper 3 \S\,IV, the level-1 substrate produced by
reinterpretation reproduces the substrate-level FCC lattice with the
same triadic confirmation rule and the same [[4,2,2]] stabilizer
structure at the new scale. We adopt the \emph{storage convention}
throughout this section: at each hierarchical level, lattice positions
are written in integer coordinates with the level-k FCC sublattice
scaled by a factor of 2 relative to \ensuremath{\text{level}-(k-1).} In this convention, the
level-0 FCC consists of integer-coordinate sites with sum i + j + k
even; the level-1 sites lie on the all-odd-coord integer positions (the
centroids of the level-0 tetrahedra after the \ensuremath{\times 2} scaling).

\textbf{The bifurcation.} The all-odd-coord level-1 sites do not form a
simple FCC. Under the squared-distance metric, the nearest-neighbor
distance among all-odd-coord integers is 2 (e.g., from (1,1,1) to
(1,1,3)) --- a \emph{simple cubic} coordination, not FCC. The
simple-cubic NN structure supports no four-mutually-nearest tetrahedra,
which would block the level-1 [[4,2,2]] stabilizer dynamics
altogether.

The structural resolution is that the all-odd-coord level-1 sites split
into two interpenetrating FCC sublattices by the sum-mod-4 invariant:

\begin{itemize}
\item
  \textbf{Sublattice A} (sum \ensuremath{\equiv} 1 mod 4): sites such as (1,1,3), (1,3,1),
  (3,1,1), (3,3,3) --- all-odd-coord positions whose coordinate-sum is
  congruent to 1 modulo 4.
\end{itemize}

\begin{itemize}
\item
  \textbf{Sublattice B} (sum \ensuremath{\equiv} 3 mod 4): the inversion partners (1,1,1),
  (1,3,3), (3,1,3), (3,3,1) --- all-odd-coord positions whose
  coordinate-sum is congruent to 3 modulo 4.
\end{itemize}

Within each sublattice, the nearest-neighbor squared distance is 8 ---
the face-diagonal of the underlying cubic cell at edge 2. Each
sublattice is therefore an FCC at NN distance \ensuremath{2\sqrt{2}} in storage
coordinates, with the standard FCC coordination number of 12 (Stage 5
verifies this empirically through n = 5; see Appendix E). The two
sublattices are related by spatial inversion through any all-even-coord
site (the level-2 sites), and they share no level-1 vertices, edges, or
faces; they share only centroids at level 2.

This structure is analogous to the \textbf{diamond cubic} crystal
familiar from solid-state physics, with the two FCC sublattices playing
the roles of the two atom positions in the diamond unit cell. In TMS
terms, the sublattice bifurcation is a structural consequence of the FCC
+ [[4,2,2]] recursion under the storage scaling --- not an
additional axiom.

\subsubsection*{6.1.2 The Minimum Level-1 Program Configuration}

By the bifurcated structure, a complete level-1 program configuration
requires elements from \emph{both} A and B sublattices, because each
sublattice independently hosts an FCC tetrahedron and a
[[4,2,2]] stabilizer code, and the level-2 centroid where their
dynamics jointly read out is the actualization site of the level-1
confirmed bit.

\begin{definition}[Stella-Octangula Level-1 Tile]
Let p be an
all-even-coord (level-2) site. The \emph{stella-octangula tile at p}
consists of:
\end{definition}

\begin{itemize}
\item
  The A-tetrahedron whose centroid is p: four sites in sublattice A at
  FCC nearest-neighbor positions to one another, surrounding p.
\end{itemize}

\begin{itemize}
\item
  The B-tetrahedron whose centroid is p: four sites in sublattice B at
  FCC nearest-neighbor positions to one another, surrounding p, related
  to the A-tetrahedron by inversion through p.
\end{itemize}

The union A-tet \ensuremath{\cup} B-tet at p consists of eight level-1 sites arranged as
a stella octangula (compound of two interpenetrating tetrahedra)
centered on the level-2 site p.

Stage 5 verifies this geometry explicitly through \ensuremath{L_{\mathrm{box}}} = 11 (see
Appendix E): A-tet centroids and B-tet centroids coincide at 100\% of
shared sites, with each pair forming an inversion partnership through
their common level-2 centroid.

\textbf{The minimum program configuration} consists of the
stella-octangula tile at p together with the boundary neighbors required
to evaluate the [[4,2,2]] stabilizer parity at each of the eight
level-1 sites. Each level-1 site within the tile has its own FCC
neighborhood within its sublattice; only some of these neighbors are
internal to the tile, and the remaining neighbors must be present in the
configuration to evaluate stabilizer parity. The boundary neighbors lie
on the same two sublattices and contribute to the program's spatial
footprint.

\subsubsection*{6.1.3 The Spatial Extent}

\begin{theorem}[Minimum Program Extent at Level 1]
The minimum
spatial extent of a stabilizer-satisfying T-closed configuration with
K(C) \textbf{\textless{}} \textbar C\textbar{} on the level-1
sublattice-bifurcated substrate is \ensuremath{L^{p(1)}} = \ensuremath{3\sqrt{2}} \ensuremath{\ell_{p}.}
\end{theorem}

\begin{proof}
We work in the storage convention of \S\,6.1.1, in which 1
storage unit equals \ensuremath{\ell_{p}/\sqrt{2}} --- fixed by Paper 1 \S\,VI, where the FCC
second-moment tensor over the 12 nearest neighbors of any agent
evaluates to \ensuremath{4\ell^{2}_{p}\delta_{\mu}\nu} and forces \ensuremath{\ell_{p}} to equal the level-0
nearest-neighbor distance (so \ensuremath{\ell_{p}} = \ensuremath{\sqrt{2}} storage units).

The minimum configuration is the inversion-partner stella-octangula tile
of \S\,6.1.2 --- 8 level-1 vertices comprising 4 on sublattice A and 4 on
sublattice B, all within the cube [1, \ensuremath{3]^{3}} in storage coords with
centroid (2, 2, 2) --- together with the boundary sites required to
evaluate the [[4,2,2]] stabilizer parity at each of the 8 tile
vertices.

For each tile vertex, 3 of its 12 in-sublattice nearest neighbors (at
storage distance \ensuremath{2\sqrt{2}} within the sublattice) lie inside the tile, namely
the other three vertices of the same tetrahedron. The remaining 9 per
vertex extend outside the tile. The cardinal range of the outward
extension is bounded by the most negative and most positive coordinates
these neighbors reach:

\begin{itemize}
\item
  From the A-tetrahedron vertex (1, 1, 3), the negative-direction
  sublattice-A neighbors include \ensuremath{(-1,} 1, 5), \ensuremath{(-1,} 3, 3), \ensuremath{(-1,} \ensuremath{-1,} 3),
  and \ensuremath{(-1,} 1, 1), all at storage distance \ensuremath{2\sqrt{2}} from (1, 1, 3) and all
  with coordinate-sum \ensuremath{\equiv} 1 mod 4 --- these reach coord \ensuremath{-1.}
\end{itemize}

\begin{itemize}
\item
  From the A-tetrahedron vertex (3, 3, 3), the positive-direction
  sublattice-A neighbors include (5, 5, 3), (5, 3, 5), and (3, 5, 5),
  all at storage distance \ensuremath{2\sqrt{2}} and sum \ensuremath{\equiv} 1 mod 4 --- these reach coord
  +5.
\end{itemize}

The cardinal bounding-box extent of the full configuration (tile plus
stabilizer-evaluation boundary) runs from coord \ensuremath{-1} to coord +5, i.e., 6
storage units along any cardinal axis. Converting to \ensuremath{\ell_{p}:}

\emph{\ensuremath{L^{p(1)}} = 6 storage units \ensuremath{\cdot} \ensuremath{(\ell_{p}/\sqrt{2}} per storage unit) = \ensuremath{(6/\sqrt{2})} \ensuremath{\ell_{p}} =
\ensuremath{3\sqrt{2}} \ensuremath{\ell_{p}} \ensuremath{\approx} 4.24 \ensuremath{\ell_{p}.}}

This sits at the lower edge of the structural range of 4--8 lattice
spacings posited in Paper 3 \S\,4.1, consistent with PGE forcing the
minimum at the lowest geometrically permissible scale.

The three program conditions are satisfied as follows.

\begin{itemize}
\item
  \emph{Stabilizer parity:} both A-tetrahedron and B-tetrahedron are in
  the +1 eigenspace of their respective level-1 [[4,2,2]]
  stabilizers; the joint configuration is the ground state of the
  compound stabilizer code under inversion symmetry through (2, 2, 2).
  Stabilizer parity is evaluable for each of the 8 tile vertices because
  all 12 in-sublattice nearest neighbors of each vertex lie within the
  configuration by construction.
\end{itemize}

\begin{itemize}
\item
  \emph{T-closure:} by the Recursion-Preservation Theorem of \S\,2.5, the
  level-1 T operator closes after a 2-flop cycle of substrate-scale
  confirmation steps. The first flop confirms a bit at the level-2
  centroid (the inversion center (2, 2, 2)) --- the unique level-2 site
  adjacent to at least three tile vertices, so the singular centroid is
  geometrically forced --- and the second flop re-distributes confirmation
  across the eight level-1 sites of the tile. This two-step cycle is the
  isolated closure floor of the configuration; the operative level-1 flop
  period \ensuremath{T_{p}^{(1)}} is set separately by the
  coupling-timescale constraint (\S\,6.2).
\end{itemize}

\begin{itemize}
\item
  \emph{K(C) \textbf{\textless{}} \textbar C\textbar:} the configuration
  is compressible because the inversion-partnership between
  A-tetrahedron and B-tetrahedron forces B from A via inversion through
  (2, 2, 2). The stella-octangula tile is \ensuremath{O_{h}}-invariant, so its
  orientation is a single orbit contributing \ensuremath{\log_2(1) = 0} bits
  (App.~B.4); specifying only the central [[4,2,2]] logical state (one of
  4, 2 bits) determines the entire tile, the boundary configuration being
  forced by the sublattice-bifurcation recursion. Total description length
  is \ensuremath{\log_2(4) = 2} bits, much smaller than the \textbar C\textbar{}
  \ensuremath{\approx} 17 bits of the raw eight-site polarization count.
\end{itemize}

Smaller candidates fail at least one condition.

\begin{itemize}
\item
  A configuration restricted to one sublattice cannot host the level-1
  [[4,2,2]] dynamics in its full form, because the level-2
  centroid at which the level-1 confirmed bit actualizes is reached only
  when both A-tet and B-tet are present. Stage 5 verifies this
  explicitly: shared A--B centroids occur at 100\% of inversion-partner
  pairs at every \ensuremath{L_{\mathrm{box}}} tested (Appendix E).
\end{itemize}

\begin{itemize}
\item
  A single tetrahedron without boundary is not closed --- no surround to
  evaluate.
\end{itemize}

\begin{itemize}
\item
  A configuration with cardinal extent fewer than 6 storage units in any
  direction cannot enclose the stella-octangula compound together with
  its boundary.
\end{itemize}

The \ensuremath{3\sqrt{2}} \ensuremath{\ell_{p}} extent is the geometric minimum at which all three program
conditions are jointly satisfied.

\end{proof}

\begin{corollary}[Minimum Subsystem Extent at Level 1]
\ensuremath{L^{sub(1)}} = 3 \ensuremath{L^{p(1)}} + \ensuremath{\delta} = \ensuremath{9\sqrt{2}} \ensuremath{\ell_{p}} + \ensuremath{\delta ,} with \ensuremath{\delta} \ensuremath{\leq} 2 \ensuremath{\ell_{p}} from the
meta-tetrahedral coupling overhead of Paper 3 \S\,5.3. Therefore \ensuremath{L^{sub(1)}}
\ensuremath{\in} \ensuremath{[9\sqrt{2}} \ensuremath{\ell_{p},} \ensuremath{9\sqrt{2}} \ensuremath{\ell_{p}} + 2 \ensuremath{\ell_{p}]} \ensuremath{\approx} [12.73 \ensuremath{\ell_{p},} 14.73 \ensuremath{\ell_{p}].}
\end{corollary}

\subsubsection*{6.1.4 Empirical Verification (Forward to Appendix E)}

The bifurcated structure underlying this construction is empirically
verified through Stages 5 and 6 of the supplementary simulation work
(Appendix E). The key results are:

\begin{itemize}
\item
  Each sublattice achieves the FCC coordination number 12 at interior
  sites, confirming the bifurcated geometry through \ensuremath{L_{\mathrm{box}}} = 11 in
  storage units (n = 5).
\end{itemize}

\begin{itemize}
\item
  The stabilizer-satisfying configuration space on each sublattice has
  dimension k = 3n over \ensuremath{F_{2},} where n = \ensuremath{(L_{\mathrm{box}}} \ensuremath{-} 1)/2 is the level-local
  lattice depth. This is the same scale-invariant k = 3n law that
  governs the substrate level (Paper 1 \S\,5.4, scale-invariant stabilizer
  null-space theorem).
\end{itemize}

\begin{itemize}
\item
  Tetrahedral count scales as \ensuremath{n^{3},} matching the bulk FCC count of unit
  cells in a volume of edge n.
\end{itemize}

\begin{itemize}
\item
  100\% of A-tet and B-tet centroids coincide at the same level-2 sites
  and form exact inversion partners through those centroids (verified at
  \ensuremath{L_{\mathrm{box}}} = 5, 7, 9, 11).
\end{itemize}

\begin{itemize}
\item
  The bifurcation structure repeats identically at level 2 (Stage 6),
  confirming the constant bifurcation factor 2 per level rather than a
  compounding factor. This is the empirical content of the structural
  correspondence between hierarchical levels established by the
  Recursion-Preservation Theorem.
\end{itemize}

\textbf{Status:} \emph{Derived} with empirical verification. The
sublattice bifurcation and the resulting \ensuremath{L^{p(1)}} = \ensuremath{3\sqrt{2}} \ensuremath{\ell_{p}} are derived
from FCC + [[4,2,2]] + recursive coarse-graining under PGE, with
the geometric closure structure (stella-octangula tile + boundary)
explicitly constructed and the cardinal bounding-box extent computed in
the storage frame. The empirical content is provided by Stages 5 and 6
of the simulation work, summarized in Appendix E.

\subsection*{6.2 The \ensuremath{c^{(k)}} Sequence at Saturation}

Paper 4 \S\,7.4 derived the upper bound \ensuremath{c^{(k+1)}} \ensuremath{\leq} \ensuremath{(1/\sqrt{2})} \ensuremath{c^{(k)}} from
the coupling-timescale constraint of \S\,5.3. At saturation, this is
realized with equality. We close the structural debt by explicit
minimum-period construction.

\begin{theorem}[\ensuremath{c^{(k)}} Sequence at Saturation]
At the
coupling-saturation bound, the hierarchical lightspeeds form the
geometric sequence
\end{theorem}

\begin{equation*}
c^{(k)} = (1/\sqrt{2})^{k} \cdot c
\end{equation*}

with c = \ensuremath{c^{(0)}} the substrate-level propagation speed. The
corresponding \ensuremath{L_{p}^{(k)}} and \ensuremath{T_{p}^{(k)}} sequences are

\begin{equation*}
L_{p}^{(k)} / L_{p}^{(k-1)} = R_{L}^{(k)}, T_{p}^{(k)} /
T_{p}^{(k-1)} = \sqrt{2} R_{L}^{(k)}
\end{equation*}

where \ensuremath{R_{L}^{(k)}} is the level-k minimum-program-extent ratio, with
\ensuremath{R_{L}^{(1)}} = \ensuremath{3\sqrt{2}} from \S\,6.1.3. The recursive scale invariance forces
\ensuremath{R_{L}^{(k)}} = \ensuremath{3\sqrt{2}} at every level k \ensuremath{\geq} 1 by the Recursion-Preservation
Theorem of \S\,2.5; the absolute level-k minimum extent is therefore
\ensuremath{L_{p}^{(k)}} = \ensuremath{(3\sqrt{2})^{k}} \ensuremath{\ell_{p}} in substrate units.

\emph{Construction.} Two temporal quantities must be distinguished. The
isolated closure floor of the minimum program is 2 substrate-scale flops
(the two-step cycle of \S\,6.1.2, verified in Appendix E). The operative
level-k flop period \ensuremath{T_{p}^{(k)}} is the larger of this floor
and the coupling-timescale constraint of Paper 4 \S\,5.3; the constraint
dominates, giving \ensuremath{T_{p}^{(1)}} = \ensuremath{\sqrt{2}}
\ensuremath{\cdot} \ensuremath{(3\sqrt{2})} \ensuremath{\cdot} \ensuremath{t_{p}} = 6
\ensuremath{t_{p}} at saturation. At each hierarchical level the same
structure recurs: the level-k program closes after a 2-flop cycle of
\ensuremath{\text{level}-(k-1)}-scale steps, and the coupling constraint
sets the operative period. The minimum
\ensuremath{L_{p}^{(k)}} is set by the level-k minimum tetrahedral configuration plus
its boundary, which has the same structure as the substrate-level
construction of \S\,6.1 scaled by \ensuremath{R_{L}^{(k)}.} At the coupling-saturation
bound, \ensuremath{T_{p}^{(k+1)}} = \ensuremath{\sqrt{2}} \ensuremath{(L_{p}(k+1)/L_{p}(k))} \ensuremath{T_{p}^{(k)},} giving the
saturation relation

\emph{\ensuremath{c^{(k+1)}} = \ensuremath{(L_{p}}\textbf{\ensuremath{(k+1)/L_{p}}}(k)) \ensuremath{\cdot}
\ensuremath{(T_{p}}\textbf{\ensuremath{(k)/T_{p}}}(k+1)) \ensuremath{\cdot} \ensuremath{c^{(k)}} = \ensuremath{(1/\sqrt{2})} \ensuremath{c^{(k)}.} \hfill\ensuremath{\square}}

The explicit \ensuremath{c^{(k)}} sequence in substrate units, normalized to c =
\ensuremath{c^{(0)}} = \ensuremath{c_{\mathrm{TMS}},} is given in the following table.

\begin{longtable}[]{@{}
  >{\raggedright\arraybackslash}p{(\columnwidth - 8\tabcolsep) * \real{0.2000}}
  >{\raggedright\arraybackslash}p{(\columnwidth - 8\tabcolsep) * \real{0.2000}}
  >{\raggedright\arraybackslash}p{(\columnwidth - 8\tabcolsep) * \real{0.2000}}
  >{\raggedright\arraybackslash}p{(\columnwidth - 8\tabcolsep) * \real{0.2000}}
  >{\raggedright\arraybackslash}p{(\columnwidth - 8\tabcolsep) * \real{0.2000}}@{}}
\toprule\noalign{}
\begin{minipage}[b]{\linewidth}\raggedright
\textbf{Level k}
\end{minipage} & \begin{minipage}[b]{\linewidth}\raggedright
\textbf{\ensuremath{L^{p(k)}} / \ensuremath{\ell_{p}}}
\end{minipage} & \begin{minipage}[b]{\linewidth}\raggedright
\textbf{\ensuremath{T^{p(k)}} / \ensuremath{t_{p}}}
\end{minipage} & \begin{minipage}[b]{\linewidth}\raggedright
\textbf{\ensuremath{c^{(k)}/c} (exact)}
\end{minipage} & \begin{minipage}[b]{\linewidth}\raggedright
\textbf{\ensuremath{c^{(k)}/c} (decimal)}
\end{minipage} \\
\midrule\noalign{}
\endhead
\bottomrule\noalign{}
\endlastfoot
0 & 1 & 1 & 1 & 1.000 \\
1 & \ensuremath{3\sqrt{2}} \ensuremath{\approx} 4.24 & 6 & \ensuremath{1/\sqrt{2}} & 0.707 \\
2 & 18 & 36 & 1/2 & 0.500 \\
3 & \ensuremath{54\sqrt{2}} \ensuremath{\approx} 76.37 & 216 & \ensuremath{1/(2\sqrt{2})} & 0.354 \\
4 & 324 & 1296 & 1/4 & 0.250 \\
5 & \ensuremath{972\sqrt{2}} \ensuremath{\approx} 1374.6 & 7776 & \ensuremath{1/(4\sqrt{2})} & 0.177 \\
\end{longtable}

where \ensuremath{T^{p(k)}} = \ensuremath{(\sqrt{2}} \ensuremath{\cdot} \ensuremath{R_{L})_{k}} \ensuremath{\cdot} \ensuremath{t_{p}} = \ensuremath{6_{k}} \ensuremath{\cdot} \ensuremath{t_{p}} (since \ensuremath{\sqrt{2}} \ensuremath{\cdot} \ensuremath{3\sqrt{2}} = 6).

These are upper bounds at saturation. In non-saturated regimes, \ensuremath{c^{(k)}}
can be smaller than the saturation value. The sequence forms a discrete
dispersion spectrum bounded above by the saturation cascade.

\subsection*{6.3 Identifying Our Universe's Level}

\subsubsection*{6.3.0 The Level-Advance Relation}

The hierarchical structure of TMS substrate requires a precise statement
of what advances the level counter k. Without such a statement, the
cross-level ratios of \S\,6.1.3 and \S\,6.2 and the level-identification
framework of \S\,7.6 use k as an undefined parameter. We give the
definition here and verify its consistency with the recursion structure
of \S\,2.5.

\begin{definition}[Level-Advance Relation]
Let \ensuremath{O^{(k)}} denote
a primitive substrate operation at level k --- a confirmation event in
the sense of Paper 1 \S\,III, taking three agents at level k as inputs and
producing one resolved bit at level k. Level k+1 begins at the scale at
which a minimum stabilizer-satisfying T-closed configuration with K(C)
\textless{} \textbar C\textbar{} forms among level-k confirmed bits
acting as agents in a new primitive operation \ensuremath{O^{(k+1)}} of the same
structural form.
\end{definition}

Three load-bearing features:

\begin{enumerate}
\def\labelenumi{(\roman{enumi})}
\item
  \emph{Same structural form.} Each level's primitive operation must be
  a confirmation event satisfying \ensuremath{\bar{B}} = \ensuremath{3\bar{\alpha}.} This excludes ``more complex
  configurations of level-k events'' from counting as level-(k+1) ---
  only genuine role-shifts where outputs become new agents in a new
  triadic closure advance k.
\end{enumerate}

\begin{enumerate}
\def\labelenumi{(\roman{enumi})}
\item
  \emph{Single bit per operation.} Level advancement is a structural
  role-shift, not a numerical accumulation. The level counter does not
  advance because more bits get resolved at level k; it advances because
  the resolved bits begin functioning as agents in a structurally new
  operation.
\end{enumerate}

\begin{enumerate}
\def\labelenumi{(\roman{enumi})}
\item
  \emph{Quantization-by-closure criterion.} The ``scale at which'' a new
  level begins is operationalized through \S\,3.5's
  quantization-by-closure: level k+1 is realized when the minimum K(C)
  \textless{} \textbar C\textbar{} configuration successfully forms.
  This makes the criterion substrate-internal --- it does not appeal to
  an external observer making the scale comparison.
\end{enumerate}

\textbf{Consistency with \S\,2.5 (Recursion-Preservation).} The
Recursion-Preservation Theorem states that confirmation-event structure
is preserved under level transitions. The Level-Advance Relation is the
formal statement of \emph{how} it is preserved: by role-shift, where the
bit-role at level k becomes the agent-role at level k+1.

\textbf{Consistency with \S\,6.1.3 (length scaling).} Under the
Level-Advance Relation, the geometric extent of a level-(k+1) operation
is determined by the spacing of level-k confirmed bits (now functioning
as level-(k+1) agents). This spacing inherits the stella-octangula tile
structure from level k, producing \ensuremath{L_{p}}\textsuperscript{\ensuremath{(k+1)/L_{p}}}(k) =
\ensuremath{3\sqrt{2}} as derived in \S\,6.1.3.

\textbf{Consistency with \S\,5.8 (reinterpretation cascades).} The
reinterpretation operator \ensuremath{\hat{R}} at substrate exhaustion performs a
level-advance at cosmological scale: level-k confirmed bits at
exhaustion become the agents of the \ensuremath{\text{post}-\hat{R}} substrate. \ensuremath{\hat{R}} therefore
advances k by exactly 1, not by the within-cycle level count and not by
resetting.

\begin{corollary}
The total level identity of an observer in TMS
is the sum of (a) within-cycle role-shifts in the current cycle's
history, plus (b) the number of \ensuremath{\hat{R}-\text{firings}} in the cosmic history leading
to the current cycle.
\end{corollary}

\subsubsection*{6.3.1 Inferring Our Level from Observation}

If our universe is itself a level-k structure for some k determined by
the chain of reinterpretation events that produced our substrate, then
our measured speed of light \ensuremath{c_{\mathrm{measured}}} is \ensuremath{c^{(k)}} in substrate units.
The substrate \ensuremath{c_{\mathrm{TMS}}} is obtained by inverting the \ensuremath{c^{(k)}} formula:

\begin{equation*}
c_{\mathrm{TMS}} = (\sqrt{2})^{k} \cdot c_{\mathrm{measured}}
\end{equation*}

at the saturation bound. For k = 1 (one reinterpretation step removed
from the substrate), \ensuremath{c_{\mathrm{TMS}}} = \ensuremath{\sqrt{2}} \ensuremath{\cdot} \ensuremath{c_{\mathrm{measured}}} \ensuremath{\approx} 4.24 \ensuremath{\times} \ensuremath{10^{8}} m/s. For k =
2, \ensuremath{c_{\mathrm{TMS}}} = 2 \ensuremath{\cdot} \ensuremath{c_{\mathrm{measured}}} \ensuremath{\approx} 6 \ensuremath{\times} \ensuremath{10^{8}} m/s. Higher k gives proportionally
larger substrate-level c.

The actual value of k for our universe is a structural prediction whose
empirical pinning requires comparison with multiple classes of
observable. GUP scaling \ensuremath{\beta} depends on k (\S\,6.4 and \S\,VIII below); precision
measurements of GUP corrections at varying preparation depths would in
principle constrain k. Cross-level relativistic effects (\S\,VII) produce
specific signatures whose magnitude depends on the local hierarchical
structure. Vacuum-energy and dark-sector observables (\S\,7.6.4) constitute
a further empirical handle. All are accessible to next-generation
precision experiments in principle.

Our universe's k is therefore a derived structural quantity whose value
is pinned by experimental signatures the framework predicts.

\subsection*{6.4 The Closed Form of f}

Paper 4 \S\,8.1 specified the functional dependences of the refined GUP
parameter:

\begin{equation*}
\beta = \beta_{0} \cdot f(d, K_{\psi,\mathrm{local}}, K_{\psi,\mathrm{surround}}, K_{\psi,\mathrm{subsystem}}, k)
\end{equation*}

with f monotonically increasing in all five arguments. We close the
structural debt by deriving the closed form under additive Kolmogorov
bounds.

\begin{theorem}[Closed Form of f]
Under the additive Kolmogorov
bound, where each algorithmic complexity contributes independently to
the substrate footprint, the closed form is
\end{theorem}

\emph{f(d, \ensuremath{K_{\mathrm{local}},} \ensuremath{K_{\mathrm{surround}},} \ensuremath{K_{\mathrm{subsys}},} k) = 1 + \ensuremath{\alpha_{d}} \ensuremath{\cdot} d +
\ensuremath{\alpha_{\mathrm{local}}} \ensuremath{\cdot} \ensuremath{K_{\mathrm{local}}} + \ensuremath{\alpha_{\mathrm{surround}}} \ensuremath{\cdot} \ensuremath{K_{\mathrm{surround}}} + \ensuremath{\alpha_{\mathrm{subsys}}} \ensuremath{\cdot} \ensuremath{K_{\mathrm{subsys}}}
+ \ensuremath{\alpha_{k}} \ensuremath{\cdot} k}

where each coefficient \ensuremath{\alpha_{X}} is the per-unit weight of its corresponding
argument's contribution to the bit's substrate footprint. The
coefficients are dimensionless geometric constants determined by the FCC
lattice connectivity and the [[4,2,2]] code distance;
specifically:

\begin{equation*}
\alpha_{d} = 1, \alpha_{\mathrm{local}} = 1, \alpha_{\mathrm{surround}} = \sqrt{2}, \alpha_{\mathrm{subsys}} = 3, \alpha_{k} =
(1/\sqrt{2})^{k}
\end{equation*}

where the per-level structural overhead of meta-tetrahedral coupling,
written \ensuremath{\delta_{k}} in earlier drafts, is derived to be identically one by
the form-invariance fixed point of Paper~3 \S\,4.7 (the meta-tetrahedral
geometry is the same at every level, so its per-level ratio is one) and is
therefore absorbed.

\emph{Proof sketch.} The substrate footprint of a bit is the lattice
volume over which its causal record extends. Each of the five arguments
contributes additively to this volume in the limit of small
contributions (independent contributions from each axis). The causal
depth d contributes one lattice unit per flop of history. The local
algorithmic complexity \ensuremath{K_{\mathrm{local}}} contributes one lattice unit per bit of
local algorithmic structure. The surround complexity \ensuremath{K_{\mathrm{surround}}}
contributes \ensuremath{\sqrt{2}} per bit of imprinting history (weighted by the \ensuremath{\sqrt{2}}
imprint-depth factor of Paper 3 \S\,2.3). The subsystem complexity
\ensuremath{K_{\mathrm{subsys}}} contributes 3 per bit of program content (weighted by the
3-component \ensuremath{(C_{D},} \ensuremath{C_{O},} \ensuremath{C_{C})} trichotomy). The hierarchical depth k
contributes \ensuremath{(1/\sqrt{2})^{k}} per level due to the \ensuremath{c^{(k)}} dispersion
(the meta-tetrahedral coupling overhead being unity by scale-invariance). The closed form
follows by combining these contributions additively, with each \ensuremath{\alpha_{X}} set
by the corresponding geometric coupling. \hfill\ensuremath{\square}

The full derivation, including the multiplicative corrections beyond the
additive bound, is given in Appendix C. The structural result here is
sufficient: f is linear in its arguments under the additive Kolmogorov
bound, with coefficients determined by geometric quantities. The
empirical content of the prediction is that bits in deeper, more
complex, more deeply hierarchically embedded configurations exhibit
measurably larger GUP corrections; the linear functional form provides
specific quantitative predictions accessible to precision GUP
phenomenology (Hossenfelder, 2013).

Beyond the additive limit --- when the contributions become large enough
that cross-terms dominate --- f acquires multiplicative structure. The
additive form is the leading-order expression; the next-order correction
is detailed in Appendix C and reserved for simulation work for full
numerical pinning.

\section*{VII. Cross-Level Relativistic Dynamics: The Substrate-Level Origin of Gravity}
\addcontentsline{toc}{section}{VII. Cross-Level Relativistic Dynamics: The Substrate-Level Origin of Gravity}

\subsection*{7.1 Three Cross-Level Phenomena}

Paper 4 \S\,8.2 flagged three structural effects of \ensuremath{c^{(k)}} variation
across levels when subsystems at different hierarchical levels share
substrate at their interfaces: variable apparent propagation,
frame-relative timing, and embedded-medium dilation. We derive each here
as substrate-level mechanisms producing the kinematical and
gravitational phenomena of general relativity in the continuum limit.

\subsection*{7.2 Variable Apparent Propagation: The Origin of Gravitational Redshift}

\begin{theorem}[Variable Apparent Propagation]
When a signal
propagates from a region of \ensuremath{\text{level}-(k-1)} hierarchical depth into a region
of level-k hierarchical depth, its apparent propagation speed shifts
from \ensuremath{c^{(k-1)}} to \ensuremath{c^{(k)}} at the interface, with \ensuremath{c^{(k)}} = \ensuremath{(1/\sqrt{2})}
\ensuremath{c^{(k-1)}} at the saturation bound.
\end{theorem}

\begin{proof}
The signal propagates at the local lattice speed of the
level it occupies --- by the lattice-relative invariance of the wave
equation (Paper 3 \S\,4.5), this is \ensuremath{1/\sqrt{2}} in level-local units. In substrate
units, this corresponds to \ensuremath{c^{(k)}} for level k and \ensuremath{c^{(k-1)}} for level
\ensuremath{(k-1).} At the interface between levels, the signal transitions from one
substrate to the other, and its apparent propagation speed in any global
coordinate system reflects the level it currently occupies. The
transition is geometric, not dynamical: no energy is added or removed;
only the local propagation rate changes because the local substrate's
lattice structure differs.

\end{proof}

The physical interpretation: when a photon-like signal propagates
outward from a deeply hierarchically embedded subsystem (a ``mass'' in
\ensuremath{\text{level}-(k-1)} substrate) into shallower substrate (level k with smaller
hierarchical embedding), it appears to slow down. The slowing is what an
external observer measures as gravitational redshift. In the TMS
picture, redshift is not a frequency shift of the photon but a change in
the local propagation speed of the substrate medium through which the
photon propagates.

The factor \ensuremath{1/\sqrt{2}} in the cross-level propagation ratio is not an empirical
bound but a structural geometric law of the FCC lattice. By the FCC
second-moment tensor \ensuremath{4\ell^{2}_{p}\delta} over the 12 nearest neighbors (Paper 1 \S\,VI),
the lattice's effective propagation rate is the face-diagonal-to-edge
ratio of the cubic cell containing the FCC, which is exactly \ensuremath{1/\sqrt{2}.} By
recursive scale invariance and the Recursion-Preservation Theorem of
\S\,2.5, the same face-diagonal-to-edge ratio reappears at every
hierarchical level. The \ensuremath{c^{(k+1)}} \ensuremath{\leq} \ensuremath{(1/\sqrt{2})} \ensuremath{c^{(k)}} bound is therefore a
\emph{structural law} of the FCC + recursive coarse-graining geometry,
not an empirical fit. At saturation (the coupling-timescale bound of
Paper 4 \S\,5.3), the bound is realized with equality at every level;
non-saturated regimes have \ensuremath{c^{(k+1)}} \textless{} \ensuremath{(1/\sqrt{2})} \ensuremath{c^{(k)}} with the
deviation controlled by the local hierarchical embedding.

This is structurally consistent with Volovik's Fermi-point scenarios
(Volovik, 2007), in which gravity emerges from defect structures in an
underlying lattice substrate, and with the general entropic-gravity
program (Verlinde, 2011; Bianconi, 2025), in which gravitational
phenomena emerge from information-theoretic properties of the underlying
substrate. The TMS provides a particularly clean instance: the
gravitational redshift is the variable apparent propagation across
hierarchical levels, with the level structure forced by the recursive
coarse-graining of Paper 3.

\subsection*{7.3 Frame-Relative Timing: The Origin of Time Dilation}

\begin{theorem}[Frame-Relative Timing]
Events that are simultaneous
in \ensuremath{\text{level}-(k-1)} flop time \ensuremath{T_{p}^{(k-1)}} may not be simultaneous in
level-k flop time \ensuremath{T_{p}^{(k)},} because the flop periods at different
hierarchical levels are independent. Two events at level-mixed locations
have ordering that depends on the level-frame of the observer.
\end{theorem}

\begin{proof}
By the \ensuremath{c^{(k)}} dispersion and the coupling-timescale
constraint of Paper 4 \S\,5.3, \ensuremath{T_{p}^{(k)}} is independent of \ensuremath{T_{p}^{(k-1)}}
modulo the saturation bound. An event at a \ensuremath{\text{level}-(k-1)} location
occurring at \ensuremath{\text{level}-(k-1)} flop time \ensuremath{\tau_{a}^{(k-1)}} and an event at a
level-k location occurring at level-k flop time \ensuremath{\tau_{b}^{(k)}} can both be
assigned absolute substrate times. If the substrate times coincide, the
events are simultaneous in absolute substrate time. But their
level-local timing depends on the level-frame: an observer at level
\ensuremath{(k-1)} sees them as occurring at \ensuremath{\text{level}-(k-1)} units that differ from the
level-k units by the \ensuremath{T_{p}(k)/T_{p}(k-1)} ratio. The ordering depends on
the level-frame.

\end{proof}

The physical interpretation: time dilation between observers in
different gravitational potentials (different hierarchical-level
configurations in TMS terms) follows from the difference in flop periods
at their respective levels. An observer in a deeper
hierarchically-embedded configuration experiences slower flop time
relative to a shallower observer; this is the substrate-level mechanism
of gravitational time dilation.

Standard general relativity describes this through the metric-tensor
coupling of time to gravitational potential. The TMS provides the
substrate-level origin: gravitational potential is a measure of local
hierarchical embedding depth, and time dilation is the difference in
flop periods at different hierarchical levels.

\subsection*{7.4 Embedded-Medium Dilation: The Mass-Geometry Coupling}

\begin{theorem}[Embedded-Medium Dilation]
A \ensuremath{\text{level}-(k-1)} subsystem
deeply embedded in a level-k meta-subsystem experiences propagation
through the meta-subsystem's substrate at the effective \ensuremath{c^{(k)}} of the
meta-medium, producing an apparent dilation of its internal dynamics
relative to the bare \ensuremath{c^{(k-1)}.}
\end{theorem}

\emph{Proof sketch.} When a \ensuremath{\text{level}-(k-1)} subsystem is embedded inside a
level-k meta-subsystem, the substrate at the subsystem's location is not
pure \ensuremath{\text{level}-(k-1)} substrate --- it is level-k substrate hosting a
\ensuremath{\text{level}-(k-1)} configuration. Wave propagation within the \ensuremath{\text{level}-(k-1)}
subsystem occurs through the level-k medium, with apparent speed
\ensuremath{c^{(k)}} rather than the bare \ensuremath{c^{(k-1)}} the subsystem would experience
in isolation. The dilation factor is the ratio \ensuremath{c(k-1)/c(k),} which at
saturation is \ensuremath{\sqrt{2}.} \hfill\ensuremath{\square}

The physical interpretation: a massive object --- a deeply embedded
subsystem at lower hierarchical level --- experiences slowed internal
dynamics relative to its bare clock because its substrate is the level-k
medium of the enclosing meta-subsystem. The slowing corresponds to the
rest-mass energy in the field-theoretic continuum limit: the energy
stored in the embedded configuration is the energy required to support
its \ensuremath{\text{level}-(k-1)} operation within a level-k substrate.

This is the structural origin of mass-energy equivalence in the TMS
picture: mass is energy that has been laid into the wake. A massive
excitation is a standing, self-re-confirming program (\S\,3.5) --- it
re-confirms itself in place with a rest period \ensuremath{T_{0},} giving it an
intrinsic rest frequency \ensuremath{\omega_{0}} = \ensuremath{2\pi /T_{0}} even at zero momentum. Its rest
energy is the activity of that re-confirmation, \ensuremath{E_{0}} = \ensuremath{\hbar \omega_{0},} and its mass
is that rest energy in inertial units, m = \ensuremath{\hbar \omega_{0}/c^{2}.} The conversion factor
\ensuremath{c^{2}} here is the within-level substrate wave speed of Paper 1 \S\,6.5 --- the
same c that propagates every excitation --- not the cross-level \ensuremath{\sqrt{2}}
ratio. This is the crucial disentanglement: the within-level \ensuremath{c^{2}} performs
the mass-energy conversion (E = \ensuremath{mc^{2}),} while the cross-level \ensuremath{\sqrt{2}-\text{per}-\text{level}}
ratio performs gravity (the cross-level dispersion of \S\,7.5). One
geometric primitive --- the FCC face-diagonal --- plays two distinct
roles, within a level and across levels.

\begin{corollary}[Relativistic Energy--Momentum Relation]
A
standing excitation with rest frequency \ensuremath{\omega_{0}} propagating on the FCC
substrate obeys, in the continuum limit, the dispersion \ensuremath{\omega^{2}(k)} = \ensuremath{\omega^{2}_{0}} +
\ensuremath{c^{2}}\textbar k\textbar\ensuremath{^{2},} with c the universal, mass-independent, isotropic
substrate speed. Identifying E = \ensuremath{\hbar \omega ,} p = \ensuremath{\hbar k,} and m = \ensuremath{\hbar \omega_{0}/c^{2}} yields \ensuremath{E^{2}} =
\ensuremath{(pc)^{2}} + \ensuremath{(mc^{2})^{2}} --- the relativistic energy--momentum relation.
\end{corollary}

This relation is not an additional postulate; it is a corollary of four
results already established in the volume. The second-order substrate
dynamics (Paper 1 \S\,6.3) force the quadrature form \ensuremath{\omega^{2}} = \ensuremath{\omega^{2}_{0}} + \ensuremath{c^{2}k^{2}} rather
than a non-relativistic \ensuremath{\omega} = \ensuremath{\omega_{0}} + \ensuremath{k^{2}/2m.} The isotropic second-moment
tensor (Paper 1 \S\,6.5) forces c to be universal and direction-independent
--- the same c for every mass. Quantization-by-closure (\S\,3.5) supplies
the rest frequency \ensuremath{\omega_{0}} of a standing program. The PGE-forced wave-form
theorem (Paper 2 \S\,VII) forces every excitation to share the one
substrate wave operator, so no mass modifies c. The relation holds
exactly in the continuum limit; its leading correction is the
\ensuremath{(E/E_{\mathrm{Planck}})^{2}} Lorentz violation of \S\,4.1 Lemma B, entering at the same
quartic-in-momentum \ensuremath{O_{h}} invariant. Group velocity is zero at rest,
rises with momentum, and remains below c, so propagation is causal.
Status: the kinematic form is Derived; the rest-frequency spectrum \ensuremath{\omega_{0}}
--- the actual mass values --- is the mass-spectrum problem deferred to
\S\,8.4.3.

\subsection*{7.5 Gravity as Cross-Level Dispersion: The Unified Thesis}

The three cross-level phenomena of \S\,7.2--\S\,7.4 --- variable apparent
propagation, frame-relative timing, embedded-medium dilation --- are not
three separate phenomena. They are three manifestations of the same
underlying substrate mechanism: the dispersion of effective propagation
speed \ensuremath{c^{(k)}} across hierarchical levels, with the dispersion rooted in the
recursive FCC face-diagonal-to-edge ratio of \S\,7.2.

\textbf{The unified thesis.} Gravity is not a fundamental force in the
TMS picture. It is the substrate-level expression of a single structural
fact: subsystems at different hierarchical levels of embedding
experience different effective propagation speeds, because each level's
FCC + [[4,2,2]] substrate has its own local \ensuremath{c^{(k)}.} The standard
phenomena of general relativity all reduce to manifestations of this
cross-level dispersion:

\begin{itemize}
\item
  \emph{Gravitational redshift (\S\,7.2)} is the change in apparent
  propagation speed when a signal crosses between substrate levels. A
  photon-like signal moving outward from a deeply hierarchically
  embedded mass into shallower substrate appears to slow because the
  local \ensuremath{c^{(k)}} changes by a factor \ensuremath{\sqrt{2}} per level transition.
\end{itemize}

\begin{itemize}
\item
  \emph{Gravitational time dilation (\S\,7.3)} is the difference in flop
  periods \ensuremath{T^{p(k)}} between substrate levels. An observer in deeper
  hierarchical embedding experiences slower local flop time relative to
  a shallower observer, because \ensuremath{T^{p(k)}} increases by a factor of \ensuremath{\sqrt{2}} \ensuremath{\cdot}
  \ensuremath{R_{L}} = 6 per level (with \ensuremath{R_{L}} = \ensuremath{3\sqrt{2}} from \S\,6.1.3).
\end{itemize}

\begin{itemize}
\item
  \emph{Mass-geometry coupling (\S\,7.4)} is the embedded-medium dilation
  when a lower-level subsystem operates within a higher-level substrate.
  The energy required to support \ensuremath{\text{level}-(k-1)} dynamics within a level-k
  medium is the rest-mass energy of the embedded configuration. Mass is
  hierarchical embedding depth, and gravitational binding is the
  substrate cost of maintaining the cross-level embedding.
\end{itemize}

\textbf{Why this unifies as a thesis.} Standard general relativity
expresses the matter-geometry coupling through the Einstein equations,
which couple stress-energy to the metric. In the TMS picture, both sides
of this coupling are the same underlying object viewed at different
scales: the matter side is the embedded subsystem at lower hierarchical
level, and the geometry side is the host substrate at higher
hierarchical level. The ``coupling'' is just the embedding relation,
with the coupling strength set by the cross-level \ensuremath{c^{(k)}} dispersion. This
is why the Newton constant G has a structurally simple form in the TMS
reading (\S\,7.6): G is the ratio \ensuremath{(L^{p(1)}/L^{p(0)})} / \ensuremath{(T^{p(1)}/T^{p(0)})^{2}} =
\ensuremath{(3\sqrt{2})/36,} in substrate-relative units, modulo the dimensional conversion
factors.

\textbf{The substrate-level origin of gravity.} Gravity emerges in the
TMS picture not as an additional force or field but as the structural
consequence of having a hierarchical FCC substrate with recursive
coarse-graining under PGE. The hierarchical structure exists because PGE
selects the minimum sufficient closure at every scale (\S\,2.3); the
cross-level \ensuremath{c^{(k)}} dispersion exists because the FCC face-diagonal-to-edge
ratio is \ensuremath{1/\sqrt{2}} at every level (\S\,7.2); the matter-geometry coupling exists
because lower-level subsystems can only operate by being embedded in
higher-level substrate (\S\,7.4). Take any one of these structural
ingredients away and the substrate either has no hierarchy or no
gravitational phenomena. Gravity is what the TMS substrate exhibits when
its hierarchical structure is viewed from inside one of its levels.

\textbf{Status of \S\,7.5:} \emph{Derived} (the unified thesis). The
structural identification of gravity as cross-level dispersion follows
from \S\,7.2, \S\,7.3, \S\,7.4 jointly under recursive scale-invariance.

\subsection*{7.6 Cross-Level Gravitational Coupling and the Level-Identification Problem}

The three cross-level phenomena of \S\,7.2--\S\,7.5 --- variable apparent
propagation, frame-relative timing, and embedded-medium dilation ---
combine in the continuum limit to produce a coupling between matter
\ensuremath{(\text{level}-(k-1)} subsystems embedded in level-k substrate) and the tensor
sector of the polarization field. This coupling has the structural form
of the Einstein equations. The structural ratio characterizing this
coupling between successive levels is forced by substrate recursion. The
SI numerical value of Newton's constant in our universe depends on which
level we sit at; this level identity is empirically determined.

\subsubsection*{7.6.1 The Structural Skeleton}

The cross-level acceleration-like ratio characterizing gravitational
coupling between successive levels of the substrate hierarchy is
determined by three derived quantities: the length-ratio
\ensuremath{L_{p}}\textsuperscript{\ensuremath{(k+1)/L_{p}}}(k), the time-ratio
\ensuremath{T_{p}}\textsuperscript{\ensuremath{(k+1)/T_{p}}}(k), and the per-level lightspeed
\ensuremath{c^{(k)}.} From \S\,6.1.3, the length-ratio is forced to
\ensuremath{L_{p}}\textsuperscript{\ensuremath{(k+1)/L_{p}}}(k) = \ensuremath{3\sqrt{2}} by the stella-octangula tile
geometry and the FCC second-moment storage convention (1 storage unit =
\ensuremath{\ell_{p}/\sqrt{2}).}

From \S\,5.3, the coupling-timescale constraint requires that three level-k
subsystems coherently meta-confirm at level (k+1). Outputs propagate at
\ensuremath{c^{(k)},} and the meta-tetrahedral edge-to-interstice distance in
level-k units is \ensuremath{\sqrt{2}} \ensuremath{\cdot} \ensuremath{L_{p}}\textsuperscript{\ensuremath{(k+1)/L_{p}}}(k). For coherent
simultaneous arrival of all three signals at the meta-interstice, the
time-ratio must satisfy \ensuremath{T_{p}}\textsuperscript{\ensuremath{(k+1)/T_{p}}}(k) \ensuremath{\geq} \ensuremath{\sqrt{2}} \ensuremath{\cdot}
\ensuremath{(L_{p}}\textsuperscript{\ensuremath{(k+1)/L_{p}}}(k)).

This is a lower bound; the substrate could in principle use a larger
time-ratio (slower meta-confirmation). By the Principle of Generative
Economy (Axiom 0), no unmotivated waiting between closures is admissible
--- the substrate selects the minimum-sufficient closure. The time-ratio
therefore saturates: \ensuremath{T_{p}}\textsuperscript{\ensuremath{(k+1)/T_{p}}}(k) = \ensuremath{\sqrt{2}} \ensuremath{\cdot} \ensuremath{3\sqrt{2}} =
\textbf{6}. Correspondingly, the lightspeed ratio saturates the
geometric bound: c\textsuperscript{(k+1)/c}(k) = \ensuremath{(3\sqrt{2})} / 6 =
\textbf{\ensuremath{1/\sqrt{2}}}.

The cross-level acceleration-like ratio that characterizes gravitational
coupling at the substrate level is then:

\ensuremath{(L_{p}}\textsuperscript{\ensuremath{(k+1)/L_{p}}}(k)) /
\ensuremath{(T_{p}}\textsuperscript{\ensuremath{(k+1)/T_{p}}}\ensuremath{(k))^{2}} = \ensuremath{3\sqrt{2}} / 36 = \ensuremath{\sqrt{2}/12} \ensuremath{\approx} 0.118

This is Newton's gravitational constant expressed as a cross-level
structural ratio. The reduction is striking: the entire skeleton reduces
to one geometric primitive (the FCC face-diagonal-to-edge ratio \ensuremath{\sqrt{2})} plus
the Principle of Generative Economy plus recursion. No dimensional
analysis, no free constants, no tuning.

\subsubsection*{7.6.2 Scope of the Substrate Prediction}

The framework predicts: (i) the cross-level ratios L = \ensuremath{3\sqrt{2},} T = 6, c =
\ensuremath{1/\sqrt{2}} between any two successive levels; (ii) level-internal
self-consistency: at each level k, the local Planck triple \ensuremath{(\ell_{p}^{(k)},}
\ensuremath{t_{p}^{(k)},} \ensuremath{m_{p}^{(k)})} satisfies the standard Planck relations using
\ensuremath{c^{(k)}} and \ensuremath{\hbar_{\mathrm{TMS}}} = 1 atomic action quantum per substrate flop; (iii)
mass-ratio \ensuremath{M_{p}}\textsuperscript{\ensuremath{(k+1)/M_{p}}}(k) = \ensuremath{3\sqrt{2}/2,} forced by
requiring local G-consistency at every level.

The absolute SI value of G in our universe depends on which level of the
substrate hierarchy our universe occupies --- the level-identification
problem.

\subsubsection*{7.6.3 The Level-Identification Problem}

If our universe occupies level k, then the substrate-level Planck length
satisfies \ensuremath{\ell_{p,\mathrm{substrate}}} = \ensuremath{\ell_{p,\mathrm{observed}}} / \ensuremath{(3\sqrt{2})^{k}.} Sample magnitudes:
at k = 1, substrate \ensuremath{\ell_{p}} \ensuremath{\approx} 3.8 \ensuremath{\times} \ensuremath{10^{-36}} m; at k = 2, substrate \ensuremath{\ell_{p}} \ensuremath{\approx} 9.0 \ensuremath{\times}
\ensuremath{10^{-37}} m. The framework predicts the structural form of physical law; the
empirical content lives in the discrete question of level identity,
whose answer determines all level-dependent quantities simultaneously.

\subsubsection*{7.6.4 Empirical Routes Toward Level Identification}

Several candidate empirical handles on k are flagged here for future
work.

\emph{CMB power spectrum signature.} The level-1 saturation cascade
(\S\,5.1--\S\,5.2) produces a discrete-substrate fingerprint at small angular
scales. The transition scale depends on \ensuremath{(3\sqrt{2})^{k},} so the position of
the fingerprint in the CMB power spectrum is a direct probe of k.

\emph{Hubble constant and dark-energy density.} These quantities scale
with cosmological time, which itself scales with \ensuremath{t_{p}^{(k)}} through the
saturation cascade. Their ratio to local atomic clock rates carries
level information.

\emph{Gravitational-wave dispersion.} If \ensuremath{c^{(k)}} has cross-level
structure, gravitational waves approaching substrate scale should show
level-dependent dispersion. Current LIGO/Virgo frequencies are far below
substrate scale, but future high-frequency detectors could probe it.

\emph{Neutrino mass hierarchy.} The three generations occupy distinct
positions in the recursive embedding structure (cf.~\S\,8.4). If their mass
eigenvalues encode the recursion depth at which each generation's
fermionic content is anchored, the observed neutrino mass hierarchy
\ensuremath{(\Delta m^{2}_{21}} \ensuremath{\ll} \textbar \ensuremath{\Delta m^{2}_{32}}\textbar) constrains substrate parameters. The
framework structurally accommodates a hierarchy but does not yet predict
the specific ratios; making this prediction explicit is a candidate
next-step calculation.

\emph{Black-hole interior structure.} If black holes host recursive
sublevels (\S\,5.8), their evaporation spectrum and Hawking temperature
would deviate from standard predictions in level-dependent ways.

\textbf{Vacuum energy, dark matter, and dark energy.} The framework
produces a structural account of cosmological-scale vacuum phenomena via
four nested results, presented here at honest epistemic statuses.

\emph{Result 1 (formally derived). Within-level \ensuremath{S_{3}} cancellation.} The 9
bilinear role-correlation operators on a triadic confirmation event
decompose under the \ensuremath{S_{3}} symmetry of agent permutations into 2
trivial-representation modes, 1 sign-representation mode, and 6
standard-representation modes. Only \ensuremath{S_{3}-\text{trivial}} modes can carry nonzero
vacuum expectation value in an \ensuremath{S_{3}-\text{invariant}} substrate vacuum. Of the 9
bilinears, 7 therefore have vanishing vacuum stress-energy contribution
by \ensuremath{S_{3}} symmetry alone. The 2 surviving modes are the U(1) trace mode
\ensuremath{T^{(k)}} (``bit-channel residual'') and the symmetric-singlet mode
\ensuremath{R^{(k)}} (``relation-channel residual''). This 7-of-9 cancellation is
forced by substrate primitives and does not require additional
structural assumptions.

\emph{Result 2 (formally derived, as corollary of \S\,6.3.0). Recursive
absorption.} Under the Level-Advance Relation (\S\,6.3.0), level-k
confirmed bits become level-(k+1) agents, and level-k confirmation
relations become level-(k+1) substrate manifold structure. Applied to
the two surviving vacuum modes: \ensuremath{T^{(k)}} is reabsorbed as matter content
at level k+1 (since the bits carrying it become matter-carrier agents),
and \ensuremath{R^{(k)}} is reabsorbed as substrate structure at level k+1 (since
the relations carrying it become the carrier lattice). The structural
consequence: within-level vacuum residuals do not propagate to higher
levels as vacuum. The ``sub-Planckian vacuum integration'' form of the
cosmological constant problem does not arise in TMS, because sub-level
vacuum modes have been reabsorbed into non-vacuum structure rather than
summed.

\emph{Result 3 (candidate, assumption-dependent). \ensuremath{\Lambda} as steady-state
equilibrium.} If the substrate maintains a dynamical equilibrium between
level-k bit-generation and bit-absorption into level-(k+1) confirmation
events, the unabsorbed-bit population reaches a steady state. The
per-level density-suppression ratio is
\ensuremath{\rho_{p}}\textsuperscript{\ensuremath{(k+1)/\rho_{p}}}(k) = \ensuremath{(M-\text{ratio})/(L-\text{ratio})^{3}} =
\ensuremath{(3\sqrt{2}/2)/(3\sqrt{2})^{3}} = 1/36. Observed \ensuremath{\Lambda /\rho_{p}} \ensuremath{\approx} \ensuremath{10^{-123}} under this scaling gives
a candidate level identity k \ensuremath{\approx} 79 for our universe. Three ingredients
combine here and should not be conflated: the per-level suppression ratio
1/36 is derived from the mass- and length-ratios; the steady-state of the
unabsorbed-bit population is an assumption from rate-balancing, not a
theorem; and the level identity k \ensuremath{\approx} 79 follows only when both are combined
with the observed \ensuremath{\Lambda /\rho_{p}}. The derived content is the 1/36 ratio alone.
k \ensuremath{\approx} 79 is a candidate empirical bridge contingent on the steady-state
assumption, and is falsified independently of 1/36 if that assumption
fails. The steady-state
assumption is asserted from structural rate-balancing, not derived from
a formal substrate-level theorem. A rigorous derivation would route
through \S\,5.3 coupling-timescale machinery and \S\,5.1's saturation cascade
dynamics.

\emph{Result 4 (candidate, structurally motivated). Bit and relation
channels as dark sector candidates.} The two surviving residual modes
have distinct distribution properties. \ensuremath{T^{(k)}} concentrates where
confirmation events have occurred --- i.e., where matter is --- and
behaves matter-like under gravitational coupling. \ensuremath{R^{(k)}} is
distributed across the substrate manifold structure itself, independent
of event concentrations, and behaves field-like. These distribution
properties match the observed phenomenologies of dark matter (clusters
with matter, gravitates) and dark energy (smooth, drives expansion)
respectively. Predicted ratio from substrate combinatorics:
relation-channel : bit-channel = 3 : 1, matching the foundational
substrate ratio \ensuremath{\bar{B}} = \ensuremath{3\bar{\alpha}.} Observed ratio: dark energy : dark matter \ensuremath{\approx} 2.65
: 1. The leading-order prediction of 3:1 differs from observation by
\textasciitilde12\%; higher-order corrections in mode-weighting could
move the prediction in either direction and have not been calculated.
The structural identification is clean; the residual numerical gap is
flagged as open.

Underlying structural identification (proposed but not yet ratified):
substrate agents correspond to the energy-side of physical content, and
substrate bits correspond to the mass-side. Under this identification,
agents : bits = 3 : 1 from \ensuremath{\bar{B}} = \ensuremath{3\bar{\alpha}} is the substrate-side origin of the
energy-mass split in observable physics. This identification has not
been formally introduced in Papers 1-4 and requires consistency audit
before ratification.

\subsubsection*{7.6.5 The Structural Reduction}

Standard physics has four independent fundamental quantities in the
gravitational sector: \ensuremath{\hbar ,} c, G, and the mass scale. TMS reduces these to
one geometric primitive (FCC face-diagonal-to-edge ratio = \ensuremath{\sqrt{2}),} one
axiom (Principle of Generative Economy, Axiom 0), and one empirical
question (level identity k). The Einstein field equations are the
continuum limit of a substrate recursion whose coupling is structurally
forced and whose absolute scale is a property of our universe's location
in the hierarchy.

\subsection*{7.7 Connection to Independent Programs}

The TMS account of gravitational phenomena connects with several
independent programs that derive gravity from underlying substrate
dynamics. Verlinde's entropic-gravity proposal (Verlinde, 2011)
identifies gravity as an entropic force arising from holographic
screens; in the TMS reading, the holographic screens correspond to the
boundaries between hierarchical levels, with the entropic content being
the algorithmic complexity of the level-coupling. Bianconi's
gravity-from-entropy framework (Bianconi, 2025) derives gravity from a
relative-entropy action between the spacetime metric and the
matter-induced metric; in the TMS reading, the two metrics correspond to
the level-k tensor sector and the level-(k-1) tensor sector that meet at
the embedding boundary. Padmanabhan (2014) derives the cosmological
constant from holographic mode-counting under emergent-gravity
assumptions; the TMS framework parallels this structurally through its
level-counting of substrate modes.

Two cellular-automaton programs warrant explicit contrast. The
QCA-to-QFT programs of D'Ariano and Perinotti (2014) and Mlodinow and
Brun (2025) develop continuum-limit field theory from discrete substrate
dynamics; the TMS framework derives the equivalent lift
substrate-internally (Section 4.1, Lemmas A and B) using only TMS
primitives, with these external programs serving as comparison
references rather than as load-bearing imports. The QICT (Quantum
Information Copy Time) framework (Quantum Rep., 2026) derives emergent
gravity from gauge-coded QCAs; structurally close to TMS in spirit, with
the level-identification problem in TMS finding a natural analog in
QICT's copy-time geometry.

A direct contrast with empirical implications: Wetterich's
cellular-automaton spinor-gravity construction (Wetterich, 2022)
predicts exact local Lorentz symmetry on the discrete level for a
non-FCC lattice designed specifically to preserve it. TMS on the FCC
lattice predicts residual Lorentz violation suppressed by
\ensuremath{(E/E_{\mathrm{Planck}})^{2}} in the inversion-symmetric (scalar and gauge-vector)
sectors, with a stronger \ensuremath{(E/E_{\mathrm{Planck}})^{1}} signal only in the
inversion-broken matter sector, per Section 4.1, Lemma B. These are
competing falsifiable predictions; future high-energy cosmic-ray
observations will distinguish the two substrate hypotheses.

None of these independent programs starts from the same geometric
primitives as TMS --- circles, triadic confirmation, and the FCC
lattice. The convergence at the level of structural results --- gravity
as an emergent, geometric, holographic, information-theoretic phenomenon
--- supports the broader picture that classical gravitational dynamics
is not a fundamental force but a structural consequence of the
underlying substrate's organization across scales.

\section*{VIII. Structural Correspondence to Standard Model Gauge Structure}
\addcontentsline{toc}{section}{VIII. Structural Correspondence to Standard Model Gauge Structure}

The Standard Model gauge group SU(3) \ensuremath{\times} SU(2) \ensuremath{\times} U(1), its representation
content for matter fields, and the fermion-boson decomposition are among
the most striking empirical inputs to modern physics. We address their
structural origin here at the level of structural correspondence. The
TMS does not derive the Standard Model in full; it provides the
substrate-level carrier structure for an SM-form gauge theory, with the
explicit representation-theoretic mapping reserved for subsequent work.

\subsection*{8.1 The Substrate-Level Stabilizer Group}

The [[4,2,2]] stabilizer code at the substrate level has
stabilizer group \ensuremath{Z_{2}} \ensuremath{\times} \ensuremath{Z_{2},} generated by \ensuremath{S_{x}} = \ensuremath{X_{1}X_{2}X_{3}X_{4}} and \ensuremath{S_{z}} = \ensuremath{Z_{1}Z_{2}Z_{3}Z_{4}}
(Paper 1 \S\,V). By the Recursion-Preservation Theorem (\S\,2.5), each
hierarchical level reproduces the [[4,2,2]] structure at its own
scale, giving a \ensuremath{Z_{2}} \ensuremath{\times} \ensuremath{Z_{2}} stabilizer group at every level k \ensuremath{\geq} 0.

The continuous gauge group is not obtained from any general theorem that
finite Abelian stabilizers must extend to continuous Lie groups --- there is no
such theorem, and many discrete gauge structures remain discrete. It is obtained
by explicit construction: the push-forward transport of \S\,8.3 builds the
generators from substrate primitives and closes them, yielding su(3) with the
trace mode carrying a separate U(1) (the commutator computation is given in
Appendix H). The role of the \ensuremath{Z_{2}} \ensuremath{\times} \ensuremath{Z_{2}} stabilizer hierarchy here is to
supply, through its cross-level structure, the non-commutativity that the
transport requires; the continuous group itself is the transport's closure, not
an assumed Abelian lift.

\subsection*{8.2 Cross-Level Commutator Mechanism (Bifurcation-Derived)}

The non-Abelian closure that converts U(1) \ensuremath{\times} U(1) into SU(2) \ensuremath{\times} U(1) and
ultimately SU(3) \ensuremath{\times} SU(2) \ensuremath{\times} U(1) is supplied, in the TMS picture, by the
sublattice bifurcation of \S\,6.1 and \S\,5.3. At each level k \ensuremath{\geq} 1, the
substrate splits into two FCC sublattices A and B related by spatial
inversion. The stabilizer dynamics on A and B are independent at level k
but couple at level k + 1 through the shared centroids that host the
level-(k+1) confirmed bits.

The cross-level coupling between A-sublattice and B-sublattice
stabilizers at level k provides the commutator structure required for
non-Abelian gauge closure at level k + 1. Specifically, the A-stabilizer
\ensuremath{S_{x}^{A}} and the B-stabilizer \ensuremath{S_{x}^{B}} at level k do not generally commute
when their actions are evaluated at the same level-(k+1) centroid,
because the A and B dynamics run independently at level k and their
joint readout at the shared centroid is sensitive to the order in which
the readout combines them. This non-commutativity is the substrate-level
origin of the non-Abelian closure of the gauge group at the meta-level.

This is the structural mechanism that the original gauge appendix (now
relabeled Appendix F) hand-waved. The bifurcation supplies the
cross-level commutator naturally, without requiring any additional
structure to be posited.

\subsection*{8.3 Hierarchical Carrier Structure for SU(3) \ensuremath{\times} SU(2) \ensuremath{\times} U(1)}

The Standard Model gauge group emerges from the substrate primitives at
distinct levels of the hierarchy. This section gives the revised
structural derivation; earlier drafts placed SU(3) at level 2 via
sublattice bifurcation acting on level-1 doubling, but that account gave
2 \ensuremath{\times} 2 = 4 sublattices rather than 3 colors. The corrected derivation
places SU(3) at the level-0 substrate primitive itself.

\textbf{SU(3) at level 0.} Every confirmation event involves three
agents closing around one bit \ensuremath{(\bar{B}} = \ensuremath{3\bar{\alpha},} Paper 1 \S\,III). The event is
invariant under \ensuremath{S_{3}} permutations of agent labels. In the QCA-to-QFT
continuum lift (\S\,4.1), this discrete \ensuremath{S_{3}} symmetry on label space lifts to
continuous local gauge symmetry on the 3-dimensional internal space of
agent labels. \ensuremath{S_{3}} is the Weyl group of SU(3); standard gauging of ``3
indistinguishable label slots'' produces SU(3) as the continuum gauge
group. The 8 gluons of QCD correspond to the 8 traceless bilinear
role-correlation operators on the 3-agent space (the 9th, trace, mode
being U(1)). Color confinement is the substrate-side requirement that
only triadic-closed configurations are observable substrate events:
single agents (single quarks) do not appear in isolation because triadic
closure is structurally required for actualization.

\textbf{SU(3), rigorously: the push-forward derivation and its global
form.} The label-permutation account above identifies the
three-dimensional internal space but does not by itself fix the gauge
group, and a literal reading of the confirmation readout does not yield
SU(3): the confirmation rule P = (1 + \ensuremath{\Pi )/2} reads a real bias and so
preserves only real-orthogonal symmetry, giving so(3) at the readout.
The gauge group does not live in the readout. It lives in the connection
--- the transport that pushes a confirmed configuration forward to the
next meta-flop. Because the substrate wave equation is second order
(Paper 1 \S\,6.3; Paper 2 \S\,VII), the transported object is not the value \ensuremath{\psi}
alone but the complex amplitude z = \ensuremath{\psi} + \ensuremath{i\Pi} carrying the momentum \ensuremath{\Pi} (the
rate of change \ensuremath{\dot{\psi});} the factor i is required by propagation, not
imported. This is the cross-level \ensuremath{\text{temporal}\to \text{spatial}} flip of \S\,2.3 acting
as the carry.

On the three-agent label space the transport decomposes into two
substrate pieces. Value transport --- the real rotations of the agent
labels --- generates the real antisymmetric matrices (so(3), three
generators). The momentum carry, to which the forced i attaches
symmetric reorientations, generates i times the real symmetric traceless
matrices (five generators). These eight close under commutation to
exactly the eight-dimensional algebra of anti-Hermitian traceless \ensuremath{3\times 3}
matrices, with no overshoot --- neither to \ensuremath{sl(3,\mathbb{C})} (dimension 16) nor to
sp(6) (dimension 21). The generators are anti-Hermitian and traceless:
anti-Hermiticity is the norm-preservation of the transport (the
stochastic resolution term is zero-mean, \ensuremath{E[\chi}\textbar \ensuremath{\Pi ]} = \ensuremath{\Pi} exactly,
and the mean-field re-confirmation spectral radius never exceeds 1;
verified in Appendix G), and tracelessness is the absence of the U(1)
trace mode.

These three properties fix the continuum algebra uniquely.
Anti-Hermiticity selects the compact real form; a non-degenerate,
negative-definite Killing form with an irreducible adjoint action
(commutant of dimension one) makes the algebra simple and compact; and
the rank is two with six non-zero roots --- the \ensuremath{A_{2}} root system. Because
the unique compact simple Lie algebra of dimension eight is su(3), the
continuum lift of the transport symmetry is forced to be su(3): the
eight gluons are the eight transport generators, with the trace mode
removed by tracelessness carrying the separate U(1). The check is non-circular, and the absence of overshoot is not an
empirical observation about the closure but a structural fact: the eight
generators are anti-Hermitian and traceless by construction, and the
commutator of any two anti-Hermitian traceless matrices is again
anti-Hermitian (\ensuremath{[X,Y]^{\dagger} = -[X,Y]}) and traceless
(\ensuremath{\operatorname{tr}[X,Y] = 0} identically). The space of anti-Hermitian
traceless \ensuremath{3\times 3} matrices is exactly eight-dimensional and is su(3); the
bracket cannot leave it, so there is nowhere to overshoot. A control algebra
of the same dimension, compactness, and anti-Hermitian character
(\ensuremath{su(2) \oplus su(2) \oplus u(1)^{2}}) fails the simplicity test, distinguishing
su(3) from same-dimension imposters. The explicit structure constants,
the negative-definite Killing form, the trivial center, and the rank-two
\ensuremath{A_{2}} root system are computed in Appendix H.

The global form is fixed by the matter content. The three agents carry
the fundamental representation 3, which is irreducible (the only
operators commuting with all eight generators are the scalars). The
centre of the group is therefore the scalars \ensuremath{\zeta \mathbf{1}} with \ensuremath{\zeta^{3}} = 1, namely \{1,
\ensuremath{\omega ,} \ensuremath{\omega^{2}}\} \ensuremath{\cong} \ensuremath{\mathbb{Z}_{3},} acting with triality \ensuremath{\omega} on a single agent, \ensuremath{\omega^{2}} on an agent
pair (the \ensuremath{\bar{3}),} and \ensuremath{\omega^{3}} = 1 on a triadically closed event. Single agents
are substrate-real primitives on which the centre acts non-trivially, so
faithful action requires the cover: the global gauge group is SU(3), not
\ensuremath{SU(3)/\mathbb{Z}_{3}.} Confinement is then exactly the statement that only \ensuremath{\mathbb{Z}_{3}-\text{neutral}}
configurations --- the triadically closed events --- actualize, the
substrate-level form of the unbroken-centre confinement criterion of
QCD. The same forced i that lifts so(3) to su(3) is what makes the
centre \ensuremath{\mathbb{Z}_{3}:} the value-only transport generates the real group SO(3),
whose centre is trivial.

\textbf{Status.} Given the QCA-to-QFT continuum lift of \S\,4.1 (itself a
structural correspondence), the identification of the colour gauge group
is \emph{Derived}: the algebra is su(3) (no overshoot, simple, compact),
the global group is SU(3) (fundamental triality of the substrate-real
agents), and confinement is the \ensuremath{\text{unbroken}-\mathbb{Z}_{3}-\text{centre}} statement. What
remains a structural correspondence is the lift itself --- that the
discrete transport symmetry promotes to a gauged continuous symmetry ---
not the group identity, which the lift now determines. The earlier
label-permutation reading \ensuremath{(S_{3}} as the Weyl group of SU(3)) is recovered
as the readout-side shadow of this transport symmetry and is consistent
with it: so(3) at the readout, su(3) in the connection.

\textbf{SU(2) at level 1.} The level-1 substrate carries two
interpenetrating FCC sublattices (the A and B sublattices of \S\,6.1.1).
Under inversion through the level-2 centroid of each stella-octangula
tile, the A and B sublattices exchange. The polarization field \ensuremath{\psi} at
level 1 decomposes as a two-component doublet \ensuremath{\psi} = \ensuremath{(\psi_{A},} \ensuremath{\psi_{B}).} Gauging
this doublet structure produces \ensuremath{SU(2)_{L}} electroweak. (This SU(2) is an
\emph{internal gauge} symmetry on the sublattice-membership label; the
fermion spinor phase is a separate structure, deriving from the triadic
oddness of \S\,4.1. Each carries its own role: gauge SU(2) here, spinor from
the threeness there.) The W and Z bosons
correspond to the 3 traceless bilinear operators on the 2-component
doublet (the 4th, trace, mode contributing to U(1)).

The group identity here is fixed by the same push-forward argument that
fixes the colour group, now applied to the two-component doublet rather
than the three-agent space. On the A/B doublet the transport again
splits into value transport --- the real rotation of the \{A, B\}
labels, generating the single real antisymmetric \ensuremath{2\times 2} matrix (so(2)) ---
and the momentum carry, the forced i attaching to the two real symmetric
traceless \ensuremath{2\times 2} reorientations. These three generators close to exactly
the anti-Hermitian traceless \ensuremath{2\times 2} algebra, su(2), with no overshoot to
\ensuremath{sl(2,\mathbb{C})} (real dimension 6); a control adding the forbidden directions
reaches dimension 6, so the closure test again detects overshoot. The
algebra is compact (anti-Hermitian), simple (non-degenerate Killing
form, irreducible adjoint), and rank one with two roots --- the \ensuremath{A_{1}}
system --- so the continuum lift is su(2) uniquely. The
norm-preservation of Appendix G applies verbatim, the doublet
propagating under the same second-order wave operator.

The global form is again fixed by the matter content. The doublet \ensuremath{(\psi_{A},}
\ensuremath{\psi_{B})} is the irreducible fundamental 2; the centre of the group is the
scalars \ensuremath{\zeta \mathbf{1}} with \ensuremath{\zeta^{2}} = 1, namely \{\ensuremath{+\mathbf{1},} \ensuremath{-\mathbf{1}}\} \ensuremath{\cong} \ensuremath{\mathbb{Z}_{2}.} The non-trivial element
\ensuremath{-\mathbf{1}} acts as \ensuremath{-1} on a single sublattice component (the 2) but trivially on
the W-boson adjoint (the 3). A single sublattice site is a
substrate-real primitive on which the centre acts non-trivially, so
faithful action requires the cover: the level-1 gauge group is SU(2),
not SO(3) = \ensuremath{SU(2)/\mathbb{Z}_{2}.} As with colour, the value-only transport would
generate the real rotation group SO(2), whose centre is trivial; it is
the momentum-carried i that produces the \ensuremath{\mathbb{Z}_{2}} centre.

\textbf{Status and the two boundaries.} Under the \S\,4.1 lift, the level-1
gauge factor is \emph{Derived} as su(2) with global form SU(2), by the
same closure and matter-content arguments as colour. Two boundaries are
stated honestly. First, \emph{chirality} --- that the gauge group acts
on left-handed components only, the \ensuremath{SU(2)_{L}} of the Standard Model ---
is not part of the gauge closure. The closure fixes the group; the
left-handedness is supplied by the inversion-odd parent-record break
\ensuremath{d_{z}} (\S\,5.3, \S\,8.4.4), the same asymmetry that carves the fermion sectors,
and is treated there rather than here. Second, the precise electroweak
\emph{global quotient} --- whether the combined electroweak group is
SU(2) \ensuremath{\times} U(1) or the quotient (SU(2) \ensuremath{\times} \ensuremath{U(1))/\mathbb{Z}_{2}} \ensuremath{\cong} U(2), which in the
Standard Model is fixed by the hypercharge assignments of the matter
multiplets --- depends on the matter representation content and remains
a structural correspondence co-located with the fermion-roster work.
What is Derived is the SU(2) factor itself and its U(1) partner as the
trace mode; the way they are glued is the reserved piece.

\textbf{U(1) at level 0.} The overall phase of the bilinear
role-correlation operators (the trace mode on the 3-agent space) carries
U(1). This is the single \ensuremath{S_{3}-\text{trivial}} component of the bilinear
decomposition that is not part of the SU(3) generators.

This is a structural derivation rather than an external import. The
mapping from discrete substrate symmetries to continuous Lie groups
follows the substrate-internal QCA-to-QFT lift of \S\,4.1 (Lemmas A and B),
which routes through Paper 1 \S\,III, Paper 1 \S\,VI, Paper 2 \S\,IV, and Paper 3
\S\,IV without inheriting external machinery. The connection to independent
QCA-to-QFT programs (Mlodinow \& Brun, 2025; D'Ariano \& Perinotti,
2014) is now a comparison citation, not a load-bearing import --- see
\S\,7.7.

\subsection*{8.4 Matter Content: Fermionic-Bosonic Sectors and Generations}

Fermionic content originates in the half-integer spin forced by the
triadic oddness of \ensuremath{\bar{B}} = \ensuremath{3\bar{\alpha}} (\S\,4.1, Lemma A): a single confirmation
node recouples three spin-\ensuremath{\tfrac{1}{2}} agents to a half-integer and is a fermion,
while a bound pair recouples to an integer and is a boson (Corollary A.1).
The A/B sublattice doublet plays a separate role, supplying the internal
gauge structure --- the \ensuremath{SU(2)_{L}} label space of \S\,8.3 on which the
electroweak break \ensuremath{d_{z}} acts. Bosonic content emerges as the coarse-grained
polarization decomposition of \S\,4.2: scalar (Higgs-like), polar vector
(gauge), symmetric-tensor geometry, and the odd-parity tensor partner.

\subsubsection*{8.4.1 The Three-Generation Count from 2D Minimum Closure}

The three generations of Standard Model fermions correspond to the three
distinguishable agent-slots in the primitive triadic confirmation event.
The argument proceeds in three steps.

\emph{Step 1: Confirmation events are intrinsically two-dimensional.}
The primitive substrate operation is \ensuremath{\bar{B}} = \ensuremath{3\bar{\alpha}} (Paper 1 \S\,III): three agents
close around one bit in a triangular configuration on the manifold \ensuremath{M_{\alpha}.}
The triangle is a 2D object. The minimum closure at 2D requires exactly
3 participating agents; the minimum closure at 3D would be tetrahedral
\ensuremath{(\bar{B}} = \ensuremath{4\bar{\alpha}^{3},} requiring 4 participating agents).

\emph{Step 2: The substrate hierarchy embeds 2D events in 3D packing,
not the other way.} The FCC lattice (Paper 1 \S\,VI) is the 3D packing
structure that supports level-1 confirmation events. The 3D structure
governs spatial packing of agent-spheres but does not change the
dimensionality of the confirmation event itself. Each level-1 site
carries triadic structure inherited from the 2D primitive beneath.

\emph{Step 3: Matter-carrier count inherits the 2D minimum, not the 3D
packing minimum.} Under recursive embedding through the hierarchy
(\S\,2.5), distinguishable agent-slots are preserved by the
recursion-preservation structure. Each agent-slot in the primitive 2D
event becomes a distinguishable matter-carrier under the recursive lift.
Three agent-slots produce three matter-carrier generations.

The framework predicts three generations, not four, because confirmation
events live in 2D, not 3D. The 3D structure governs spatial packing
rather than the count of distinguishable confirmation roles.

The nature of this result should be stated precisely, because it is easy to
overread. The ``three'' of the generation count is the \emph{same} three as the
\ensuremath{3} of \ensuremath{\bar{B}} = \ensuremath{3\bar{\alpha}} --- the founding valence, re-expressed through the
recursion as a count of distinguishable matter-carrier slots. This is not an
independent parameter that happened to come out equal to three; it is the
founding number reappearing in a new role. Whether that recurrence is a deep
structural echo (the same threeness manifesting at two levels of description) or a
tautology (three put in, three read out) is not something the framework can
adjudicate from the inside, and it is not claimed to be an independent
prediction of the number three. What \emph{is} substantive is the mechanism: that
the generation count is tied to the confirmation valence at all --- that ``why
three generations'' and ``why triadic confirmation'' are the same question rather
than two --- and that the count is forced to follow the 2D confirmation minimum
rather than the 3D packing minimum. The result is therefore labeled a structural
correspondence tying generation number to confirmation valence, not a
first-principles derivation of the numeral three from independent structure.

The informal invocation of \S\,2.5 in Step 3 above can be promoted to a
formal theorem, completing the derivation of the count ``three'' within
the framework.

\begin{theorem}[Recursion-Preservation for Matter Content]
Under \ensuremath{\hat{R},} the structural form of matter content --- the count of
generation slots forced by \ensuremath{\bar{B}} = \ensuremath{3\bar{\alpha},} the \ensuremath{O_{h}} irrep decomposition of the
polarization field on the FCC nearest-neighbor space, and the
inversion-symmetry character that distinguishes fermion from boson
sectors --- is preserved at every recursion level. The matter content at
level k+1 has the same structural form as at level k.
\end{theorem}

\begin{proof}
Three components.

\emph{(i) \ensuremath{\bar{B}} = \ensuremath{3\bar{\alpha}} at every level.} By the recursion-preservation theorem
of \S\,2.5, \ensuremath{\Xi (R)} at level k+1 reproduces the level-k substrate's triadic
confirmation rule under the same five axioms. The triadic minimum \ensuremath{\bar{B}} = \ensuremath{3\bar{\alpha}}
in 2D, derived from the angular-closure argument of Paper 1 \S\,II.1,
applies to the substrate primitive at any scale. Therefore the count
three of the triadic minimum is preserved at every level.

\emph{(ii) \ensuremath{O_{h}} irrep decomposition at every level.} Section \S\,4.2 gives
the complete five-irrep decomposition of the polarization field under
\ensuremath{O_{h}} on the FCC nearest-neighbor space, exhausting all twelve
dimensions. The FCC connectivity and \ensuremath{O_{h}} symmetry are preserved by \S\,2.5
at every level. Therefore the irrep decomposition is preserved at every
level.

\emph{(iii) Fermion-versus-boson character at every level.} The
substrate-level distinction between fermion sectors and boson sectors in
the polarization field is determined by two structures: the \ensuremath{O_{h}} irrep
decomposition (which classifies modes by their transformation under the
substrate's full point group, \S\,4.2), and the action of
inversion-breaking from the parent-record on those modes (which carves
the multifold sectors into chiral fermion content, \S\,5.3, while leaving
inversion-symmetric sectors as bosons, \S\,9.6). The irrep decomposition is
preserved at every level by component (ii) above. The parent-record
break is structurally guaranteed at every \ensuremath{\hat{R}-\text{produced}} level by the
recursion structure itself: by the definition of reinterpretation
(\S\,2.3), every level has a parent whose biographical record provides the
inversion-breaking environment for the child substrate. The existence of
a break is therefore not a feature specific to our level but an
automatic structural consequence of being an \ensuremath{\hat{R}-\text{produced}} substrate. Both
the classifier (the \ensuremath{O_{h}} irrep decomposition) and the actor (the
parent-record break) are preserved at every recursion level; therefore
the fermion-versus-boson distinction is preserved at every level.

The theorem preserves \emph{structural form}, not specific
instantiation: the matter content at level k+1 has the same count of
generation slots, the same irrep decomposition, and the same
fermion-boson character as at level k. The specific configurations
realizing matter content at each level differ --- this is what makes
hierarchical recursion non-trivial --- but the structural form they
realize is identical.

\end{proof}

\begin{corollary}[Structural Count Preservation]
The structural
count ``three matter-carrier slots'' is preserved at every recursion
level. The observed three fermion generations in our substrate result
from the combination of (a) this structural preservation at every level
and (b) the \ensuremath{\text{depth}-\geq -3} arrow-tower realization that selects three
frozen-arrow carriers from the accumulated recursion (\S\,8.4.2). Together,
(a) and (b) upgrade the count ``three'' from substrate-forced to
formally derived.
\end{corollary}

\begin{remark}
The theorem makes the inheritance explicit but does not
yet identify \emph{which} three generations are realized --- that
question concerns specific configurations and mass values, both of which
are governed by the biographical record \ensuremath{\mathcal{B}(R)} and remain the subject of
\S\,8.4.3.
\end{remark}

\textbf{Status.} With this theorem the three-generation count of \S\,8.4.1
is promoted from substrate-forced to Derived; the identification of the
three structural carriers with the three observed fermion generations
remains a structural conjecture (\S\,8.4.2), and the mass values remain
open (\S\,8.4.3).

\subsubsection*{8.4.2 The Same 3 Thrice}

The ``3 of three colors'' (\S\,8.3, SU(3) gauging) and the ``3 of three
generations'' (\S\,8.4.1) are the same 3, distinguished by which
symmetry-axis is projected. Color is \emph{which-agent-now}: the \ensuremath{S_{3}}
action within a single confirmation event. Generation is
\emph{which-agent-slot-as-matter-carrier}: the \ensuremath{S_{3}} action across the
recursive embedding that produces observable matter content. Two
structurally distinct symmetry-axes acting on the same underlying
triadic primitive.

Independent corroboration: Jegerlehner (2019) argues that three fermion
generations are required to make CP violation and baryogenesis emerge
naturally in low-energy effective Standard Model scenarios. This is a
different reason for the count 3, derived from low-energy phenomenology
rather than substrate combinatorics, arriving at the same number.

A third independent route comes from the arrow tower under recursive
reinterpretation. Each reinterpretation event freezes its parent's
causal-succession axis (\S\,2.3, applied at the meta level via the
recursion-preservation theorem of \S\,2.5) into the child's spatial
geometry as a frozen directional structure. Successive reinterpretation
events accumulate a tower of frozen arrows in the descendant substrate,
one per ancestor level. The Principle of Generative Economy compresses
linearly-dependent directions as not-minimum-sufficient, so the tower
saturates at the linear-independence ceiling of the substrate's spatial
dimension. With substrate dimension 3 (Paper 1 \S\,II.4), \textbf{the tower
saturates at exactly three frozen arrows for any recursion depth \ensuremath{\geq} 3.}
If each fermion generation corresponds to dominant alignment with one
frozen arrow, the count of generations equals the spatial dimension,
locked.

Three structurally distinct routes --- substrate combinatorics,
low-energy phenomenology, and recursion-tower saturation --- converging
on the same integer indicate that the count is overdetermined rather
than coincidental.

Two side-predictions fall out of the arrow-tower route. First, a
recursion-depth lower bound: the observation of three generations
implies the cosmological recursion has depth \ensuremath{\geq} 3 ancestor substrates,
sharpening \S\,5.8's cosmological natural selection from structural
prediction to a count with empirical traction. Second, a
dimensionality--generations lock: the generation count is not a free
number; it is identical to dim(space). A substrate framework deriving a
different spatial dimension would predict a correspondingly different
generation count.

\textbf{Status of Results for \S\,8.4.2:}

\begin{itemize}
\item
  \emph{Derived:} Tower saturation at min(depth, dim) under PGE
  compression; for our substrate (dim = 3, depth \ensuremath{\geq} 3), the count is 3.
\end{itemize}

\begin{itemize}
\item
  \emph{Structural correspondence:} The convergence of three independent
  geometric routes on the count 3.
\end{itemize}

\begin{itemize}
\item
  \emph{Prediction / Conjecture:} The mapping from frozen arrows to
  specific fermion generations via dominant alignment; the mass values
  themselves reside in \ensuremath{\mathcal{B}(R)} and are not derivable from the tower alone
  (cf. \S\,8.4.3 and keystone idealization 2).
\end{itemize}

\subsubsection*{8.4.3 Mass Spectrum}

The mass spectrum's dimensionful values are set by the band-crossing
dispersion together with the biographical record, and are calibrated
against experiment in the manner of the Standard Model's Yukawa
couplings (developed below and in \S\,8.4.3's status block).

Mass spectrum values follow from (i) the detailed band-crossing structure
on the electroweak A/B gauge doublet under the FCC dynamics (\S\,8.3),
(ii) the substrate-side coupling structure J(k) = \ensuremath{J_{0}} \ensuremath{\cdot}
\textbar F(k)\textbar\ensuremath{^{2}} where F(k) is the A-tetrahedron structure factor
(\S\,6.1.2), and (iii) the interaction of band-crossing dispersion with the
cross-level coupling structure of \S\,7.5.

This calculation is developed as future work. The structural claim is
that mass differences between generations correspond to differences in
how each generation's agent-slot is anchored in the recursive embedding,
with the mass spectrum following from band-structure analysis under the
substrate dynamics.

A structural decomposition of the spectrum is nonetheless available,
separating what the framework forces from what the biography supplies.
The fermion rest frequencies \ensuremath{\omega_{0}} are the protected standing-mode
frequencies of the re-confirmation dynamics: a standing program
re-confirms in place with rest frequency \ensuremath{\omega_{0}} = \ensuremath{\Omega (0),} giving mass m =
\ensuremath{\hbar \omega_{0}/c^{2}} (\S\,7.4). The inversion-odd structured break \ensuremath{d_{z}} (\S\,8.4.5, \S\,5.3)
acts on the vector sector of the \ensuremath{O_{h}} field decomposition (\S\,4.2), and
each frozen arrow of the recursion tower (\S\,8.4.2) contributes an
inversion-odd break aligned with its axis. The break-induced
mass-squared operator on the three-dimensional vector sector is
therefore \ensuremath{M^{2}} = \ensuremath{\Sigma_{g}} \ensuremath{w_{g}} \ensuremath{\hat{n}_{g}} \ensuremath{\hat{n}_{\mathrm{gT}},} summed over the tower's frozen
arrows \ensuremath{\hat{n}_{g}} with biographical amplitudes \ensuremath{w_{g}} \ensuremath{\geq} 0; the rest frequencies
are the square roots of its eigenvalues.

Two consequences are structural. First, the count: because the
PGE-compressed tower saturates at exactly dim(space) = 3 linearly
independent arrows (\S\,8.4.2), \ensuremath{M^{2}} is a 3 \ensuremath{\times} 3 matrix and the spectrum has
exactly three rest frequencies --- locked to the spatial dimension, so a
substrate of different dimension would predict a correspondingly
different count. Second, the existence of a hierarchy: for any linearly
independent arrow triad the three eigenvalues of \ensuremath{M^{2}} are generically
distinct, so a non-degenerate mass splitting is forced rather than
assumed --- degeneracy occurs on a measure-zero set that includes both
the orthonormal equal-amplitude case and symmetric triads --- such as a
tetrahedron-vertex triad along the \ensuremath{\langle 110\rangle} directions --- at equal
amplitudes; a non-degenerate three-level hierarchy therefore requires
amplitude asymmetry or triad asymmetry, both supplied by the biography
\ensuremath{\mathcal{B}(R)}. Direct evaluation on a generic saturated triad gives
rest-frequency ratios of order a few even at equal amplitude, widening as
the biographical amplitudes spread.

One premise underlies both consequences and deserves explicit statement:
they assume the tower's frozen arrows are linearly independent. The
arrow-tower argument (\S\,8.4.2) establishes independence for generic
arrows under PGE compression, but whether the arrows a given biography
actually freezes are independent --- rather than collapsing onto a
single repeated succession axis --- is a property of the
reinterpretation operator \ensuremath{\hat{R}} that has not yet been demonstrated from the
substrate dynamics. The count and the hierarchy are accordingly
conditional on arrow independence, which a faithful construction of \ensuremath{\hat{R}}
must establish.

\emph{Localized matter is massive; matter has a minimum mass.} Two
structural facts sharpen the mass picture below the level of specific
values. First, a localized triadic confirmation must re-confirm in place
to persist --- an isolated confirmation that does not re-confirm decays
back into the \ensuremath{U_{\mathrm{sh}}} background --- so any standing, localized program has a
non-zero rest re-confirmation frequency \ensuremath{\omega_0 > 0} and hence, through
\ensuremath{m = \hbar\omega_0/c^2} (\S\,7.4), a non-zero rest mass. Localization forces
mass. Second, the minimum stable standing program under the operational
flop dynamics has size four --- a bent four-cluster with alternating
chirality --- so there is a smallest standing re-confirmation frequency
and matter carries a minimum mass.

\emph{The re-confirmation amplitude spectrum is discrete and rational.}
The self-re-confirmation amplitude \ensuremath{|\Pi|} of a standing program --- the
net agreement of a node with its own prior confirmation, divided by its
coordination --- is a ratio of integers and takes only a small set of
robust values under the flop dynamics: \ensuremath{|\Pi| \in \{1, \tfrac{1}{2}, \tfrac{1}{3}, 0\}}. The
value \ensuremath{|\Pi| = 1} is maximal freeze (the most tightly bound, heaviest
standing mode); \ensuremath{|\Pi| = \tfrac{1}{3}} is the triadic signature of \ensuremath{\bar{B}} = \ensuremath{3\bar{\alpha}}
directly; \ensuremath{|\Pi| = \tfrac{1}{2}} is the paired (bosonic) amplitude; and \ensuremath{|\Pi| = 0}
is the massless limit. These are the harmonic values \ensuremath{1/n}, so the
re-confirmation amplitudes are mode numbers rather than a continuous
parameter. This is the structural content of the mass spectrum: a mode
structure in \ensuremath{\Pi}-amplitude space, discrete and rational, forced by the
counting.

\emph{The map from \ensuremath{|\Pi|} to particle mass values is open.} The discrete
\ensuremath{\Pi}-structure fixes the mode architecture; the dimensionful mass values in
MeV are set by the full band-crossing dispersion together with empirical
calibration, exactly as the Standard Model takes its Yukawa couplings from
experiment. The framework derives the mode structure and leaves the value
map open, at the honest boundary between the two.



\textbf{Status of Results for \S\,8.4.3:}

\emph{Derived:} The three-rest-frequency count and the mass-squared
operator \ensuremath{M^{2}} = \ensuremath{\Sigma_{g}} \ensuremath{w_{g}} \ensuremath{\hat{n}_{g}} \ensuremath{\hat{n}_{\mathrm{gT}}} on the vector sector, with the count
locked to dim(space) = 3 through the arrow tower (\S\,8.4.2); that localized
matter is massive (\ensuremath{\omega_0 > 0} from the persistence-requires-re-confirmation
argument); that matter has a minimum mass (smallest stable standing
program, size four); and that the re-confirmation amplitude spectrum is
discrete and rational, \ensuremath{|\Pi| \in \{1, \tfrac{1}{2}, \tfrac{1}{3}, 0\}}, a mode structure in
\ensuremath{\Pi}-amplitude space.

\emph{Structural correspondence:} The existence of a non-degenerate mass
hierarchy as a generic (non-fine-tuned) property of \ensuremath{M^{2}.}

\emph{Prediction / Conjecture:} The specific rest-frequency ratios and
the overall mass scale, set by the arrow geometry and amplitudes of
\ensuremath{\mathcal{B}(R);} co-located with the biographical simulation (keystone idealization
2).

\emph{Open:} The map from the discrete \ensuremath{\Pi}-structure to particle mass
\emph{values} in MeV, set by the full band-crossing dispersion plus
empirical calibration, as the Standard Model takes its Yukawa couplings
from experiment. The framework derives the mass \emph{architecture} (the
counting, the minimum, the mode structure); the dimensionful values live
at the calibration boundary.

\subsubsection*{8.4.4 Open Items}

\emph{Chirality and parity violation.} The mechanism by which the
parent-record break breaks inversion is now structurally identifiable:
under \S\,2.5's recursion-preservation theorem, a reinterpretation event
has the structure of a meta-confirmation, and its meta-bit inherits the
inversion-odd cyclic chirality that every confirmed bit carries (cf.
\S\,4.1, Lemma A). \emph{The break breaks inversion because it is a bit,
and bits are chiral.} The local mechanism by which this break produces
chirality is now demonstrated. The break \ensuremath{d_{z}} = \ensuremath{\langle \mathcal{B}_{A}} \ensuremath{-} \ensuremath{\mathcal{B}_{B}\rangle} is odd
under the \ensuremath{A\leftrightarrow B} inversion, and the \ensuremath{\hat{R}} boundary relates its two sides by
exactly that inversion, so \ensuremath{d_{z}} changes sign across the boundary without
further assumption. A sign-changing mass across a codimension-1 surface
is the Jackiw--Rebbi / Kaplan domain-wall construction: it binds exactly
one zero-energy mode, localized on the wall, of definite chirality.
Direct computation confirms this on the \ensuremath{d_{z}} profile --- a single
mid-gap mode at zero energy bound to the wall, carrying a definite
chirality eigenvalue, with its handedness tracking \ensuremath{\text{sign}(d_{z}):} reversing
the break reverses the bound mode's group velocity. The substrate's
inversion-odd break is therefore what selects left over right,
converting the SU(2) of \S\,8.3 into the chiral \ensuremath{SU(2)_{L}} of the Standard
Model at the level of the local wall mode. The companion script
\ensuremath{(dz_{\mathrm{chirality}})} records the verification together with the
Nielsen--Ninomiya control (a naive lattice doubles, netting zero) that
shows the codimension-1 mechanism is necessary. Global weak parity
violation (net L \ensuremath{\neq} R coupling) remains open: Nielsen--Ninomiya doubling
pins net chirality to zero on any local Hermitian translation-invariant
lattice, so the L \ensuremath{\neq} R asymmetry must come through a codimension-1
mechanism --- a domain-wall or overlap-fermion construction on the \ensuremath{\hat{R}}
boundary, with the hierarchy direction as the extra dimension. The local
wall mode is now established; what remains open is specifically the
\emph{net global} asymmetry: on a closed hierarchy a wall and an
anti-wall bind modes of opposite handedness that cancel
(Nielsen--Ninomiya reasserting globally), so a net chiral count requires
the hierarchy direction to be effectively one-sided at \ensuremath{\hat{R}}. The existence
and sign of a net asymmetry are established from results already in hand; only
the explicit integer magnitude is deferred. Two inherited facts are decisive.
First, \ensuremath{\hat{R}} is not a generic mass domain wall but a role-promotion map: by its
definition (\S\,2.3) each confirmed bit becomes an agent one level up, a
center-to-vertex, outcome-to-confirmer promotion whose build direction is fixed by
irreversibility (parent to child, never the reverse; Axiom~3, the frozen arrow).
Second, and decisively for the geometric obstruction, Corollary~A.0 established that
the substrate carries \emph{exactly one orientation invariant} --- a single global
handedness --- and this is a fact about the lattice geometry (the sign of the triad
triple product), derived from the recoupling geometry with no reference to the
fermion spectrum. The Nielsen--Ninomiya cancellation requires a handedness-symmetric
lattice, on which every mode of one handedness has a mirror partner of the other;
that symmetry is the theorem's premise. The substrate's geometry denies that
premise at the root: there is no opposite global orientation for a mirror partner to
carry. This identifies a genuine geometric obstruction to the global cancellation
--- the wall and anti-wall that would cancel bind modes referred to a lattice
orientation that is not symmetric. What this does \emph{not} do, and the distinction
is essential, is establish the net \emph{spectral} index: the step from ``the
lattice orientation is single-valued'' (a geometric fact, Corollary~A.0) to ``the
fermion spectrum carries a nonzero net chiral charge, of definite sign'' (a spectral
fact) is an inference that requires the explicit overlap-operator index on the
\ensuremath{\hat{R}} boundary lattice, and that computation is not carried out here. The
framework therefore establishes the \emph{geometric precondition} for a net global
chirality --- \ensuremath{\hat{R}}'s irreversibility gives the one-sidedness, and Corollary~A.0
removes the mirror symmetry that would force cancellation --- while the existence,
sign, and magnitude of the resulting net spectral index remain to be demonstrated by
that index calculation. This is a narrower and more honest claim than a derived net
chirality, and it is worth stating precisely what the framework's results do and do
not require. The corpus's substantive claim about the matter sector is that the
framework produces \emph{chiral} matter rather than vector-like matter --- that left
and right couple differently, as in the Standard Model --- and that claim rests on
the \emph{local} chiral wall mechanism (\S\,8.4.4), which is established, together
with the identification of the geometric obstruction just given (Corollary~A.0
removes the mirror symmetry that a no-go theorem would otherwise use to force the
matter back to vector-like). Neither of those depends on the value of the net global
integer. The exact net chiral count --- whether it is one unit per generation or
another value --- is a well-posed further question, answerable in principle by the
overlap-operator index on the explicit \ensuremath{\hat{R}} boundary lattice, but it is not
load-bearing for the framework's claim that the matter it produces is chiral. It is
therefore recorded as a genuine but non-essential open computation: the local
chirality is derived, the obstruction to global cancellation is identified and is
not circular, and the precise net integer is neither claimed nor required here.

\emph{Neutrino masses and their hierarchy.} As flagged in \S\,7.6.4, the
neutrino mass hierarchy is potentially a probe of level identification
through how each generation's agent-slot anchors into the recursive
embedding. Specific predictions deferred to future work.

\subsubsection*{8.4.5 Higgs Mechanism Structural Correspondence}

The Higgs mechanism of standard electroweak theory --- a scalar field
acquires a non-zero vacuum expectation value, giving mass to gauge
bosons and (through Yukawa couplings) to fermions --- has a structural
correspondence in the TMS framework. We identify the substrate-level
Higgs analog using the gauge structure of \S\,8.3, in which SU(3) is the
continuum lift of the within-event \ensuremath{S_{3}} on agent labels (level 0) and
SU(2) is the \emph{internal gauge} symmetry acting on the A/B label space
of the level-1 doublet. This internal gauge-SU(2) is the object on which
the Higgs mechanism depends. It is a distinct structure from the fermion
spinor phase, which derives from the triadic oddness (\S\,4.1, Lemma A):
the electroweak SU(2) here acts on the sublattice-membership label, and it
is this gauge symmetry that the Higgs analog uses below.

\begin{theorem}[Existence of a Substrate-Inherited Scalar Background]
Every \ensuremath{\hat{R}-\text{produced}} substrate inherits, in addition to the
inversion-odd structured break of \S\,5.3, a generically non-zero
inversion-even scalar background --- the A1g (trivial irrep) component
of the parent biography's projection onto the polarization field
decomposition (\S\,4.2).
\end{theorem}

\begin{proof}
By \S\,4.2, the polarization field decomposes under \ensuremath{O_{h}} on
the FCC nearest-neighbor space into five irreps including the trivial
A1g. The biographical record \ensuremath{\mathcal{B}(R),} projected onto this decomposition,
contains generically non-zero components in all five irreps. The
zero-A1g case is measure-zero and not preserved under generic substrate
dynamics; direct simulation of the TMS triadic confirmation rule
(companion materials) confirms a non-zero A1g component in the late-time
record.

\end{proof}

\textbf{The Higgs is \ensuremath{d_{z},} not the A1g background.} The natural first
guess --- that the inversion-even A1g background plays the role of the
Higgs VEV --- fails by symmetry. A1g is the trivial irrep of the full
\ensuremath{O_{h}} point group, and a fortiori a singlet under every subgroup of \ensuremath{O_{h},}
including both the within-event \ensuremath{S_{3}} that lifts to SU(3) color and the A/B
doublet structure on which SU(2) acts (\S\,8.3). A scalar in the trivial
representation of any gauge group does not contribute to gauge-boson
mass terms by the standard \ensuremath{\langle H\dagger T^{a}} \ensuremath{H\rangle} = 0 mechanism. The A1g
background therefore gives mass to no gauge boson.

The Higgs of the TMS framework is the inversion-odd structured break
\ensuremath{d_{z}} of \S\,5.3. The break \ensuremath{d_{z}} = \ensuremath{\langle \mathcal{B}_{A}} \ensuremath{-} \ensuremath{\mathcal{B}_{B}\rangle} is by construction
non-trivial under A \ensuremath{\leftrightarrow} B exchange, placing it in the SU(2)-doublet
direction of the level-1 sublattice doublet structure. It is therefore a
non-trivial SU(2) representation, and its non-zero spatial-mean
component plays the role of an SU(2)-doublet Higgs vacuum expectation
value, measured directly in the FCC + A/B biographical simulation
(companion materials) as a \ensuremath{6\sigma} non-zero, time-stable, spatially
structured mean: it breaks SU(2) \ensuremath{\times} U(1) to a U(1) subgroup determined by
the specific direction of \ensuremath{d_{z}} in the doublet space, giving mass to the
SU(2) gauge bosons that do not preserve the unbroken U(1) subgroup.

\textbf{Color stays massless.} The within-event \ensuremath{S_{3}} that lifts to SU(3)
color acts on the three agent labels within a single confirmation event
(\S\,8.3), independently of the level-1 A/B sublattice structure on which
\ensuremath{d_{z}} lives. \ensuremath{d_{z}} transforms trivially under this within-event \ensuremath{S_{3}:} it is
a property of the sublattice membership at level 1, not of the agent
labels within an event at level 0. By the same \ensuremath{\langle H\dagger T^{a}} \ensuremath{H\rangle} = 0
mechanism, an SU(3)-trivial scalar gives no mass to the SU(3) gauge
bosons. The color gauge bosons therefore remain massless, consistent
with their observed character.

\textbf{Two roles, one object.} The structural identification \ensuremath{d_{z}} =
Higgs assigns a single substrate object two distinct field-theoretic
roles: (a) the band-structure carver that produces fermion sector
content (\S\,5.3, \S\,8.4.4), and (b) the SU(2)-doublet VEV that breaks
electroweak symmetry. Both roles emerge from the same band-structure
analysis of \ensuremath{d_{z}} acting on the A/B sublattice doublet. This is a single
substrate object playing both roles, verified by the existence and
stability of \ensuremath{d_{z}} in the substrate's own dynamics, not a correspondence
between two separate substrate quantities and two separate Standard
Model operators.

\textbf{Status of the A1g background.} The non-zero A1g scalar
background established by the theorem above exists but does not play the
role of the electroweak Higgs. Its substrate-level role is open:
candidate identifications include a \ensuremath{U(1)_{Y}} hypercharge background, a
substrate-level vacuum energy contribution, or a quantity without an
immediate Standard Model analog at this level of the framework.

\textbf{Status of Results for \S\,8.4.5:}

\begin{itemize}
\item
  \emph{Derived (under \S\,8.3's SU(2)-from-A/B-doublet structural
  correspondence):} \ensuremath{d_{z}} exists, is time-stable, and is spatially
  structured in the substrate's own dynamics. The identification \ensuremath{d_{z}} =
  Higgs VEV holds modulo the \S\,8.3 continuum lift.
\end{itemize}

\begin{itemize}
\item
  \emph{Derived:} The existence of a non-zero inversion-even scalar
  background (A1g component) inherited via \ensuremath{\hat{R}} from the parent biography.
\end{itemize}

\begin{itemize}
\item
  \emph{Structural correspondence:} The dual role of \ensuremath{d_{z}} as both
  fermion-content carver (\S\,5.3) and electroweak symmetry breaker (this
  section), as one substrate object in two field-theoretic roles. The
  fermion-carver role is established in \S\,5.3; the EW-breaker role is
  established here.
\end{itemize}

\begin{itemize}
\item
  \emph{Prediction / Conjecture:} The specific direction of \ensuremath{d_{z}} in the
  SU(2)-doublet that determines the unbroken \ensuremath{U(1)_{\mathrm{EM}}} subgroup; the W
  and Z mass values; the substrate-level role of the A1g scalar
  background. All co-located with the keystone idealization 2
  biographical simulation and the \S\,8.3 representation-theoretic mapping.
\end{itemize}

\subsubsection*{8.4.6 The Fermion Roster from Anomaly Freedom}

With the gauge algebra Derived (\S\,8.3) and the chirality mechanism in
place (\S\,8.4.4), the matter content can be assembled. The substrate fixes
the representation structure of one generation: quarks carry colour,
participating directly in the three-agent triadic confirmation, and so
transform as SU(3) triplets (confined, \S\,8.3); leptons are triadically
neutral and transform as colour singlets; the left-handed components are
the chiral A/B doublet bound to the \ensuremath{d_{z}} wall (\S\,8.4.4) and so transform
as SU(2) doublets, while the right-handed components do not couple to
the wall and are SU(2) singlets. In left-handed Weyl form the multiplets
are Q = (3, 2, \ensuremath{Y_{Q}),} \ensuremath{u_{c}} = \ensuremath{(\bar{3},} 1, \ensuremath{Y_{u}),} \ensuremath{d_{c}} = \ensuremath{(\bar{3},} 1, \ensuremath{Y_{d}),} L = (1, 2,
\ensuremath{Y_{L}),} and \ensuremath{e_{c}} = (1, 1, \ensuremath{Y_{e}),} with the hypercharges Y left open at this
stage. This representation assignment is a structural correspondence:
the substrate motivates each rep, but the identification of a substrate
excitation with an SM multiplet is not a theorem.

The hypercharges, however, are not free. A chiral gauge theory is
quantum-mechanically consistent only if all gauge and gravitational
anomalies cancel; an uncancelled anomaly breaks gauge invariance at one
loop and renders the theory ill-defined. Because \S\,8.4.4 establishes that
the substrate produces a genuinely chiral gauge theory, anomaly-freedom
is not an optional feature to be imposed but a consistency requirement
the substrate must already satisfy. Imposing it on the rep structure
above fixes the hypercharges. The mixed \ensuremath{SU(2)^{2}--U(1)} condition gives
\ensuremath{3Y_{Q}} + \ensuremath{Y_{L}} = 0; the gravitational U(1) condition then gives \ensuremath{Y_{e}} =
\ensuremath{6Y_{Q};} the mixed \ensuremath{SU(3)^{2}--U(1)} condition gives \ensuremath{Y_{u}} + \ensuremath{Y_{d}} = \ensuremath{-2Y_{Q};} and
the cubic \ensuremath{U(1)^{3}} condition fixes the remaining split, its only solutions
being \ensuremath{Y_{u}} = \ensuremath{-4Y_{Q},} \ensuremath{Y_{d}} = \ensuremath{2Y_{Q}} or the \ensuremath{up\leftrightarrow \text{down}} reversal. Setting the
single overall normalization \ensuremath{Y_{Q}} = 1/6 reproduces the Standard-Model
hypercharges exactly: Y = (1/6, \ensuremath{-2/3,} 1/3, \ensuremath{-1/2,} 1) for (Q, \ensuremath{u_{c},} \ensuremath{d_{c},} L,
\ensuremath{e_{c}).} The cubic condition has isolated roots, not a continuum, so the
split is genuinely pinned. Every hypercharge is therefore fixed up to
one overall normalization and the convention of which weak partner is
named ``up'' --- no free numerical parameter, consistent with the
volume's no-tuning stance.

One dependency is stated honestly rather than hidden. The anomaly-cancellation
conditions used here are the one-loop triangle-anomaly polynomials of continuum
quantum field theory; the substrate does not yet derive them internally, but
inherits them as the consistency requirement any chiral gauge theory it produces
in the continuum limit must satisfy. This is therefore a \emph{structural
correspondence with a continuum-QFT dependency}: the hypercharges are pinned
\emph{given} that the continuum theory obeys the standard anomaly conditions, and
deriving those conditions from substrate primitives --- rather than importing
them --- is an open item. What is internal is the chirality (\S\,8.4.4) that makes
anomaly-freedom binding in the first place; what is imported is the anomaly
polynomial whose vanishing then fixes the charges.

The content that results is exactly one Standard-Model generation:
fifteen Weyl fermions (six in Q, three each in \ensuremath{u_{c}} and \ensuremath{d_{c},} two in L, one
in \ensuremath{e_{c}),} replicated three times by the generation count of \S\,8.4.2. Two
consequences are worth stating. First, the right-handed neutrino \ensuremath{\nu_{c}} =
(1, 1, 0) is a total gauge singlet that contributes to no anomaly; it is
therefore neither required nor forbidden by consistency --- the
framework predicts it as an optional sterile state, matching its
anomalous absence-or-presence in the Standard Model. Second, the content
is the minimal anomaly-free \emph{chiral} completion of the substrate
rep structure: dropping the right-handed electron, recolouring the
leptons, altering a hypercharge, or rendering the content vector-like
each breaks anomaly-freedom (and the vector-like case additionally
destroys the parity violation the substrate produces). The
Standard-Model generation is thus singled out, not assumed.

\textbf{Status of Results for \S\,8.4.6.} \emph{Derived} (given the rep
structure and the chirality of \S\,8.4.4): anomaly-freedom fixes all
hypercharges to the SM values up to overall normalization and the
\ensuremath{up\leftrightarrow \text{down}} label, with no free parameter; the optional \ensuremath{\nu_{c}} as a
consistency-invisible singlet. \emph{Structural correspondence}: the
assignment of substrate excitations to the colour-triplet /
colour-singlet and weak-doublet / weak-singlet representations.
\emph{Open}: the overall hypercharge normalization (the analog of the
level index k elsewhere in the volume), and the net global parity
violation flagged in \S\,8.4.4, on which the chirality --- and therefore
the requirement of anomaly-freedom itself --- ultimately rests.

\subsubsection*{8.4.7 The Periodic-Table Skeleton: Subshell Capacities,
Madelung Ordering, and the Noble-Gas Atomic Numbers}

The same nearest-neighbor \ensuremath{O_{h}} decomposition that organizes the field
sectors (\S\,4.2), combined with the exclusion principle derived in \S\,4.1
(Corollary A.2) and the single flop-clock of the substrate dynamics,
fixes the combinatorial skeleton of atomic structure --- the subshell
capacities, the filling order, and the noble-gas atomic numbers --- as
pure counting, with no free parameter. What follows is the architecture of
the periodic table; the dimensionful energies that set the actual
term values are not claimed here, exactly as the mass \emph{values} are
deferred in \S\,8.4.3.

\emph{Subshell capacities \ensuremath{2(2\ell+1)}.} The angular content of an atomic
shell is the set of standing polarization modes of definite angular momentum on
the sphere of confirmation directions. The number of independent modes of
angular momentum \ensuremath{\ell} is the dimension of the \ensuremath{\ell}-th irreducible
representation of the rotation group, \ensuremath{2\ell + 1} --- and Paper 5 \S\,4.1
derives that the triadic recoupling reconstructs the full rotation group SO(3),
so this count is available at every \ensuremath{\ell}, not merely the small values a single
shell resolves. The count \ensuremath{2\ell + 1} is the number of independent degree-\ensuremath{\ell}
harmonic polynomials on \ensuremath{\mathbb{R}^3} (the angular standing modes at that degree),
which is \ensuremath{1, 3, 5, 7, 9, \dots} for \ensuremath{\ell = 0, 1, 2, 3, 4, \dots}. Multiplying by
the two-fold spin polarity (the derived spin factor of \S\,4.1, Corollary A.2)
gives the subshell capacities \ensuremath{2(2\ell+1)}: \ensuremath{s: 2}, \ensuremath{p: 6}, \ensuremath{d: 10}, \ensuremath{f: 14},
\ensuremath{g: 18}. The twelve nearest-neighbor directions of a single FCC shell resolve
only the low-angular-momentum part of this tower: under the subduction
\ensuremath{SO(3) \downarrow O_{h}}, \ensuremath{L = 0} subduces to \ensuremath{A_{1}g}, \ensuremath{L = 1} to \ensuremath{T_{1}u}, and
\ensuremath{L = 2} to the pair \ensuremath{E_{g} \oplus T_{2}g} (the standard five-fold \ensuremath{d}-shell
splitting), reproducing \ensuremath{s: 2}, \ensuremath{p: 6}, \ensuremath{d: 10} and matching the harmonic count
\ensuremath{1, 3, 5} exactly. The \ensuremath{f}-shell capacity 14 is the seven angular modes of
\ensuremath{\ell = 3} --- present in the reconstructed group though beyond the single
shell's resolution --- times the spin factor two. No value is fitted.

\emph{Madelung \ensuremath{(n+\ell)} ordering.} The substrate advances one lattice step
per flop tick; radial closure of a shell costs on the order of \ensuremath{n} ticks
and angular closure on the order of \ensuremath{\ell} ticks, counted in the \emph{same}
unit because there is one clock. The filling cost is therefore additive in
\ensuremath{n} and \ensuremath{\ell}, which is the Madelung rule (Madelung 1936; Klechkowski 1962):
subshells fill in order of
increasing \ensuremath{n + \ell}. This dissolves the standard puzzle of why the radial
and angular quantum numbers enter the filling rule with equal weight ---
they are ticks of one clock. The tie-break among equal \ensuremath{n + \ell} (fill
lower \ensuremath{n} first) follows from loop compactness: the more radially compact
loop is more tightly bound. That tie-break is a sound physical tendency
rather than a rigorous consequence, and is labeled Derived-conditional
accordingly.

A point of honesty about what this result is and is not. In ordinary atomic
physics the \ensuremath{n+\ell} ordering is the empirical signature of electron
\emph{screening}: in a pure Coulomb \ensuremath{1/r} potential the energy depends on \ensuremath{n}
alone (the accidental degeneracy), which would fill strictly by \ensuremath{n} and give the
wrong closures --- the third noble gas would fall at \ensuremath{Z=28} rather than 18. It is
the screened many-electron potential that lifts the degeneracy into the \ensuremath{n+\ell}
order actually observed. The substrate reproduces that same \ensuremath{n+\ell} order from a
different and more primitive source: the additivity of radial and angular closure
costs under a single clock. That the clock-counting yields precisely the screened
ordering, rather than the bare-Coulomb ordering, is a genuine and non-trivial
structural match --- pure \ensuremath{1/r} does \emph{not} give the periodic table, and the
substrate does. Whether the single-clock additivity is, at a deeper level, the
substrate's expression of screening physics, or an independent route that happens
to coincide with it, is not settled here. The honest status is therefore a
derived counting rule that matches the empirically screened order exactly, with
the physical identification of the two mechanisms marked as open.

\emph{Noble-gas atomic numbers and period lengths.} Filling the derived
capacities in the derived order gives cumulative electron counts closing
each period at
\begin{center}
\ensuremath{2,\ 10,\ 18,\ 36,\ 54,\ 86,\ 118}
\end{center}
--- exactly the atomic numbers of helium, neon, argon, krypton, xenon, radon,
and oganesson, the heaviest element yet synthesized --- with period lengths
\ensuremath{2, 8, 8, 18, 18, 32, 32}. The two length-32 periods are exactly the two that
contain an \ensuremath{f}-block: the \ensuremath{4f} subshell fills fourteen electrons across the
lanthanides (ending at \ensuremath{Z = 70}) and the \ensuremath{5f} subshell fourteen across the
actinides (ending at \ensuremath{Z = 102}). This is pure counting of the two preceding
results, with no adjustable input: the real atomic numbers of the entire
periodic table --- heavy elements, lanthanides, and actinides included ---
emerging from \ensuremath{\bar{B}} = \ensuremath{3\bar{\alpha}}. It is the strongest concrete arithmetic
anchor in the matter sector.

\textbf{Status of Results for \S\,8.4.7.} \emph{Derived}: the subshell
capacities \ensuremath{2(2\ell+1)} --- \ensuremath{s: 2}, \ensuremath{p: 6}, \ensuremath{d: 10}, \ensuremath{f: 14} --- from the SO(3)
angular-mode count (Paper 5 \S\,4.1) times the spin factor; the Madelung
\ensuremath{n + \ell} additivity from the single flop-clock; and the noble-gas atomic
numbers \ensuremath{(2, 10, 18, 36, 54, 86, 118)} and period lengths
\ensuremath{(2, 8, 8, 18, 18, 32, 32)}, including the \ensuremath{4f} lanthanide and \ensuremath{5f} actinide
blocks, by counting. \emph{Derived-conditional}: the lower-\ensuremath{n} Madelung
tie-break, a sound compactness tendency. \emph{Open}: the dimensionful orbital
energies (the analog of the mass values of \S\,8.4.3). The framework derives the
periodic table's full skeleton --- capacities, ordering, the noble gases, and the
\ensuremath{f}-block placement across all seven periods --- not its spectroscopic energies.

\subsection*{8.5 Status Summary}

The substrate-level \ensuremath{Z_{2}} \ensuremath{\times} \ensuremath{Z_{2}} at each level and the cross-level commutator
mechanism from the sublattice bifurcation are \emph{Derived} --- they
follow from the [[4,2,2]] stabilizer structure of Paper 1 \S\,V and
the bifurcation of \S\,6.1 respectively. Given the QCA-to-QFT continuum
lift of \S\,4.1, the colour gauge group is now \emph{Derived} (\S\,8.3): the
algebra is su(3) by a non-circular closure of substrate-built generators
with no overshoot, simple and compact; the global group is SU(3) ---
fixed by the fundamental triality of the substrate-real agents ---
rather than \ensuremath{SU(3)/\mathbb{Z}_{3};} and confinement is the \ensuremath{\text{unbroken}-\mathbb{Z}_{3}-\text{centre}}
statement. The level-1 SU(2) factor and its U(1) trace-mode partner are
now Derived on the same footing (\S\,8.3): su(2) closes from the A/B
doublet transport with no overshoot, with global form SU(2) rather than
SO(3) by the faithful fundamental doublet. What remains reserved is the
electroweak gluing --- the quotient SU(2) \ensuremath{\times} U(1) versus (SU(2) \ensuremath{\times}
\ensuremath{U(1))/\mathbb{Z}_{2}} \ensuremath{\cong} U(2), fixed by hypercharge --- the chirality \ensuremath{SU(2)_{L}}
supplied by the \ensuremath{d_{z}} break (\S\,8.4.4), and the three-generation fermion
content; these are \emph{Structural correspondences} in the sense
developed throughout the volume: the TMS framework produces the
structural skeleton of the SM, with the explicit
representation-theoretic mapping reserved for those factors. The
quantitative consequences (mass spectrum, mixing angles, etc.) are
\emph{Predictions / Conjectures} whose empirical bridges require
subsequent simulation work. One quantitative layer below the gauge group
is, however, now fixed without tuning: given the substrate rep
structure, the requirement that the chiral theory (\S\,8.4.4) be
anomaly-free Derives the entire one-generation fermion roster and every
hypercharge to the Standard-Model values, up to one overall
normalization and the \ensuremath{up\leftrightarrow \text{down}} label, with the right-handed neutrino
predicted as an optional gauge singlet (\S\,8.4.6). What stays open at the
numerical level is the overall hypercharge normalization and, beneath
it, the net global parity violation on which the chirality rests.

\section*{IX. Substrate-Forced Predictions and Empirical Bridges}
\addcontentsline{toc}{section}{IX. Substrate-Forced Predictions and Empirical Bridges}

The structural content of Papers 1-5 produces a number of predictions
about the observable universe that follow from substrate primitives
without free parameters. We organize them by epistemic status:
\emph{formally derived} (proof chain complete within stated framework),
\emph{substrate-forced} (follows from primitives but with a stated
structural assumption not yet a theorem), and \emph{falsifiable}
(specific quantitative form that observation could decisively test).
Each prediction is annotated with its substrate provenance.

\subsection*{9.1 Gauge Group Structure (substrate-forced)}

The Standard Model gauge group SU(3) \ensuremath{\times} SU(2) \ensuremath{\times} U(1) emerges from the
triadic confirmation primitive at distinct levels of the substrate
hierarchy.

\emph{SU(3) at level 0.} Every confirmation event involves three agents
closing around one bit (B-bar = 3 alpha-bar, Paper 1 Section III). The
\ensuremath{S_{3}} permutation symmetry of agent labels lifts in the QCA-to-QFT
continuum limit (Section 4.1) to continuous local gauge symmetry on the
3-dimensional internal space. \ensuremath{S_{3}} is the Weyl group of SU(3); the
standard gauging produces SU(3) as the continuum symmetry. The group
identity is fixed not by this label-permutation reading alone but by the
push-forward derivation of \S\,8.3: the second-order transport carries the
complex amplitude z = \ensuremath{\psi} + \ensuremath{i\Pi ,} whose substrate-built generators close to
su(3) (anti-Hermitian, traceless, simple, compact, dimension 8, no
overshoot) and whose fundamental-triality matter content fixes the
global group as SU(3) with an unbroken \ensuremath{\mathbb{Z}_{3}} centre (confinement). With
that derivation the colour factor is Derived under the \S\,4.1 lift, not
merely substrate-forced.

\emph{SU(2) at level 1.} The level-1 substrate carries two
interpenetrating FCC sublattices (the A and B sublattices of Sections
6.1.1, 6.1.2). The doublet structure under inversion through the level-2
centroid produces \ensuremath{SU(2)_{L}} electroweak as the gauging of the sublattice
exchange.

\emph{U(1) at level 0.} The overall phase of the bilinear
role-correlation operators on the 3-agent space (the trace mode of
Section 9.2) carries U(1).

\textbf{Provenance:} Paper 1 Section III, Paper 5 Sections 4.1, 6.1.1,
6.1.2, 8.3. \textbf{Status:} Substrate-forced under the QCA-to-QFT lift
(lift itself substrate-internal per Section 4.1 with Lemmas A and B).

\subsection*{9.2 Gluon Count (formally derived)}

For a triadic event with 3 agents, the bilinear role-correlation
operators number 3 x 3 = 9. One is the trace \ensuremath{(U(1)_{\mathrm{baryon}}).} The
remaining 8 are traceless and form the SU(3) adjoint generators. The 8
gluons of QCD correspond to these 8 traceless bilinears. The count
\ensuremath{3^{2}} - 1 = 8 is forced once the substrate primitive is triadic.

\textbf{Provenance:} Paper 1 Section III, Paper 5 Section 8.3.
\textbf{Status:} Formally derived.

\subsection*{9.3 Fractional Electric Charge (formally derived)}

Three agents share one bit's U(1) charge in a triadic confirmation
event. By \ensuremath{S_{3}} symmetry no agent is privileged, so each carries 1/3 of
the bit's charge. Anomaly cancellation (3 \ensuremath{Y_{q}} + \ensuremath{Y_{L}} = 0, with the 3
being the number of colors --- i.e., the same 3 as triadic confirmation)
combined with \ensuremath{Y_{L}} = -1 forces \ensuremath{Y_{q}} = +1/3, then Q = \ensuremath{T_{3}} + Y/2 gives Q
= +/-1/2 + 1/6, producing +2/3 (up) or -1/3 (down). The ``1/3'' is
literally ``one of three'' --- the agent count itself.

With the full hypercharge assignment now fixed by anomaly-freedom
(\S\,8.4.6), the electric charge Q = \ensuremath{T_{3}} + Y of every fermion follows with
no further input: u, c, t at +2/3; d, s, b at \ensuremath{-1/3;} neutrinos at 0;
charged leptons at \ensuremath{-1.} Two consequences are worth drawing out. First,
charge is quantized in units of 1/3, and the denominator is the colour
count three --- the same three agents of \ensuremath{\bar{B}} = \ensuremath{3\bar{\alpha};} leptons, being colour
singlets, carry integer charge. Second, the proton (uud) has charge
exactly +1 and the electron exactly \ensuremath{-1,} so a hydrogen atom is exactly
neutral. This exactness is not numerical luck: the proton--electron
cancellation is the same per-generation charge sum that the
gravitational--U(1) anomaly condition sets to zero, so charge
quantization and the neutrality of atoms are corollaries of the very
consistency requirement that fixed the roster. The measured bound
\textbar \ensuremath{Q_{p}} + \ensuremath{Q_{e}}\textbar/e \textless{} \ensuremath{10^{-21}} is, within the
framework, an exact algebraic identity rather than a fine-tuned
coincidence.

\textbf{Provenance:} Paper 1 Section III, Paper 5 Section 8.3.
\textbf{Status:} Formally derived modulo the import of anomaly
cancellation from standard QFT.

\subsection*{9.4 Three-Generation Fermion Count (formally derived)}

Confirmation events are intrinsically 2-dimensional: three agents close
around one bit in a triangular configuration on the manifold \ensuremath{M_{\mathrm{alpha}}}
(Paper 1 Section III, B-bar = 3 alpha-bar). The 3D FCC structure (Paper
1 Section VI) is the embedding lattice for the events, not the events
themselves. The minimum-closure count for distinguishable agent-slots is
therefore the 2D minimum: 3.

Three generations correspond to three distinguishable agent-slots in the
primitive triadic event, instantiated as matter carriers under recursive
embedding through the hierarchy. The ``3 of three colors'' and the ``3
of three generations'' are the same 3, distinguished by which
symmetry-axis is projected. The framework predicts not 4 generations
(which would be the 3D tetrahedral minimum B-bar = 4 \ensuremath{\text{alpha}-\text{bar}^{3})}
because confirmation events live in 2D, not 3D.

Independent corroboration: Jegerlehner (2019) argues that three fermion
generations are required to make CP violation and baryogenesis emerge
naturally in low-energy effective Standard Model scenarios --- a
different reason for the count 3, arriving at the same number.

\textbf{Provenance:} Paper 1 Section III, Paper 1 Section VI, Paper 5
Sections 2.5, 8.4. \textbf{Status:} Formally derived (count), via the
\S\,8.4.1 recursion-preservation theorem and the \S\,8.4.2 arrow-tower
saturation; the identification of the three structural carriers with the
three observed fermion generations remains a structural conjecture
(\S\,8.4.2, \S\,8.4.3).

\subsection*{9.5 Cross-Level Gravitational Coupling (formally derived)}

The cross-level acceleration-like ratio characterizing gravitational
coupling at the substrate level is \ensuremath{G_{\mathrm{skeleton}}} =
\ensuremath{(L_{p}}\textsuperscript{\ensuremath{(k+1)/L_{p}}}(k)) /
\ensuremath{(T_{p}}\textsuperscript{\ensuremath{(k+1)/T_{p}}}\ensuremath{(k))^{2}} = (3 sqrt 2)/36, with T-ratio
= 6 forced by PGE saturating the c-bound inequality of Sections 5.3 +
6.2 to equality. The framework predicts the structural ratio between
successive levels and the internal self-consistency of Planck units at
each level. The SI value of G depends on which level we sit at --- see
Section 9.10.

\textbf{Provenance:} Paper 1 Section VI, Paper 5 Sections 5.3, 6.1.3,
6.2, 7.5, 7.6. \textbf{Status:} Formally derived as a cross-level
structural ratio.

\subsection*{9.6 Lorentz-Violation Residual Scaling (falsifiable)}

The substrate-internal continuum-symmetry argument (Section 4.1, Lemma
B) predicts residual Lorentz violation suppressed by
\textbf{\ensuremath{(E/E_{\mathrm{Planck}})^{2}}} for all inversion-symmetric sectors ---
scalar and gauge-vector alike --- because lattice inversion symmetry
forbids odd-momentum terms and the leading anisotropy enters at the
quartic-in-momentum \ensuremath{O_{h}} invariant. A stronger
\textbf{\ensuremath{(E/E_{\mathrm{Planck}})^{1}}} contribution is predicted only in the
inversion-broken matter sector (\S\,5.3), with magnitude set by the
parent-record asymmetry. Future cosmic-ray observations measuring a
fractional Lorentz-violation exponent inconsistent with these
predictions --- quadratic in the boson sectors, with any linear signal
confined to and scaling with the matter-sector asymmetry --- would
falsify the framework's QCA-to-QFT lift.

\textbf{Provenance:} Paper 1 Section VI, Paper 5 Section 4.1 (Lemma B).
\textbf{Status:} Falsifiable. Contrast with Wetterich (2022), whose
cellular-automaton spinor-gravity construction predicts exact Lorentz
symmetry on the discrete level for a non-FCC lattice; observation will
distinguish the two substrate hypotheses.

\subsection*{9.7 GUP Scaling with Hierarchical Depth (falsifiable)}

The Generalized Uncertainty Principle deformation parameter beta scales
with causal depth d rather than being universal: beta(d) = \ensuremath{beta_{0}} *
f(d, \ensuremath{K_{\mathrm{psi},\mathrm{local}},} \ensuremath{K_{\mathrm{psi},\mathrm{surround}},} \ensuremath{K_{\mathrm{psi},\mathrm{subsystem}},} k), with closed
form in Appendix C. Distinguishes TMS from string-theory and
loop-quantum-gravity GUP predictions; specific experimental targets
include quantum decoherence measurements at varying gravitational
potentials (matter-wave interferometry at varying altitudes, atomic
clock decoherence comparisons between ground-level and ISS-level
platforms, gravitational-redshift-corrected decoherence times).

\textbf{Provenance:} Paper 1 Section VII, Paper 5 Section 6.4, Appendix
C. \textbf{Status:} Falsifiable.

\subsection*{9.8 Dark Sector Identification (candidate, structurally motivated)}

Within-level vacuum stress-energy decomposes under \ensuremath{S_{3}} into 7 vanishing
modes and 2 surviving residual modes: \ensuremath{T^{(k)}} (bit-channel) and
\ensuremath{R^{(k)}} (relation-channel). These have distinct distribution
properties: \ensuremath{T^{(k)}} concentrates with matter (clusters, gravitates like
dark matter); \ensuremath{R^{(k)}} is distributed across substrate structure
(smooth, field-like, behaves like dark energy). Predicted ratio from
substrate combinatorics: relation-channel : bit-channel = 3 : 1,
matching the foundational substrate ratio B-bar = 3 alpha-bar. Observed
ratio: dark energy : dark matter \textasciitilde{} 2.65 : 1.
Leading-order prediction differs from observation by
\textasciitilde12\%; higher-order corrections in mode-weighting could
move the prediction in either direction and have not been calculated.

\textbf{Provenance:} Paper 5 Sections 7.6.4, 8.3. \textbf{Status:}
Candidate, structurally motivated. The 12\% gap is a real residual
flagged as open.

\subsection*{9.9 Recursive Vacuum Absorption (formally derived)}

Under the Level-Advance Relation (Section 6.3.0), within-level vacuum
residuals are reabsorbed as non-vacuum structure at higher levels:
\ensuremath{T^{(k)}} becomes matter at level k+1, \ensuremath{R^{(k)}} becomes substrate
structure at level k+1. The ``sub-Planckian vacuum integration'' form of
the cosmological constant problem does not arise in TMS.

\textbf{Provenance:} Paper 5 Sections 6.3.0, 7.6.4. \textbf{Status:}
Formally derived as corollary of the Level-Advance Relation.

\subsection*{9.10 Level-Identification (empirical bridge)}

The framework predicts cross-level structural ratios and level-internal
Planck consistency, but not the SI values of fundamental constants
directly. Asking ``what is G in SI?'' reduces to ``what level k are we
at?'' Empirical handles on k include: CMB power spectrum signature
(Section 7.6.4); Hubble constant and dark-energy density ratio to atomic
clock rates; gravitational-wave dispersion at future high-frequency
detectors; neutrino mass hierarchy as probe of recursive embedding
depth; black-hole interior evaporation spectrum deviations from standard
predictions. Under steady-state vacuum dynamics (candidate, Section
7.6.4 Result 3), observed Lambda corresponds to k \textasciitilde{} 79.

\textbf{Provenance:} Paper 5 Sections 6.3, 7.6. \textbf{Status:}
Empirical bridge --- primary point of contact with cosmological data.

\subsection*{9.11 CMB Anisotropy Structural Signature (substrate-forced)}

The level-1 saturation cascade produces a discrete-substrate fingerprint
in the CMB power spectrum at small angular scales, distinguishable from
inflationary alternatives. Two structural features predicted: structural
maximum at characteristic late-phase configuration scale (empirically
consistent with first acoustic peak at l \textasciitilde{} 220,
structural identification requires simulation), and super-horizon
correlations from parent-level causal structure (empirically consistent
with observed super-horizon correlations that motivate inflationary
cosmology).

\textbf{Provenance:} Paper 5 Sections 5.1-5.2. \textbf{Status:}
Substrate-forced; specific signature form requires numerical simulation.

\subsection*{9.12 Open Items}

Items where the framework has structural commitments but the formal
derivation is incomplete:

\begin{itemize}
\tightlist
\item
  Anomaly cancellation as consequence of QCA-to-QFT consistency (used in
  Section 9.3)
\end{itemize}

\begin{itemize}
\tightlist
\item
  Formal SU(2)-vs-SO(3) representation verification for Section 4.1
  Lemma A
\end{itemize}

\begin{itemize}
\tightlist
\item
  Steady-state vacuum dynamics derivation from substrate primitives
  (closes Section 9.8 Result 3 conditionality)
\end{itemize}

\begin{itemize}
\tightlist
\item
  Higher-order mode-weighting calculation for dark sector 3:1 ratio
  (resolves Section 9.8 gap)
\end{itemize}

\begin{itemize}
\tightlist
\item
  Consistency audit of ``agents = energy-side, bits = mass-side''
  foundational identification against Papers 1-4
\end{itemize}

\begin{itemize}
\tightlist
\item
  Chirality / parity violation in \ensuremath{SU(2)_{L}}
\end{itemize}

\begin{itemize}
\tightlist
\item
  Higgs mechanism: the structural correspondence \ensuremath{(d_{z}} as the
  SU(2)-doublet vacuum expectation value) is established in Section
  8.4.5; open are the W and Z mass values, the precise \ensuremath{U(1)_{\mathrm{EM}}}
  embedding, and the substrate-level role of the A1g scalar background
\end{itemize}

\begin{itemize}
\tightlist
\item
  Asymptotic freedom and running couplings from triadic-confirmation
  behavior at different substrate scales
\end{itemize}

\begin{itemize}
\tightlist
\item
  Specific numerical CMB signature form (Section 9.11)
\end{itemize}

\subsection*{9.13 Summary Table}

\begin{longtable}[]{@{}
  >{\raggedright\arraybackslash}p{(\columnwidth - 4\tabcolsep) * \real{0.3333}}
  >{\raggedright\arraybackslash}p{(\columnwidth - 4\tabcolsep) * \real{0.3333}}
  >{\raggedright\arraybackslash}p{(\columnwidth - 4\tabcolsep) * \real{0.3333}}@{}}
\toprule\noalign{}
\begin{minipage}[b]{\linewidth}\raggedright
Prediction
\end{minipage} & \begin{minipage}[b]{\linewidth}\raggedright
Status
\end{minipage} & \begin{minipage}[b]{\linewidth}\raggedright
Substrate Origin
\end{minipage} \\
\midrule\noalign{}
\endhead
\bottomrule\noalign{}
\endlastfoot
SU(3) color from triadic \ensuremath{S_{3}} & Substrate-forced & 3 agents per
confirmation \\
8 gluons & Formally derived & \ensuremath{3^{2}} - 1 traceless bilinears \\
Quark charges +/-1/3, +/-2/3 & Formally derived* & 1 bit charge / 3
agents + anomaly \\
3 generations & Derived (count) & 2D minimum + recursion-preservation
(\S\,8.4.1) \\
\ensuremath{SU(2)_{L}} from sublattice doublet & Substrate-forced & A/B inversion
structure \\
3 weak bosons (W+/-, Z) & Formally derived & \ensuremath{2^{2}} - 1 traceless
bilinears \\
G structural skeleton (3 sqrt 2)/36 & Formally derived & Sections 5.3 +
PGE + 6.1.3 \\
Lorentz-violation residual \ensuremath{(E/E_{\mathrm{Planck}})^{2}} (boson sectors);
\ensuremath{(E/E_{\mathrm{Planck}})^{1}} only if P-broken (matter) & Falsifiable & Inversion
symmetry + quartic \ensuremath{O_{h}} invariant \\
GUP beta scaling with depth & Falsifiable & Paper 1 Section VII \\
Dark energy / dark matter \textasciitilde{} 3:1 & Candidate & Substrate
combinatorics \\
Recursive vacuum absorption & Formally derived & Level-Advance
Relation \\
Level-identification & Empirical bridge & Sections 6.3, 7.6 \\
\end{longtable}

*modulo imported anomaly cancellation

\section*{X. Conclusion: From Circles to a Universe}
\addcontentsline{toc}{section}{X. Conclusion: From Circles to a Universe}

Volume 1 of the Triadic Measurement Substrate has, across five papers,
derived a generative chain from geometric primitives to cosmological
initial conditions. The chain has no free \emph{dimensionless} numerical
parameters and no extraneous axioms: every dimensionless quantity is
forced by the five axioms and the Principle of Generative Economy, and a
single dimensionful input (one action scale) sets the physical scale.

Paper 1 established the geometric and ontological architecture: the
sphere primitive of Axiom 2, the hexagonal manifold \ensuremath{M\alpha ,} the triadic
confirmation rule, the \ensuremath{1\to 4} multiplier, FCC necessity, the
[[4,2,2]] stabilizer structural correspondence, the
identification \ensuremath{\ell_{0}} \ensuremath{\equiv} \ensuremath{\ell_{p}} and \ensuremath{\tau_{\mathrm{flop}}} = \ensuremath{t_{p},} the conjugate field structure,
and the wave equation at propagation speed \ensuremath{c_{\mathrm{TMS}}} = \ensuremath{1/\sqrt{2}} forced by FCC
face-diagonal geometry.

Paper 2 developed the computational dynamics: the polarization field \ensuremath{\psi}
with its wave equation, the conjugate pair \ensuremath{(\psi ,} \ensuremath{\Pi_{\psi}),} the
polarization-biased resolution rule \ensuremath{P(\bar{B}} = 1 \textbar{} \ensuremath{\Pi_{\psi})} = (1 +
\ensuremath{\Pi_{\psi})/2} as the Born-rule structural correspondence, the 3-input majority
gate at every confirmation event, universal computation from \{MAJ,
constant\}, and the program formalism with stabilizer-based selection.

Paper 3 established the emergence of computation: surround-driven repair
as the [[4,2,2]] stabilizer projection at boundary sites, the
Imprinting Theorem, the Hebbian Corollary, the coherence bistability
threshold, the recursive coarse-graining operator G producing meta-level
primitives, and the first computing subsystem as the \ensuremath{(C_{D},} \ensuremath{C_{O},} \ensuremath{C_{C})}
trichotomy at meta-tetrahedral extent.

Paper 4 developed the macroscopic consequences: the recursive
coarse-graining \ensuremath{G^{(k)}} at arbitrary hierarchical depth, three growth
modes (annexation, imprinting, incorporation), self-modeling as a
structural property of every subsystem, three stressors (mortality,
predation, starvation), the Borel-Cantelli inevitability of subsystem
formation, region biography, novelty exhaustion as the monotonic decay
of \ensuremath{\nu (R,} \ensuremath{\tau ),} and the hierarchical lightspeed dispersion bounded above by
\ensuremath{c^{(k+1)}} \ensuremath{\leq} \ensuremath{(1/\sqrt{2})} \ensuremath{c^{(k)}.}

Paper 5 has closed the structural loop. The reinterpretation operator \ensuremath{\hat{R}}
acting on the exhausted biography \ensuremath{\mathcal{B}(R)} is the unique PGE-selection
forced by Axioms 3 and 4 jointly closing every other continuation. The
new substrate \ensuremath{\Xi (R)} is structurally identical to the parent substrate
under the same axioms applied at the new scale, with confirmed bits
playing the role of agents and the polarization landscape playing the
role of \ensuremath{U_{\mathrm{sh}}'.} The two-wave structure that appeared to threaten
consistency is resolved as a single wave equation at successive
recursion levels. The fundamental fields of the new substrate are the
irreducible components of the polarization field on the FCC lattice ---
scalar, vector, tensor --- with the substrate-internal QCA-to-QFT lift
of \S\,4.1 (Lemmas A and B) providing the structural bridge to standard
QFT. The initial conditions of the new substrate exhibit the observed
structural signatures of the Big Bang: maximum informational density,
near-uniform background, small fluctuations, large-scale flatness, and a
horizon-connected past --- the Big Bang Correspondence. The numerical
closures of Papers 1--4 are supplied: \ensuremath{L_{p}^{(1)}} = \ensuremath{3\sqrt{2}} \ensuremath{\ell_{p},} the \ensuremath{c^{(k)}}
saturation cascade \ensuremath{c^{(k)}} = \ensuremath{(1/\sqrt{2})^{k}} c, and the closed form of f
under additive Kolmogorov bounds. The cross-level relativistic phenomena
are derived as the substrate-level origins of variable apparent
propagation (gravitational redshift), frame-relative timing
(gravitational time dilation), and embedded-medium dilation (mass-energy
equivalence). The Standard Model gauge group SU(3) \ensuremath{\times} SU(2) \ensuremath{\times} U(1)
emerges from substrate primitives at distinct levels: SU(3) from the
triadic \ensuremath{S_{3}} symmetry of confirmation events at level 0, \ensuremath{SU(2)_{L}} from the
A/B sublattice doublet at level 1, and U(1) from the trace mode of the
bilinear role-correlation operators. Three generations follow from the
2D-minimum-closure count \ensuremath{\bar{B}} = \ensuremath{3\bar{\alpha}.} The dark sector is structurally
identified with the two surviving \ensuremath{S_{3}-\text{trivial}} vacuum residual modes, with
predicted dark-energy-to-dark-matter ratio 3:1 (observed
\textasciitilde2.65:1, 12\% gap flagged as open). The full predictions
consolidation is given in \S\,IX.

The full arc --- from the sphere primitive of Paper 1 \S\,1.2 to the
cosmological initial conditions of Paper 5 \S\,V --- is one unbroken
geometric chain. Axiom 0 supplies the maximum-entropy ground state.
Axiom 2 supplies the sphere. The Principle of Generative Economy selects
the unique minimum-sufficient generative geometry at every step. The
five axioms apply at every recursion level by axiom scale-independence.
The reinterpretation operator \ensuremath{\hat{R}} produces new substrates that inherit the
same structural content at coarser scales. The chain produces universes
inside the substrate because the geometry of confirmation, applied
recursively at the structural problem of exhaustion, has no other
available move.

The promise of the volume --- that a discrete, informationally closed
confirmation substrate built from geometric primitives is both necessary
and sufficient to generate wave dynamics, computational substrate,
hierarchical organization, and cosmological emergence --- is brought to
its conclusion here. The substrate is closed under its own dynamics. The
model has no numerical degrees of freedom left to tune. The remaining
open questions --- explicit numerical pinning of \ensuremath{L_{p}^{(k)}} and
\ensuremath{c^{(k)}} at higher levels, full representation-theoretic mapping to the
Standard Model gauge group, detailed CMB-spectrum computations, and the
agent-side framing reserved for Volume 2 --- are continuations within
the closed framework, not fundamental gaps.

The TMS proposes that physics is what happens when the simplest possible
structure that can measure itself is forced to keep measuring. Volume 1
has shown that this proposition, fully unpacked, generates a universe.

\section*{Appendix A: Detailed Derivation of the Reinterpretation Operator \ensuremath{\hat{R}}}
\addcontentsline{toc}{section}{Appendix A: Detailed Derivation of the Reinterpretation Operator \ensuremath{\hat{R}}}

\subsection*{A.1 The Constraint Set K at the Exhaustion Boundary}

We give the formal constraint set K invoked in \S\,2.2 in explicit form. At
\ensuremath{\tau} \ensuremath{\geq} \ensuremath{\tau *(R),} the substrate's continuation is constrained by:

\begin{itemize}
\item
  \ensuremath{K_{A0}} (Axiom 0): The new substrate's ground state must be
  informationally fundamental --- a maximum-entropy field with no
  preferred direction, scale, or target.
\end{itemize}

\begin{itemize}
\item
  \ensuremath{K_{A1}} (Axiom 1): The new substrate must support measurement.
\end{itemize}

\begin{itemize}
\item
  \ensuremath{K_{A2}} (Axiom 2): The new substrate's primitives must be isotropic
  immutable spheres (or their structural analogs).
\end{itemize}

\begin{itemize}
\item
  \ensuremath{K_{A3}} (Axiom 3): The transition from parent to daughter substrate must
  preserve the signals that produced the parent's confirmed bits.
\end{itemize}

\begin{itemize}
\item
  \ensuremath{K_{A4}} (Axiom 4): The transition must preserve the parent's confirmed
  bits as immutable causal record.
\end{itemize}

\begin{itemize}
\item
  \ensuremath{K_{\mathrm{Continuity}}:} The substrate must continue operating; there is no
  mechanism by which any region can halt.
\end{itemize}

\begin{itemize}
\item
  \ensuremath{K_{\mathrm{PGE},i}:} The continuation must be minimum-specification --- no
  structural commitment beyond what is forced by \ensuremath{K_{A0}} through
  \ensuremath{K_{\mathrm{Continuity}}.}
\end{itemize}

\begin{itemize}
\item
  \ensuremath{K_{\mathrm{PGE},\mathrm{ii}}:} The continuation must be generatively non-trivial --- it
  must admit downstream structure under the new substrate's dynamics.
\end{itemize}

\subsection*{A.2 Verification That \ensuremath{\hat{R}} Satisfies K}

We verify each constraint:

\begin{itemize}
\item
  \ensuremath{K_{A0}:} \ensuremath{\Xi (R)} has ground state \ensuremath{U_{\mathrm{sh}}'} = the polarization landscape across
  \ensuremath{\mathcal{B}(R).} This is a maximum-entropy field by construction --- it has no
  preferred direction (the polarization values are scalar), no preferred
  scale (the lattice structure inherited from R is regular), and no
  preferred target \ensuremath{(\Sigma (\Xi (R))} is the full configuration space at the new
  scale). \ensuremath{\checkmark}
\end{itemize}

\begin{itemize}
\item
  \ensuremath{K_{A1}:} The new substrate supports measurement by inheriting the
  triadic confirmation rule at the new scale. \ensuremath{\checkmark}
\end{itemize}

\begin{itemize}
\item
  \ensuremath{K_{A2}:} The new agents (confirmed bits of \ensuremath{\mathcal{B}(R))} are isotropic (no
  orientation parameters) and immutable (Axiom 4 of the parent). The
  structural analogy to spheres is via the T-void centroid
  identification of \S\,2.4. \ensuremath{\checkmark}
\end{itemize}

\begin{itemize}
\item
  \ensuremath{K_{A3}:} Signals that produced \ensuremath{\mathcal{B}(R)'s} confirmed bits are preserved ---
  they remain consumed and locked in the parent record, which is now the
  new substrate's pre-geometric substrate. No signal is lost. \ensuremath{\checkmark}
\end{itemize}

\begin{itemize}
\item
  \ensuremath{K_{A4}:} \ensuremath{\mathcal{B}(R)'s} confirmed bits are preserved as immutable substrate of
  \ensuremath{\Xi (R).} They cannot be overwritten by the daughter substrate's dynamics.
  \ensuremath{\checkmark}
\end{itemize}

\begin{itemize}
\item
  \ensuremath{K_{\mathrm{Continuity}}:} The transition is atomic at the regional scale, and the
  daughter substrate begins operating immediately upon reinterpretation.
  \ensuremath{\checkmark}
\end{itemize}

\begin{itemize}
\item
  \ensuremath{K_{\mathrm{PGE},i}:} \ensuremath{\hat{R}} introduces no new mechanism beyond the identification of
  confirmed bits as new agents. This is the minimum specification --- no
  other identification adds fewer structural commitments while
  satisfying \ensuremath{K_{A2}.} \ensuremath{\checkmark}
\end{itemize}

\begin{itemize}
\item
  \ensuremath{K_{\mathrm{PGE},\mathrm{ii}}:} \ensuremath{\Xi (R)} is generatively non-trivial --- it has its own
  \ensuremath{\Sigma (\Xi (R))} which is initially unsampled, so the new substrate's
  direct-read protocol can extract novelty for at least the time until
  \ensuremath{\Xi (R)} itself exhausts. \ensuremath{\checkmark}
\end{itemize}

All eight constraints are satisfied by \ensuremath{\hat{R}.} No other candidate
continuation satisfies all eight simultaneously, as established by
elimination in \S\,2.2. \ensuremath{\hat{R}} is therefore the unique element of \ensuremath{\mathcal{C}_{K}}
satisfying both PGE conditions, and the unique PGE-selection.

\subsection*{A.3 Uniqueness Up to Equivalence}

Strictly speaking, \ensuremath{\hat{R}} as defined in \S\,2.3 is unique up to permutation of
agent labels and translation of the new lattice's origin --- both
trivial equivalences that do not affect the structural content of \ensuremath{\Xi (R).}
Up to these equivalences, \ensuremath{\hat{R}} is the unique mapping satisfying K and PGE.

This uniqueness statement is the formal content of the structural claim
that the Big Bang Correspondence (\S\,V) is not a contingent fact but a
forced consequence of the framework. The daughter substrate is what it
is because \ensuremath{\hat{R}} is what it is; the initial conditions of the daughter are
forced by the parent's terminal state under the unique PGE-selection.

\section*{Appendix B: Combinatorial Enumeration of \ensuremath{L_{p}} at the Substrate Level}
\addcontentsline{toc}{section}{Appendix B: Combinatorial Enumeration of \ensuremath{L_{p}} at the Substrate Level}

\subsection*{B.1 The Minimum Stabilizer-Satisfying T-Closed Configuration}

The load-bearing derivation of \ensuremath{L_{p}^{(1)}} = \ensuremath{3\sqrt{2}} \ensuremath{\ell_{p}} is given in \S\,6.1.3,
which uses the storage convention (1 storage unit = \ensuremath{\ell_{p}/\sqrt{2},} forced by
Paper 1 \S\,VI's FCC second-moment tensor) and the stella-octangula tile
construction to produce the minimum configuration. This appendix
provides supporting combinatorial enumeration: the compressibility check
K(C) \textless{} \textbar C\textbar, elimination of smaller candidates,
and counting of distinct minimum configurations. The geometric
construction itself is given in \S\,6.1.3 and is not duplicated here.

Earlier drafts of this appendix began the construction at positions
\{(0,0,0), \ensuremath{(\ell_{p},} 0, 0), (0, \ensuremath{\ell_{p},} 0), (0, 0, \ensuremath{\ell_{p})}\} and described this as an
FCC tetrahedron. This is geometrically incorrect --- in the FCC lattice,
the nearest-neighbor distance is \ensuremath{\sqrt{2}} (not 1) in unit-edge cubic
coordinates, so those positions are cubic-lattice positions rather than
FCC nearest-neighbor positions. The \ensuremath{``L_{p}^{(1)}} = \ensuremath{4\ell_{p}''} result
obtained from that construction was a downstream artifact of the wrong
starting geometry. The current appendix has been revised to defer to
\S\,6.1.3's correct construction.

The configuration described by \S\,6.1.3 has the structural properties
relevant to this appendix: agent count 16 (one stella-octangula tile
worth of agent-sphere vertices, including the boundary needed for
stabilizer parity evaluation), one confirmed bit at the centroid, and
spatial extent \ensuremath{L_{p}^{(1)}} = \ensuremath{3\sqrt{2}} \ensuremath{\ell_{p}} \ensuremath{\approx} 4.24 \ensuremath{\ell_{p}.}

\subsection*{B.2 Compressibility (K(C) < | C|)}

The configuration is described completely by specifying: (a) the FCC
orientation of the central tile, which contributes \ensuremath{\log_2(1) = 0} bits,
because the stella-octangula tile is \ensuremath{O_{h}}-invariant --- its orientation
orbit has size 1 (see \S\,B.4); (b) the logical state of the
[[4,2,2]] code at the central tetrahedron (one of 4 states, 2 bits); (c)
the boundary configuration, which is forced by the higher-order
tetrahedral closure of \S\,2.5 once the central logical state is fixed.
The boundary is not determined by stabilizer parity alone --- the
Z-stabilizer constraint at the boundary box admits k = \ensuremath{3\cdot L_{\mathrm{box}}}
free parameters (see \S\,B.5) --- but by the T-closure dynamics, which
restrict the \ensuremath{2^{12}} stabilizer-satisfying configurations to a much smaller
set characterized by (a) and (b). The total description length is
therefore \ensuremath{\log_2(4) = 2} bits.

The information content \textbar C\textbar{} of the configuration is the
polarization values at 16 agents plus the confirmed bit, giving
\textbar C\textbar{} \ensuremath{\approx} 17 bits if each polarization is one bit. Since
K(C) = 2 \textless{} 17 = \textbar C\textbar, the configuration is
compressible: K(C) \textless{} \textbar C\textbar. The corrected
\ensuremath{K(C) = 2} bits coincides exactly with the per-confirmation richness of 2
bits established independently; the minimum program's description length
equals one confirmation's worth of form. \ensuremath{\checkmark}

\subsection*{B.3 Smaller Candidates Eliminated}

We verify that no smaller configuration satisfies all three conditions.
Consider candidates by their spatial extent:

\begin{itemize}
\item
  L = \ensuremath{1\ell_{p}:} a single agent. No tetrahedral structure; stabilizer parity
  not applicable. Eliminated.
\end{itemize}

\begin{itemize}
\item
  L = \ensuremath{2\ell_{p}:} a pair of agents. Cannot form a tetrahedron in 3D.
  Eliminated.
\end{itemize}

\begin{itemize}
\item
  L = \ensuremath{3\ell_{p}:} a triangular triple of agents. Forms a 2-simplex but not a
  3-simplex (tetrahedron); the [[4,2,2]] code requires 4
  vertices. Eliminated.
\end{itemize}

\begin{itemize}
\item
  L = \ensuremath{3\sqrt{2}} \ensuremath{\ell_{p}:} the minimum stella-octangula tile configuration of \S\,6.1.3.
  All three conditions satisfied. \ensuremath{\checkmark}
\end{itemize}

Therefore \ensuremath{L_{p}^{(1)}} = \ensuremath{3\sqrt{2}} \ensuremath{\ell_{p}} is the geometric minimum. The
construction is forced; the answer is unique up to symmetry.

\subsection*{B.4 The Combinatorial Count of Distinct Minimum Configurations}

The minimum configuration admits four distinct logical states under the
[[4,2,2]] code. Its orientation multiplicity is a single orbit: the
stella-octangula tile's eight vertices \emph{are} the eight corners of the
cube, so the tile is invariant under the full \ensuremath{O_{h}} point group and its
orbit under \ensuremath{O_{h}} has size 1. The A/B sublattice exchange (the red/blue
polarization swap) is the internal polarization degree of freedom already
counted among the four logical states.

The count is therefore 4 \ensuremath{\times} 1 = 4 distinct minimum configurations up to
translation and \ensuremath{O_{h}} symmetry: the four [[4,2,2]] logical states of the
single \ensuremath{O_{h}}-invariant stella-octangula tile. This is the substrate's
first ``program zoo'' --- the set of distinct minimum programs that can
exist at the substrate level. (The explicit two-flop closure is verified in
Appendix E: period exactly 2 in confirmed support in every tested
realization of the corpus cycle, and period 2 in the signed pattern in the
coherent-surround limit of the Paper 2 \S\,6.2 definition.)

\subsection*{B.5 The k = \ensuremath{3\cdot L_{\mathrm{box}}} Null-Space Regularity}

The dimension of the \ensuremath{F_{2}} null space of the Z-stabilizer constraint on an
FCC patch of integer extent \ensuremath{L_{\mathrm{box}}} satisfies k = \ensuremath{3\cdot L_{\mathrm{box}}} exactly for
\ensuremath{L_{\mathrm{box}}} \ensuremath{\in} \{1, 2, 3, 4\}, with k = 3, 6, 9, 12 respectively. This linear
relation expresses a structural fact about the FCC + [[4,2,2]]
combinatorics: each unit of patch extent admits exactly three additional
independent stabilizer-satisfying degrees of freedom, after accounting
for the rank-(N \ensuremath{-} \ensuremath{3\cdot L_{\mathrm{box}})} coupling among the tetrahedral parity
constraints.

The regularity is the small-L manifestation of the inclusion-exclusion
among shared-vertex tetrahedra: as the patch grows, additional
tetrahedra are added but their parity constraints become increasingly
redundant with existing ones. At \ensuremath{L_{\mathrm{box}}} \ensuremath{\in} \{1, 2, 3, 4\}, the redundancy
structure yields exactly the rank N \ensuremath{-} \ensuremath{3\cdot L_{\mathrm{box}}.} The site and tetrahedron
counts at these extents are (N, T) = (4, 1), (14, 8), (32, 27), (63,
64), with constraint ranks 1, 8, 23, 51 respectively.

The regularity persists beyond \ensuremath{L_{\mathrm{box}}} = 4: direct computation of the
\ensuremath{F_{2}} null-space rank continues to give k = \ensuremath{3\cdot L_{\mathrm{box}}} through
\ensuremath{L_{\mathrm{box}}} = 7, with ranks 93, 154, 235 yielding k = 15, 18, 21
respectively (Appendix E, Stage 12). This establishes k = \ensuremath{3\cdot L_{\mathrm{box}}}
as a stable geometric property of FCC + [[4,2,2]] across the verified
range rather than a small-L coincidence, and it belongs in Paper 1's
combinatorial inventory as such. Whether the linear relation holds for
all \ensuremath{L_{\mathrm{box}}} remains a conjecture; were it ever to break, the breakdown
point would itself be structurally informative, indicating the extent at
which higher-order constraint redundancies (plausibly tied to closed
geodesic loops on the FCC lattice) enter.

\section*{Appendix C: Closed Form of f(d, \ensuremath{K_{\psi,\mathrm{local}},} \ensuremath{K_{\psi,\mathrm{surround}},} \ensuremath{K_{\psi,\mathrm{subsystem}},} k)}
\addcontentsline{toc}{section}{Appendix C: Closed Form of f(d, \ensuremath{K_{\psi,\mathrm{local}},} \ensuremath{K_{\psi,\mathrm{surround}},} \ensuremath{K_{\psi,\mathrm{subsystem}},} k)}

\subsection*{C.1 The Additive Kolmogorov Bound}

The GUP deformation parameter \ensuremath{\beta} was introduced in Paper 1 \S\,7.2 as
scaling with the substrate footprint of a confirmed bit. The footprint
accumulates contributions from five independent sources, each identified
in Papers 1--4. Under the additive Kolmogorov bound, the total footprint
is the sum of the individual contributions.

Each source contributes to the substrate footprint a number of lattice
units proportional to its information content:

\begin{itemize}
\item
  Causal depth d contributes d lattice units (one per flop of history).
\end{itemize}

\begin{itemize}
\item
  Local algorithmic complexity \ensuremath{K_{\mathrm{local}}} contributes \ensuremath{K_{\mathrm{local}}} lattice
  units (one per bit of local algorithmic structure).
\end{itemize}

\begin{itemize}
\item
  Surround complexity \ensuremath{K_{\mathrm{surround}}} contributes \ensuremath{K_{\mathrm{surround}}} \ensuremath{\times} \ensuremath{\sqrt{2}} lattice
  units (one per bit of imprinting history, weighted by the \ensuremath{\sqrt{2}}
  imprint-depth factor of Paper 3 \S\,2.3).
\end{itemize}

\begin{itemize}
\item
  Subsystem complexity \ensuremath{K_{\mathrm{subsys}}} contributes \ensuremath{K_{\mathrm{subsys}}} \ensuremath{\times} 3 lattice units
  (one per bit of program content, weighted by the 3-component \ensuremath{(C_{D},}
  \ensuremath{C_{O},} \ensuremath{C_{C})} trichotomy of Paper 3 \S\,V).
\end{itemize}

\begin{itemize}
\item
  Hierarchical level k contributes \ensuremath{(1/\sqrt{2})^{k}} lattice units, the
  per-level meta-tetrahedral structural overhead being unity by the
  form-invariance fixed point (Paper~3 \S\,4.7).
\end{itemize}

\subsection*{C.2 The Closed Form}

Combining the five contributions:

\emph{f(d, \ensuremath{K_{\mathrm{local}},} \ensuremath{K_{\mathrm{surround}},} \ensuremath{K_{\mathrm{subsys}},} k) = 1 + \ensuremath{\alpha_{d}} \ensuremath{\cdot} d +
\ensuremath{\alpha_{\mathrm{local}}} \ensuremath{\cdot} \ensuremath{K_{\mathrm{local}}} + \ensuremath{\alpha_{\mathrm{surround}}} \ensuremath{\cdot} \ensuremath{K_{\mathrm{surround}}} + \ensuremath{\alpha_{\mathrm{subsys}}} \ensuremath{\cdot} \ensuremath{K_{\mathrm{subsys}}}
+ \ensuremath{\alpha_{k}} \ensuremath{\cdot} k}

with the coefficients identified by the geometric factors above:

\begin{equation*}
\alpha_{d} = 1, \alpha_{\mathrm{local}} = 1, \alpha_{\mathrm{surround}} = \sqrt{2}, \alpha_{\mathrm{subsys}} = 3, \alpha_{k} =
(1/\sqrt{2})^{k}
\end{equation*}

These are dimensionless geometric constants forced by the FCC lattice
connectivity and the [[4,2,2]] code distance. The leading
constant (1) ensures f \ensuremath{\geq} 1 for all values of the arguments, recovering
the bare GUP minimum length scale in the limit of vanishing arguments.

\subsection*{C.3 The \ensuremath{\beta} Formula}

The GUP deformation parameter \ensuremath{\beta} is related to f by

\begin{equation*}
\beta = \beta_{0} \cdot f(d, K_{\mathrm{local}}, K_{\mathrm{surround}}, K_{\mathrm{subsys}}, k)
\end{equation*}

where \ensuremath{\beta_{0}} is the bare deformation parameter at d = \ensuremath{K_{\mathrm{local}}} =
\ensuremath{K_{\mathrm{surround}}} = \ensuremath{K_{\mathrm{subsys}}} = 0, k = 0. In the literature on Generalized
Uncertainty Principles (Hossenfelder, 2013), \ensuremath{\beta_{0}} is typically a
dimensionless number of order unity. In the TMS framework, \ensuremath{\beta_{0}} is
forced by the FCC lattice structure at the substrate level and is
computable in principle from the per-flop noise statistics; we estimate
\ensuremath{\beta_{0}} \ensuremath{\approx} 1 from the dimensionality of the FCC lattice.

\subsection*{C.4 Beyond the Additive Bound}

The additive closed form is the leading-order expression valid when the
individual contributions are small relative to one another. When two or
more arguments become large enough that their cross-correlations
dominate, the form acquires multiplicative structure. For example, when
both d and \ensuremath{K_{\mathrm{local}}} are large, the joint contribution to f is bounded
above by

\emph{d \ensuremath{\cdot} \ensuremath{K_{\mathrm{local}}} / ln 2}

via the Kolmogorov bound K(string of length n with structure) \ensuremath{\leq} n \ensuremath{\cdot}
\ensuremath{log_{2}(\text{alphabet})} / ln 2. The full f beyond the additive limit is a
piecewise-linear function in each argument with additional
multiplicative cross-terms at large argument values.

The leading additive form is sufficient for empirical predictions at
moderate GUP corrections \ensuremath{(\beta} corrections of order \ensuremath{10^{-3}} or smaller). The
full multiplicative form is reserved for simulation work.

\section*{Appendix D: Numerical Estimates of \ensuremath{c^{(k)}} and Identification of Our Level}
\addcontentsline{toc}{section}{Appendix D: Numerical Estimates of \ensuremath{c^{(k)}} and Identification of Our Level}

\subsection*{D.1 The Saturation Cascade}

From \S\,6.2, the hierarchical lightspeed sequence at saturation is
\ensuremath{c^{(k)}} = \ensuremath{(1/\sqrt{2})^{k}} \ensuremath{\cdot} c.~In SI units, taking \ensuremath{c_{\mathrm{measured}}} =
299,792,458 m/s as our observed lightspeed, the substrate-level \ensuremath{c_{\mathrm{TMS}}}
for various candidate hierarchical depths of our universe is:

\begin{itemize}
\item
  If our universe is at k = 0 (we are the substrate level): \ensuremath{c_{\mathrm{TMS}}} =
  \ensuremath{c_{\mathrm{measured}}} \ensuremath{\approx} 3 \ensuremath{\times} \ensuremath{10^{8}} m/s. This is the minimum possible value.
\end{itemize}

\begin{itemize}
\item
  If our universe is at k = 1: \ensuremath{c_{\mathrm{TMS}}} = \ensuremath{\sqrt{2}} \ensuremath{\cdot} \ensuremath{c_{\mathrm{measured}}} \ensuremath{\approx} 4.24 \ensuremath{\times} \ensuremath{10^{8}}
  m/s.
\end{itemize}

\begin{itemize}
\item
  If our universe is at k = 2: \ensuremath{c_{\mathrm{TMS}}} = 2 \ensuremath{\cdot} \ensuremath{c_{\mathrm{measured}}} \ensuremath{\approx} 6 \ensuremath{\times} \ensuremath{10^{8}} m/s.
\end{itemize}

\begin{itemize}
\item
  If our universe is at k = 3: \ensuremath{c_{\mathrm{TMS}}} = \ensuremath{2\sqrt{2}} \ensuremath{\cdot} \ensuremath{c_{\mathrm{measured}}} \ensuremath{\approx} 8.48 \ensuremath{\times} \ensuremath{10^{8}}
  m/s.
\end{itemize}

\begin{itemize}
\item
  If our universe is at k = 4: \ensuremath{c_{\mathrm{TMS}}} = 4 \ensuremath{\cdot} \ensuremath{c_{\mathrm{measured}}} \ensuremath{\approx} 12 \ensuremath{\times} \ensuremath{10^{8}} m/s.
\end{itemize}

The k = 3 case is particularly intriguing if the SU(3) \ensuremath{\times} SU(2) \ensuremath{\times} U(1)
gauge structure of \S\,8.3 corresponds to three hierarchical levels above
substrate. Under this identification, our universe is at hierarchical
depth k = 3, and the substrate-level \ensuremath{c_{\mathrm{TMS}}} is \ensuremath{2\sqrt{2}} \ensuremath{\approx} 2.83 times our
observed lightspeed.

\subsection*{D.2 Consistency with Cosmological Observations}

The TMS framework predicts specific structural features of cosmological
observations under different hierarchical-depth assignments. Each
prediction provides a partial constraint:

\begin{itemize}
\item
  CMB anisotropy spectrum: the position of the first acoustic peak at \ensuremath{\ell}
  \ensuremath{\approx} 220 corresponds to a characteristic angular scale of
  \textasciitilde\ensuremath{1^{\circ},} which under standard cosmology is interpreted as
  the sound horizon at recombination. In the TMS reading, this scale
  corresponds to the typical late-phase configuration scale of \ensuremath{\mathcal{B}(R),} set
  by the parent-level program structure.
\end{itemize}

\begin{itemize}
\item
  Large-scale flatness: the observed near-flatness \ensuremath{(\Omega} \ensuremath{\approx} 1 to one part in
  100) is consistent with the structural flatness derived in \S\,5.4 for
  any hierarchical depth, since flatness is a structural property of the
  reinterpretation event rather than a numerical fine-tuning.
\end{itemize}

\begin{itemize}
\item
  Horizon problem: the structural horizon-problem resolution of \S\,5.3
  applies at any hierarchical depth.
\end{itemize}

None of these constraints uniquely pins our universe's k. The structural
framework is compatible with multiple hierarchical-depth assignments;
pinning requires the combinatorial enumeration of \ensuremath{L_{p}^{(k)}} and the
comparison to cosmological observables at sufficient precision.

\subsection*{D.3 The Newton Constant and Planck Scale}

Under the hierarchical-depth assignment, the relationship between
substrate-level and observed Planck length is:

\begin{equation*}
\ell_{p}^{(\text{observed})} = (\sqrt{2})^{k} \cdot \ell_{p}^{(\text{substrate})}
\end{equation*}

at saturation. For k = 3, the observed Planck length is \ensuremath{2\sqrt{2}} \ensuremath{\approx} 2.83 times
the substrate-level Planck length. This implies that the substrate-level
Planck length is \ensuremath{\ell_{p}^{(\text{substrate})}} \ensuremath{\approx} 5.7 \ensuremath{\times} \ensuremath{10^{-36}} m, which is the same
order as the observed Planck length and structurally compatible with the
identification \ensuremath{\ell_{0}} \ensuremath{\equiv} \ensuremath{\ell_{p}} of Paper 1 \S\,3.2.

The Newton constant G satisfies a similar level-scaling relation under
the cross-level relativistic phenomena of \S\,VII: G corresponds in the TMS
picture to the coupling between hierarchical levels, with the observed
value reflecting the level-coupling at our hierarchical depth. The
detailed identification requires the full continuum-limit
Einstein-equation derivation flagged in \S\,7.5.

\section*{Appendix E: Empirical Verification of the Sublattice Bifurcation Structure}
\addcontentsline{toc}{section}{Appendix E: Empirical Verification of the Sublattice Bifurcation Structure}

The sublattice bifurcation of \S\,6.1 and the bifurcation-derived maximum
novelty output theorem of \S\,2.6, the matter-antimatter structural
correspondence of \S\,5.3, and the cross-level commutator mechanism of \S\,8.2
all depend on the level-1 substrate having a specific geometric
structure: two interpenetrating FCC sublattices A and B on the
all-odd-coord integer lattice (in the storage convention with \ensuremath{\times 2} scaling
per hierarchical level), with each sublattice independently supporting
[[4,2,2]] stabilizer dynamics, and the two sublattices related
by spatial inversion through the level-2 centroids.

This appendix documents the empirical verification of this structure
through two simulation programs (Stages 5 and 6) that test the
bifurcation hypothesis at level 1 and at level 2, respectively. The
simulations are documented in the supplementary files \ensuremath{lp1_{stage5}.py} and
\ensuremath{lp1_{stage6}.py} (available as ancillary files with this preprint).

\subsection*{E.1 Lattice Convention}

The simulations work in the \emph{storage convention} used throughout
\S\,6.1. The level-0 agent FCC consists of integer-coordinate sites with
sum i + j + k even. Each hierarchical level scales by a factor of 2
relative to the previous level's lattice spacing, so the level-1 sites
lie at all-odd-coord integer positions and the level-2 sites lie at
all-even-coord integer positions. In substrate units, one storage unit
equals \ensuremath{\ell_{p}/\sqrt{2},} since the level-0 nearest-neighbor distance is \ensuremath{\sqrt{2}} storage
units and \ensuremath{\ell_{p}} is fixed to that nearest-neighbor distance by Paper 1's
wave-equation derivation (second-moment tensor \ensuremath{4\ell^{2}_{p}\delta_{\mu}\nu} over 12 nearest
neighbors).

\subsection*{E.2 Stage 5: Per-Sublattice FCC Verification at Level 1}

\textbf{Hypothesis.} The all-odd-coord level-1 sites split by sum-mod-4
into two interpenetrating FCC sublattices A (sum \ensuremath{\equiv} 1 mod 4) and B (sum \ensuremath{\equiv}
3 mod 4), each with FCC nearest-neighbor distance \ensuremath{2\sqrt{2}} storage units
(face-diagonal of the underlying cubic cell of side 2 storage units).
Each sublattice independently supports tetrahedral closure (four-site
clusters at mutual NN distance squared = 8) and an associated
[[4,2,2]] stabilizer structure with null-space scaling k = 3n in
the level-local depth n.

\textbf{Method.} For each \ensuremath{L_{\mathrm{box}}} \ensuremath{\in} \{3, 5, 7, 9, 11\}, enumerate the
sublattice-A sites (all-odd-coord integers in [1, \ensuremath{L_{\mathrm{box}}]^{3}} with sum
\ensuremath{\equiv} 1 mod 4), count nearest neighbors at squared distance 8, identify all
four-site clusters with mutual NN at this distance (the tetrahedra), and
compute the \ensuremath{F_{2}} null-space dimension of the parity matrix relating sites
to tetrahedra. Repeat for sublattice B (sum \ensuremath{\equiv} 3 mod 4). For shared
centroids, verify that each A-tetrahedron's centroid coincides with a
B-tetrahedron's centroid and that the two tetrahedra are exact inversion
partners through the shared centroid.

\textbf{Results.} Stage 5 verifies the bifurcation hypothesis across all
tested \ensuremath{L_{\mathrm{box}}.} The per-sublattice structural numbers:

\begin{longtable}[]{@{}
  >{\raggedright\arraybackslash}p{(\columnwidth - 16\tabcolsep) * \real{0.1108}}
  >{\raggedright\arraybackslash}p{(\columnwidth - 16\tabcolsep) * \real{0.1111}}
  >{\raggedright\arraybackslash}p{(\columnwidth - 16\tabcolsep) * \real{0.1113}}
  >{\raggedright\arraybackslash}p{(\columnwidth - 16\tabcolsep) * \real{0.1111}}
  >{\raggedright\arraybackslash}p{(\columnwidth - 16\tabcolsep) * \real{0.1112}}
  >{\raggedright\arraybackslash}p{(\columnwidth - 16\tabcolsep) * \real{0.1111}}
  >{\raggedright\arraybackslash}p{(\columnwidth - 16\tabcolsep) * \real{0.1113}}
  >{\raggedright\arraybackslash}p{(\columnwidth - 16\tabcolsep) * \real{0.1109}}
  >{\raggedright\arraybackslash}p{(\columnwidth - 16\tabcolsep) * \real{0.1110}}@{}}
\toprule\noalign{}
\begin{minipage}[b]{\linewidth}\raggedright
\ensuremath{L_{\mathrm{box}}}
\end{minipage} & \begin{minipage}[b]{\linewidth}\raggedright
A: sites
\end{minipage} & \begin{minipage}[b]{\linewidth}\raggedright
A: tets
\end{minipage} & \begin{minipage}[b]{\linewidth}\raggedright
A: k
\end{minipage} & \begin{minipage}[b]{\linewidth}\raggedright
A: max NN
\end{minipage} & \begin{minipage}[b]{\linewidth}\raggedright
B: sites
\end{minipage} & \begin{minipage}[b]{\linewidth}\raggedright
B: tets
\end{minipage} & \begin{minipage}[b]{\linewidth}\raggedright
B: k
\end{minipage} & \begin{minipage}[b]{\linewidth}\raggedright
B: max NN
\end{minipage} \\
\midrule\noalign{}
\endhead
\bottomrule\noalign{}
\endlastfoot
3 & 4 & 1 & 3 & 3 & 4 & 1 & 3 & 3 \\
5 & 13 & 8 & 6 & 12 & 14 & 8 & 6 & 8 \\
7 & 32 & 27 & 9 & 12 & 32 & 27 & 9 & 12 \\
9 & 62 & 64 & 12 & 12 & 63 & 64 & 12 & 12 \\
11 & 108 & 125 & 15 & 12 & 108 & 125 & 15 & 12 \\
\end{longtable}

Three structural features are confirmed:

\begin{itemize}
\item
  \textbf{FCC coordination 12 at interior sites.} Each sublattice
  achieves the standard FCC coordination number of 12 at interior sites
  for \ensuremath{L_{\mathrm{box}}} \ensuremath{\geq} 7, confirming that each sublattice is a genuine FCC
  lattice at NN distance \ensuremath{2\sqrt{2}} storage units. The low max-NN values at
  small \ensuremath{L_{\mathrm{box}}} reflect boundary effects (no sites are fully interior in
  the small box).
\end{itemize}

\begin{itemize}
\item
  \textbf{\ensuremath{n^{3}} tetrahedral count.} The tetrahedron count scales as \ensuremath{n^{3}}
  where n = \ensuremath{(L_{\mathrm{box}}} \ensuremath{-} 1)/2 is the level-local depth: \ensuremath{n=1\to 1,} \ensuremath{n=2\to 8,}
  \ensuremath{n=3\to 27,} \ensuremath{n=4\to 64,} \ensuremath{n=5\to 125.} This is the standard FCC cubic-volume scaling
  of regular tetrahedral cells.
\end{itemize}

\begin{itemize}
\item
  \textbf{k = 3n null-space scaling.} The stabilizer-satisfying
  configuration space dimension over \ensuremath{F_{2}} scales linearly with depth: k =
  3n through n = 5. This is the same scale-invariant k = 3n law that
  governs the substrate level (Paper 1 \S\,5.4, scale-invariant stabilizer
  null-space theorem), now empirically confirmed at level 1 on each
  sublattice independently.
\end{itemize}

\textbf{Inversion-partnership verification.} For each \ensuremath{L_{\mathrm{box}}} \ensuremath{\in} \{5, 7,
9\}, every A-tetrahedron centroid coincides with a B-tetrahedron
centroid, and (sampled at 3 centroids per \ensuremath{L_{\mathrm{box}})} the A-tetrahedron and
B-tetrahedron are exact inversion partners through their shared
centroid:

\begin{longtable}[]{@{}
  >{\raggedright\arraybackslash}p{(\columnwidth - 12\tabcolsep) * \real{0.1429}}
  >{\raggedright\arraybackslash}p{(\columnwidth - 12\tabcolsep) * \real{0.1429}}
  >{\raggedright\arraybackslash}p{(\columnwidth - 12\tabcolsep) * \real{0.1429}}
  >{\raggedright\arraybackslash}p{(\columnwidth - 12\tabcolsep) * \real{0.1429}}
  >{\raggedright\arraybackslash}p{(\columnwidth - 12\tabcolsep) * \real{0.1429}}
  >{\raggedright\arraybackslash}p{(\columnwidth - 12\tabcolsep) * \real{0.1429}}
  >{\raggedright\arraybackslash}p{(\columnwidth - 12\tabcolsep) * \real{0.1429}}@{}}
\toprule\noalign{}
\begin{minipage}[b]{\linewidth}\raggedright
\ensuremath{L_{\mathrm{box}}}
\end{minipage} & \begin{minipage}[b]{\linewidth}\raggedright
A centroids
\end{minipage} & \begin{minipage}[b]{\linewidth}\raggedright
B centroids
\end{minipage} & \begin{minipage}[b]{\linewidth}\raggedright
Shared
\end{minipage} & \begin{minipage}[b]{\linewidth}\raggedright
A-only
\end{minipage} & \begin{minipage}[b]{\linewidth}\raggedright
B-only
\end{minipage} & \begin{minipage}[b]{\linewidth}\raggedright
Sample inversion verified
\end{minipage} \\
\midrule\noalign{}
\endhead
\bottomrule\noalign{}
\endlastfoot
5 & 8 & 8 & 8 & 0 & 0 & 3/3 \ensuremath{\checkmark} \\
7 & 27 & 27 & 27 & 0 & 0 & 3/3 \ensuremath{\checkmark} \\
9 & 64 & 64 & 64 & 0 & 0 & 3/3 \ensuremath{\checkmark} \\
\end{longtable}

100\% of A-tetrahedron centroids coincide with B-tetrahedron centroids;
100\% of sampled pairs are exact inversion partners. The structural
claim of \S\,5.3 Claim 2 --- that every level-1 spatial site hosts an
inversion-paired (A-tet, B-tet) pair through the level-2 centroid --- is
verified empirically.

\subsection*{E.3 Stage 6: Level-1 vs Level-2 Structural Identity}

\textbf{Hypothesis.} The bifurcation pattern repeats at level 2 with
identical structural numbers. Specifically, the level-2 sites
(all-even-coord integers in the storage frame) split by sum mod 4 \ensuremath{\in} \{0,
2\} into two interpenetrating FCC sublattices A' (sum \ensuremath{\equiv} 0 mod 4) and B'
(sum \ensuremath{\equiv} 2 mod 4), each independently supporting tetrahedral closure with
the same \ensuremath{n^{3}} count and the same k = 3n null-space scaling as level 1.

\textbf{Method.} Repeat Stage 5's per-sublattice and
inversion-partnership analyses, this time at level 2 (all-even-coord
integers in [0, \ensuremath{L_{\mathrm{box}}]^{3}} with the appropriate sum-mod-4 filter for
sublattices A' and B'). For comparison, rerun the level-1 analyses at
matched \ensuremath{L_{\mathrm{box}}} values and compare the structural numbers side-by-side.

\textbf{Results.} At every matched level-local depth n, the level-1 and
level-2 structural numbers are identical:

\begin{longtable}[]{@{}
  >{\raggedright\arraybackslash}p{(\columnwidth - 12\tabcolsep) * \real{0.1429}}
  >{\raggedright\arraybackslash}p{(\columnwidth - 12\tabcolsep) * \real{0.1429}}
  >{\raggedright\arraybackslash}p{(\columnwidth - 12\tabcolsep) * \real{0.1429}}
  >{\raggedright\arraybackslash}p{(\columnwidth - 12\tabcolsep) * \real{0.1429}}
  >{\raggedright\arraybackslash}p{(\columnwidth - 12\tabcolsep) * \real{0.1429}}
  >{\raggedright\arraybackslash}p{(\columnwidth - 12\tabcolsep) * \real{0.1429}}
  >{\raggedright\arraybackslash}p{(\columnwidth - 12\tabcolsep) * \real{0.1429}}@{}}
\toprule\noalign{}
\begin{minipage}[b]{\linewidth}\raggedright
n
\end{minipage} & \begin{minipage}[b]{\linewidth}\raggedright
Level-1: tets
\end{minipage} & \begin{minipage}[b]{\linewidth}\raggedright
Level-2: tets
\end{minipage} & \begin{minipage}[b]{\linewidth}\raggedright
Match
\end{minipage} & \begin{minipage}[b]{\linewidth}\raggedright
Level-1: k
\end{minipage} & \begin{minipage}[b]{\linewidth}\raggedright
Level-2: k
\end{minipage} & \begin{minipage}[b]{\linewidth}\raggedright
Match
\end{minipage} \\
\midrule\noalign{}
\endhead
\bottomrule\noalign{}
\endlastfoot
2 & 8 & 8 & \ensuremath{\checkmark} & 6 & 6 & \ensuremath{\checkmark} \\
3 & 27 & 27 & \ensuremath{\checkmark} & 9 & 9 & \ensuremath{\checkmark} \\
4 & 64 & 64 & \ensuremath{\checkmark} & 12 & 12 & \ensuremath{\checkmark} \\
5 & 125 & 125 & \ensuremath{\checkmark} & 15 & 15 & \ensuremath{\checkmark} \\
\end{longtable}

The bifurcation pattern is \textbf{structurally identical} at level 1
and level 2. This confirms the constant bifurcation factor of 2 per
level (not a compounding factor \ensuremath{2^{k}),} and provides empirical content
to the Recursion-Preservation Theorem of \S\,2.5: the level-1 substrate and
the level-2 substrate have the same geometric structure at the same
level-local depth.

\textbf{Inversion-partnership at level 2.} At every \ensuremath{L_{\mathrm{box}}} tested at
level 2, 100\% of A'-tetrahedron centroids coincide with B'-tetrahedron
centroids, and sampled pairs are exact inversion partners:

\begin{longtable}[]{@{}
  >{\raggedright\arraybackslash}p{(\columnwidth - 8\tabcolsep) * \real{0.2000}}
  >{\raggedright\arraybackslash}p{(\columnwidth - 8\tabcolsep) * \real{0.2000}}
  >{\raggedright\arraybackslash}p{(\columnwidth - 8\tabcolsep) * \real{0.2000}}
  >{\raggedright\arraybackslash}p{(\columnwidth - 8\tabcolsep) * \real{0.2000}}
  >{\raggedright\arraybackslash}p{(\columnwidth - 8\tabcolsep) * \real{0.2000}}@{}}
\toprule\noalign{}
\begin{minipage}[b]{\linewidth}\raggedright
\ensuremath{L_{\mathrm{box}}}
\end{minipage} & \begin{minipage}[b]{\linewidth}\raggedright
A' centroids
\end{minipage} & \begin{minipage}[b]{\linewidth}\raggedright
B' centroids
\end{minipage} & \begin{minipage}[b]{\linewidth}\raggedright
Shared
\end{minipage} & \begin{minipage}[b]{\linewidth}\raggedright
Sample inversion verified
\end{minipage} \\
\midrule\noalign{}
\endhead
\bottomrule\noalign{}
\endlastfoot
6 & 27 & 27 & 27 & 5/5 \ensuremath{\checkmark} \\
8 & 64 & 64 & 64 & 5/5 \ensuremath{\checkmark} \\
10 & 125 & 125 & 125 & 5/5 \ensuremath{\checkmark} \\
12 & 216 & 216 & 216 & 5/5 \ensuremath{\checkmark} \\
\end{longtable}

The level-2 bifurcation is empirically as exact as the level-1
bifurcation.

\subsection*{E.4 Summary of Empirical Content}

The simulation evidence supports the following structural claims of the
main text:

\begin{itemize}
\item
  \textbf{\S\,6.1.3} \ensuremath{(L^{p(1)}} = \ensuremath{3\sqrt{2}} \ensuremath{\ell_{p}):} the bifurcated structure on which
  the derivation rests is geometrically realized at level 1 (Stage 5).
\end{itemize}

\begin{itemize}
\item
  \textbf{\S\,2.6} (Maximum Novelty Output Theorem): the bifurcation
  doubles the agent count per level (one A-flavor + one B-flavor agent
  per shared centroid), empirically verified.
\end{itemize}

\begin{itemize}
\item
  \textbf{\S\,5.3 Claim 1} (two sublattices at every level k \ensuremath{\geq} 1): verified
  at level 1 (Stage 5) and level 2 (Stage 6).
\end{itemize}

\begin{itemize}
\item
  \textbf{\S\,5.3 Claim 2} (inversion-partner pairing): 100\% verified at
  both levels.
\end{itemize}

\begin{itemize}
\item
  \textbf{Scale-invariant k = 3n null-space theorem} (Paper 1 \S\,5.4):
  empirically verified through n = 5 at level 1 and through n = 5 at
  level 2 on matched depth.
\end{itemize}

\begin{itemize}
\item
  \textbf{\S\,8.2} (cross-level commutator mechanism): the bifurcation
  structure on which \S\,8.2 builds is empirically the actual structure of
  the level-1 substrate.
\end{itemize}

\subsection*{E.5 Code Availability}

The Python source for the simulations is available as ancillary files
with this preprint: \ensuremath{lp1_{stage5}.py} (Stage 5, per-sublattice FCC
verification at level 1) and \ensuremath{lp1_{stage6}.py} (Stage 6, level-1 vs level-2
structural identity). The raw output logs are similarly available as
\ensuremath{lp1_{stage5,\mathrm{results}}.\text{txt}} and \ensuremath{lp1_{stage6,\mathrm{results}}.\text{txt}.} The simulations
are deterministic and reproducible; reviewers and readers can verify the
structural numbers tabulated above by direct execution. The full stage
suite backing the remaining structural results of this volume (\S\,E.6) is
provided on the same basis --- each script standalone under Python 3 with
NumPy and SciPy, each accompanied by its captured output log.

\subsection*{E.6 The Full Verification Suite (Reproducibility)}

The structural results introduced in this volume are each backed by a
standalone, deterministic script with a captured output log, provided as
ancillary files. Each recomputes a structural quantity from the substrate
rules and checks it against the value quoted in the text; none fits a
parameter to a target.

\emph{Geometry and code structure.} \ensuremath{lp1_{stage5}}--\ensuremath{lp1_{stage12}} verify the
\ensuremath{L_p} bifurcation numbers, the \ensuremath{[[4,2,2]]} classical and quantum distance,
the stella-octangula tile with \ensuremath{L_p = 3\sqrt{2}\,\ell_p}, and the
\ensuremath{k = 3\cdot L_{\mathrm{box}}} null-space regularity, which Stage 12 extends through
\ensuremath{L_{\mathrm{box}} = 7} (ranks \ensuremath{93, 154, 235} giving \ensuremath{k = 15, 18, 21}; \S\,B.5).
\ensuremath{lp1_{stage13}} computes the minimum-configuration count and returns \ensuremath{4}: the
stella tile is \ensuremath{O_{h}}-invariant (orbit size \ensuremath{1}), so the only multiplicity is
the four \ensuremath{[[4,2,2]]} logical states (\S\,B.4). \ensuremath{tms_{dynamics}.py} defines the
operational flop operator \ensuremath{T} imported by the later stages;
\ensuremath{spec_{verify}\_T.py} checks the T specification against the corpus's own
arithmetic (the four \ensuremath{\Pi_{\psi}} values and both resolution branches; the
\ensuremath{1\to 4} transport signal accounting). The two-flop T-closure of the minimum
program is verified directly: the corpus-described cycle (tile \ensuremath{\to} level-2
centroid \ensuremath{\to} tile, tetrahedral adjacency) closes with period exactly 2 in
confirmed support in every tested variant and branch, with the singular
centroid geometrically forced, and closes with period 2 in the full signed
polarization pattern in the coherent-surround limit --- the regime named by
the Paper 2 \S\,6.2 program definition --- with finite-coherence runs
deviating and recovering exactly as the probability-approaching-1 clause
states (\ensuremath{claim5_{corrected}.py}; \ensuremath{lp1_{stage11}} retains the
within-sublattice realization as a negative control, which correctly fails
to cycle because its domain omits the centroid).

\emph{Matter sector.} \ensuremath{lp1_{stage18}} verifies the spinor result: the
resolution rule \ensuremath{P(+) = (1+\Pi)/2} equals \ensuremath{\cos^2(\theta/2)} to machine
precision across the sampled angles (agents are SU(2) spin-\ensuremath{\tfrac{1}{2}}), and the
product of three \ensuremath{2\pi} rotations is \ensuremath{-I} (the fermion sign is the oddness of
three; \S\,4.1).

\emph{Closed constants.} \ensuremath{lp1_{stage19}} verifies the area-coupling factor of
two by the exact up/down triangular-face bipartition of the horizon
manifold, and \ensuremath{\varepsilon^{(k)} = 0} by the integer-coordinate fit of the regular
tetrahedral closure at every scale (\S\,5.5).

\emph{Mass and chemistry.} \ensuremath{lp1_{stage20}} verifies the minimum standing
program (size four, a bent cluster) and the discrete rational amplitude
set \ensuremath{|\Pi| \in \{1, \tfrac{1}{2}, \tfrac{1}{3}, 0\}} (\S\,8.4.3). \ensuremath{lp1_{stage17}} verifies the
subshell capacities \ensuremath{(2, 6, 10)} from the \ensuremath{SO(3)\downarrow O_{h}} subduction, the
Madelung \ensuremath{n+\ell} order, and the noble-gas atomic numbers
\ensuremath{(2, 10, 18, 36, 54, 86)} by direct counting.

\emph{Spin geometry, dynamics, form-invariance, heavy elements.}
\ensuremath{lp1_{stage21}} verifies the spin-geometry reconstruction (\S\,4.1): angles
defined from recoupling overlaps embed in exactly three dimensions
(positive-semidefinite, rank three, to \ensuremath{N = 200}, with the spin-1 and
non-inner-product controls failing as required); three is the minimal valence
carrying a fermionic top representation and a single orientation; and the
trivalent node is a faithful SU(2) representation with the \ensuremath{2\pi \to -I} double
cover. \ensuremath{lp1_{stage22}} verifies that the diagonal whole-event action is forced by
the flop --- the \ensuremath{S_{3}}-symmetric resolution rule admits only the diagonal
action, and the real probabilistic flop is covariant with no preferred axis (the
covariance holds to machine precision while the mismatched-axis control differs by
order one). \ensuremath{lp1_{stage23}} verifies the coarse-graining fixed point (\S\,4.7):
the meta-lattice is FCC with coordination 12 and 24 triangles and 8 tetrahedra per
site, identical across levels 0 through 3. \ensuremath{lp1_{stage24}} verifies the full
periodic table (\S\,8.4.7): the capacity \ensuremath{2(2\ell+1)} from the harmonic
angular-mode count (giving \ensuremath{f = 14}), and the noble-gas atomic numbers
\ensuremath{(2, 10, 18, 36, 54, 86, 118)} with period lengths \ensuremath{(2, 8, 8, 18, 18, 32, 32)} and
the \ensuremath{4f}/\ensuremath{5f} blocks by counting.

\emph{Dispersion, gauge closure, and the selection functional.}
\ensuremath{dispersion_{25}} computes the FCC dispersion relation of \S\,4.1 in closed form:
the exact small-momentum expansion, the isotropic quadratic term, the anisotropic
quartic remainder with coefficient \ensuremath{1/48}, the [100] \ensuremath{<} [110] \ensuremath{<} [111] angular
signature, and the \ensuremath{(E/E_{\mathrm{Planck}})^{2}} magnitude against experimental
Lorentz-violation bounds. \ensuremath{su3_{closure}} verifies the su(3) closure of Appendix H:
eight independent generators, all commutators closing to machine precision,
trivial center, negative-definite Killing form (eigenvalues all \ensuremath{-12}), rank two,
irreducible fundamental. \ensuremath{pge_{moduli}} verifies the Generative-Economy selection
functional: the moduli-dimension (specification cost) of the sphere and the
alternative primitives, the integer triadic-enclosure count, and the derivation
of the FCC-selecting constraint set from Axiom~4 on the contested HCP case.

Every script is deterministic and standalone; the quoted structural
numbers can be reproduced by direct execution, and the captured logs are
provided alongside the source.

\section*{Appendix F: Cross-Level Commutator Structure (Slim Sketch)}
\addcontentsline{toc}{section}{Appendix F: Cross-Level Commutator Structure (Slim Sketch)}

This appendix sketches the two-sublattice cross-level commutator
structure of \S\,8.2 in slightly more detail, as a target for subsequent
representation-theoretic work. Specifically: the A-stabilizer \ensuremath{S_{x}^{A}}
and the B-stabilizer \ensuremath{S_{x}^{B}} at level k can be defined as operators on
the level-1 polarization landscape that act independently on their
respective sublattices. Their product \ensuremath{S_{x}^{A}} \ensuremath{\cdot} \ensuremath{S_{x}^{B}} acts at the
level-(k+1) shared centroid as a combined operator whose value depends
on the order of action when A and B are not in a symmetric
configuration. The non-vanishing commutator \ensuremath{[S_{x}^{A},} \ensuremath{S_{x}^{B}]} at
the shared centroid is the substrate-level origin of the non-Abelian
gauge closure of SU(2) at level k + 1.

The full continuum-limit derivation of this commutator structure as the
SU(2) algebra of electroweak theory requires the QCA-to-QFT apparatus,
with the specific identification depending on detailed analysis of the
polarization-field representation theory under the bifurcated FCC +
[[4,2,2]] symmetry. This is a target for subsequent work in
Volume 1 (or Volume 2) and is sketched here as the structural direction
in which the representation-theoretic completion should proceed.

\section*{Appendix G: Norm-Preservation of the Transport (the Anti-Hermiticity Plank of \S\,8.3)}
\addcontentsline{toc}{section}{Appendix G: Norm-Preservation of the Transport (the Anti-Hermiticity Plank of \S\,8.3)}

The su(3) closure of \S\,8.3 rests on the transport generators being
anti-Hermitian and traceless: anti-Hermiticity selects the compact real
form and bounds the closure to u(3)/su(3) rather than the non-compact
\ensuremath{sl(3,\mathbb{C}),} and tracelessness removes the U(1) trace mode. Anti-Hermiticity
is equivalent to the statement that the transport preserves the norm of
the propagated amplitude --- that the re-confirmation dynamics neither
amplify nor damp the internal carry. This appendix establishes that
norm-preservation directly from the substrate primitives, in three
independent checks. The mean-field map M whose spectral radius is
computed here is the FCC nearest-neighbour re-confirmation averaging of
Paper 3 \S\,IV; the resolution step is the confirmation primitive P = (1 +
\ensuremath{\Pi )/2} of Paper 2 \S\,IV.

\textbf{G.1 The resolution term is zero-mean.} The confirmation rule
resolves a bias \ensuremath{\Pi} \ensuremath{\in} \ensuremath{[-1,} 1] into a bit \ensuremath{\chi} \ensuremath{\in} \{\ensuremath{-1,} +1\} with \ensuremath{P(\chi} = +1)
= (1 + \ensuremath{\Pi )/2.} The conditional expectation is then \ensuremath{E[\chi} \textbar{} \ensuremath{\Pi ]} =
(+1)(1 + \ensuremath{\Pi )/2} + \ensuremath{(-1)(1} \ensuremath{-} \ensuremath{\Pi )/2} = \ensuremath{\Pi ,} exactly and for every \ensuremath{\Pi .} The
resolution therefore injects no mean bias: the stochastic part of
re-confirmation is a martingale increment, contributing variance but not
drift. A Monte-Carlo check across twenty-one biases spanning \ensuremath{[-1,} 1]
at four million samples each gives a worst-case deviation \textbar\ensuremath{〉\chi 〉}
\ensuremath{-} \ensuremath{\Pi}\textbar{} of 8 \ensuremath{\times} \ensuremath{10^{-4},} consistent with the analytic identity. This
closes the term previously flagged in the volume as ``asserted, not
proved'': it is provable, and the proof is one line.

\textbf{G.2 The mean-field re-confirmation has spectral radius exactly
one, saturated only by the gauge mode.} The re-confirmation map averages
a site's value over its twelve FCC nearest neighbours (the face-diagonal
directions of Paper 1 \S\,VI). Because the FCC lattice is
inversion-symmetric, the Fourier symbol of this averaging operator is
real: \ensuremath{\lambda (k)} = (1/12) \ensuremath{\Sigma_{d}} cos(k \ensuremath{\cdot} d), summed over the twelve neighbour
vectors d. At the zone centre \ensuremath{\lambda (0)} = 1 exactly; a scan over the
Brillouin zone gives a maximum modulus of 1 attained only at the \ensuremath{\Gamma}
point, with \textbar \ensuremath{\lambda (k)}\textbar{} \ensuremath{\leq} 0.996 at every grid point away
from it and a representative finite-k value \ensuremath{\lambda (0.4,} 0, 0) \ensuremath{\approx} 0.97. The
spectral radius is therefore exactly 1, never greater. The unit
eigenvalue is the uniform (k = 0) internal mode --- the gauge carry that
the transport is meant to preserve --- while all finite-k spatial modes
are strictly contracted: this is ordinary lattice dispersion, isotropic
at leading order by the isotropic FCC second-moment tensor (Paper 1
\S\,6.5), and is not the gauge direction.

\textbf{G.3 The full stochastic dynamics never amplify the signal.}
Combining the two steps --- resolve each site, then re-confirm by
neighbour averaging with reinforcement weight r --- and sweeping r from
0 to 0.95, the growth rate of the mean field is 1 at the \ensuremath{\Gamma} mode
(measured 1.000 \ensuremath{\pm} 0.002 across the sweep) and strictly below 1 at every
finite wavevector tested. The effective spectral radius does not exceed
1. One methodological caution is recorded because it was a live trap:
measuring the growth of the total field norm from a near-zero initial
condition gives a spurious rate of about 1.21, because the zero-mean
resolution noise (G.1) inflates the field variance without carrying any
net signal. That figure is noise power, not a spectral radius; the
correct object is the growth rate of the mean field with the noise
ensemble-averaged out, which is what the \ensuremath{\Gamma -\text{mode}} measurement reports.
With that correction the three checks agree.

The three checks together establish that the re-confirmation transport
is norm-preserving, with the single conserved mode being the uniform
internal carry. A norm-preserving transport has unitary propagation and
hence anti-Hermitian generators; together with tracelessness this
confines the generated algebra to su(3), with no overshoot to the
non-compact \ensuremath{sl(3,\mathbb{C})} or the larger sp(6) --- the plank on which \S\,8.3
stands. \textbf{Status: Derived.} The reproducibility script
\ensuremath{(normpres_{G})} accompanies the companion materials.

\section*{Appendix H: Explicit su(3) Closure of the Transport Generators}
\addcontentsline{toc}{section}{Appendix H: Explicit su(3) Closure of the Transport Generators}

Section 8.3 derives the colour algebra as the closure of the eight transport
generators built from substrate primitives: the three real antisymmetric
matrices from value transport (the so(3) rotations of the agent labels) and the
five \ensuremath{i}-times-symmetric-traceless matrices from the momentum carry. This
appendix records the explicit computation confirming that these eight generators
close to exactly su(3), so that the claim rests on a displayed calculation rather
than an assertion. All quantities are reproducible by direct evaluation
\ensuremath{(su3_{closure})} in the companion materials.

\textbf{H.1 The generators.} Write the three antisymmetric generators
\ensuremath{A_{1}, A_{2}, A_{3}} (the so(3) basis) and the five \ensuremath{S_{1},\dots,S_{5}} = \ensuremath{i} times the
real symmetric traceless \ensuremath{3\times 3} basis (two diagonal, three off-diagonal).
Each of the eight is anti-Hermitian (\ensuremath{G^{\dagger} = -G}) and traceless by
construction.

\textbf{H.2 Closure, by construction and by computation.} The commutator of any
two anti-Hermitian matrices is anti-Hermitian, and every commutator is traceless
identically. The real vector space of anti-Hermitian traceless \ensuremath{3\times 3}
matrices has dimension eight (three imaginary-diagonal parameters constrained to
sum to zero, giving two, plus three complex off-diagonal entries, giving six) and
is precisely su(3). The bracket of two generators therefore cannot leave this
eight-dimensional space: closure is automatic, with no overshoot. Computing all
\ensuremath{8\times 8} commutators and projecting onto the generator span confirms this
directly, with a maximum residual outside the span of \ensuremath{9\times 10^{-16}} --- zero to
machine precision. The eight generators are linearly independent (the rank of the
flattened set is eight), so none is redundant and none is missing.

\textbf{H.3 The algebra is su(3), not a same-dimension imposter.} Three
invariants fix the algebra uniquely among eight-dimensional real Lie algebras.
(i) The Killing form \ensuremath{K_{ab} = \operatorname{tr}(\operatorname{ad}_{a}\operatorname{ad}_{b})}
is negative-definite --- all eight eigenvalues equal \ensuremath{-12} --- so the algebra is
compact, and nondegenerate, so by Cartan's criterion it is semisimple. (ii) The
center of the algebra is trivial: the common kernel of all eight adjoint maps has
dimension zero, so there is no hidden central extension and the algebra is simple.
(iii) The Cartan subalgebra has dimension two (the two diagonal generators
commute), so the rank is two, with the six remaining generators supplying the six
roots of the \ensuremath{A_{2}} root system. The unique compact simple real Lie algebra of
dimension eight with rank two is su(3). A control algebra of the same dimension
and compactness, \ensuremath{su(2)\oplus su(2)\oplus u(1)^{2}}, has a two-dimensional center
and fails the simplicity test, confirming that the invariants discriminate.

\textbf{H.4 The fundamental is faithful and irreducible.} The three agents carry
the defining three-dimensional representation. By Schur's lemma the only matrices
commuting with all eight generators are scalars (the commutant has dimension
one), so the representation is irreducible; the center of the group is then the
cube roots of unity \ensuremath{\mathbb{Z}_{3}} acting by triality, and faithful action on single
substrate-real agents requires the cover, fixing the global group as SU(3) rather
than \ensuremath{SU(3)/\mathbb{Z}_{3}} (\S\,8.3). \textbf{Status: Derived.}

\section*{Acknowledgments}
\addcontentsline{toc}{section}{Acknowledgments}

The development of this volume was substantially aided by extended
collaboration with several large language models used as critical
sounding boards, pedagogical resources, formal-rigor checkers, and
external working memory across many sessions during 2024--2026.

\textbf{Gemini (Google DeepMind)} was the long-running primary
collaborator across the framework's conceptual development, providing
physics pedagogy, structural critique, and ongoing review of the ideas
that became Volume 1. The author's path into the technical literature,
and the working-out of the framework's core conceptual moves over a long
period preceding formalization, were carried by extended dialogue with
Gemini.

\textbf{Claude (Anthropic)} was the primary collaborator across the
formalization, derivation tightening, simulation work, citation
auditing, and editorial passes that brought the volume to its present
form, including the sublattice bifurcation analysis empirically verified
in Paper 5 Appendix E, the reconstruction of the matter sector (the
spin-geometry verification and minimal-valence corollary of \S\,4.1, the
coarse-graining fixed point of \S\,4.7, and the periodic table of
\S\,8.4.7), the closed-form FCC dispersion relation and its
Lorentz-violation prediction (\S\,4.1), the explicit su(3) closure
(Appendix~H), and the formalization of the Principle of Generative
Economy as a moduli-dimension selection functional.

\textbf{ChatGPT (OpenAI)}, \textbf{DeepSeek}, \textbf{Grok (xAI)}, and
\textbf{Granite (IBM)} contributed additional structural feedback and
review at various points, including a round of adversarial hostile-referee
review whose objections --- on the continuum limit and Lorentz invariance,
the uniqueness of the lattice selection, and the rigor of the gauge-sector
derivations --- directly shaped the present revision. Where those
objections identified real gaps, the gaps were addressed in the text;
where they rested on misreadings, the analysis was strengthened to
foreclose them.

All theoretical commitments, the choice of axioms, the Principle of
Generative Economy, the sublattice bifurcation insight, and the
framework's overall architecture are the author's; the AI systems' role
was to teach, formalize, critique, verify, and document.

\section*{References}
\addcontentsline{toc}{section}{References}
\begin{reflist}
Harlow, D., Usatyuk, M., \& Zhao, Y. (2025). Quantum mechanics and
observers for gravity in a closed universe. \emph{arXiv:2501.02359
[hep-th]}; published in \emph{Journal of High Energy Physics} 02
(2026) 108.

Sakharov, A. D. (1967). Violation of CP invariance, C asymmetry, and
baryon asymmetry of the universe. \emph{Pis'ma Zh. Eksp. Teor. Fiz.} 5,
32--35 [\emph{JETP Letters} 5, 24--27].

Aitchison, C. T., \& Béri, B. (2025). Spacetime Spins: Statistical
mechanics for error correction with stabilizer circuits. arXiv preprint.
https://arxiv.org/abs/2512.21991

Amelino-Camelia, G. (1998). Testable scenario for relativity with
minimum length. Physics Letters B, 510(1--4), 255--263; and
Amelino-Camelia, G. (2013). Quantum-spacetime phenomenology. Living
Reviews in Relativity, 16, 5.

Albanese, L., Barra, A., Bianco, P., Durante, F., \& Pallara, D. (2024).
Hebbian Learning from First Principles. Journal of Mathematical Physics,
65, 113302. doi:10.1063/5.0197652.

Albantakis, L., Barbosa, L. S., Findlay, G., Grasso, M., Haun, A. M.,
Marshall, W., Mayner, W. G. P., et al.~(2023). Integrated information
theory (IIT) 4.0: Formulating the properties of phenomenal existence in
physical terms. PLoS Computational Biology, 19(10), e1011465.

Bianconi, G. (2025). Gravity from entropy. Physical Review D, 111,
066001. doi:10.1103/PhysRevD.111.066001.

Bisio, A., D'Ariano, G. M., \& Tosini, A. (2013). Dirac quantum cellular
automaton in one dimension: Zitterbewegung and scattering from
potential. Physical Review A, 88, 032301. doi:10.1103/PhysRevA.88.032301.

Bisio, A., et al. (2025). Bounds on QCA Lattice Spacing from Data on
Lorentz Violation. arXiv preprint. https://arxiv.org/abs/2506.20136

Bombelli, L., Lee, J., Meyer, D., \& Sorkin, R. D. (1987). Space-time as
a causal set. Physical Review Letters, 59(5), 521--524.

Born, M. (1926). Zur Quantenmechanik der Stoßvorgänge. Zeitschrift für
Physik, 37(12), 863--867.

Brun, T. A., \& Mlodinow, L. (2025). Quantum Electrodynamics from
Quantum Cellular Automata, and the Tension Between Symmetry, Locality,
and Positive Energy. Entropy, 27(5), 492. doi:10.3390/e27050492.
https://arxiv.org/abs/2503.05998

D'Ariano, G. M., \& Perinotti, P. (2014). Derivation of the Dirac
equation from principles of information processing. Physical Review A,
90(6), 062106.

Dowker, F., \& Zalel, S. (2017). Evolution of universes in causal set
cosmology. Comptes Rendus Physique, 18(3-4), 246--253.
https://arxiv.org/abs/1703.07556

Eon, N., Di Molfetta, G., Magnifico, G., \& Arrighi, P. (2023). A
relativistic discrete spacetime formulation of 3+1 QED. Quantum, 7,
1179. doi:10.22331/q-2023-11-08-1179.

Giné, J. (2026). On the Possibility of Quantum Gravity Emerging from
Geometry. arXiv preprint. https://arxiv.org/abs/2602.16219

Gorard, J. (2020). Some Quantum Mechanical Properties of the Wolfram
Model. Complex Systems, 29(2), 537--598.

Gorard, J. (2020). Some Relativistic and Gravitational Properties of the
Wolfram Model. Complex Systems, 29(2), 599--654.

Gottesman, D. (1997). Stabilizer Codes and Quantum Error Correction.
Ph.D.~Thesis, California Institute of Technology.

Hamilton, W. R. (1834). On a General Method in Dynamics. Philosophical
Transactions of the Royal Society, 124, 247--308.

Hebb, D. O. (1949). The Organization of Behavior: A Neuropsychological
Theory. Wiley.

Hoffman, D. D., Prakash, C., \& Prentner, R. (2023). Fusions of
Consciousness. Entropy, 25(1), 129.

Hossenfelder, S. (2013). Minimal Length Scale Scenarios for Quantum
Gravity. Living Reviews in Relativity, 16(1), 2.

Jacobson, T. (1995). Thermodynamics of spacetime: The Einstein equation
of state. Physical Review Letters, 75(7), 1260--1263.

Jacobson, T., Liberati, S., \& Mattingly, D. (2002). TeV astrophysics
constraints on Planck scale Lorentz violation. Physical Review D, 66,
081302; and (2006). Lorentz violation at high energy: Concepts, phenomena
and astrophysical constraints. Annals of Physics, 321(1), 150--196.

Jaynes, E. T. (1957). Information Theory and Statistical Mechanics.
Physical Review, 106(4), 620--630.

Kauffman, S. (1993). The Origins of Order. Oxford University Press.

Klechkowski, V. M. (1962). Distribution of atomic electrons and the rule
of successive filling of \ensuremath{(n+\ell)}-groups. Zhurnal Eksperimental'noi i
Teoreticheskoi Fiziki.

Madelung, E. (1936). Mathematische Hilfsmittel des Physikers, 3rd ed.
Springer, Berlin (statement of the \ensuremath{(n+\ell)} filling rule).

Kolmogorov, A. N. (1965). Three approaches to the quantitative
definition of information. Problems of Information Transmission, 1(1),
1--7.

Mlodinow, L., \& Brun, T. A. (2021). Fermionic and bosonic quantum field
theories from quantum cellular automata in three spatial dimensions.
Physical Review A, 103, 052203.

Padmanabhan, T. (2010). Thermodynamical aspects of gravity: New
insights. Reports on Progress in Physics, 73, 046901.

Pastawski, F., Yoshida, B., Harlow, D., \& Preskill, J. (2015).
Holographic quantum error-correcting codes: toy models for the
bulk/boundary correspondence. Journal of High Energy Physics, 2015(6),
149.

Penrose, R. (1979). Singularities and time-asymmetry. In General
Relativity: An Einstein Centenary Survey, S. W. Hawking \& W. Israel,
eds., Cambridge University Press, pp.~581--638.

Penrose, R. (1971). Angular momentum: an approach to combinatorial
space-time. In Quantum Theory and Beyond, T. Bastin, ed., Cambridge
University Press, pp.~151--180.

Moussouris, J. P. (1983). Quantum models of space-time based on
recoupling theory. D.Phil.\ thesis, University of Oxford.

Szabados, L. B. (2022). A Note on Penrose's Spin-Geometry Theorem and the
Geometry of `Empirical Quantum Angles'. Foundations of Physics, 52(5),
102.

Rissanen, J. (1978). Modeling by shortest data description. Automatica,
14(5), 465--471.

Rideout, D. P., \& Sorkin, R. D. (1999). Classical sequential growth
dynamics for causal sets. Physical Review D, 61, 024002.

Rosas, F. E., Mediano, P. A. M., Luppi, A. I., et al.~(2024). Software
in the natural world: A computational approach to hierarchical
emergence. arXiv preprint. https://arxiv.org/abs/2402.09090

Shannon, C. E. (1948). A mathematical theory of communication. Bell
System Technical Journal, 27(3), 379--423.

Smolin, L. (1992). Did the universe evolve? Classical and Quantum
Gravity, 9, 173--191.

Sorkin, R. D. (2000). Indications of causal set cosmology. International
Journal of Theoretical Physics, 39, 1731.

Surya, S. (2019). The causal set approach to quantum gravity. Living
Reviews in Relativity, 22(1), 5.

Switzer, A. (2026a). The Triadic Measurement Substrate (TMS), Volume 1,
Paper 1: Emergent 3D Information Topology and the Binary Alphabet via
[[4,2,2]] Quantum Error Detection Structural Correspondence.

Switzer, A. (2026b). The Triadic Measurement Substrate (TMS), Volume 1,
Paper 2: Computational Dynamics of the Binary Alphabet.

Switzer, A. (2026c). The Triadic Measurement Substrate (TMS), Volume 1,
Paper 3: Computation Emerges.

Switzer, A. (2026d). The Triadic Measurement Substrate (TMS), Volume 1,
Paper 4: Hierarchies, Self-Modeling Measurement Nodes, Multi-Subsystem
Dynamics, and Novelty Exhaustion.

Takayanagi, T. (2025). Emergent Holographic Spacetime from Quantum
Information. Physical Review Letters, 134, 240001.
doi:10.1103/pg4r-fy8n.

't Hooft, G. (2016). The Cellular Automaton Interpretation of Quantum
Mechanics. Fundamental Theories of Physics 185. Springer.

Verlinde, E. P. (2011). On the origin of gravity and the laws of Newton.
Journal of High Energy Physics, 2011(4), 29.

Volovik, G. E. (2007). Fermi-point scenario for emergent gravity.
Proceedings of Science (QG-Ph)043. arXiv:0709.1258.

Wheeler, J. A. (1990). Information, physics, quantum: The search for
links. In W. H. Zurek (Ed.), Complexity, Entropy and the Physics of
Information (pp.~3--28). Addison-Wesley.

Wolfram, S. (2002). A New Kind of Science. Wolfram Media.

Wolfram, S. (2020). A Class of Models with the Potential to Represent
Fundamental Physics. Complex Systems, 29(2), 107--536.

Zhao, E. Y., Harlow, D., \& Usatyuk, M. (2025). Observers in a closed
universe. arXiv preprint.
\end{reflist}

\clearpage

\cleardoublepage
\thispagestyle{empty}
\vspace*{0pt}\vfill
\begin{center}
{\sffamily\bfseries\Huge Volume 2\par}
\vspace{1cm}
{\large\itshape The Confirmation Relation: From the Now-Surface to the Reader in the Frame\par}
\end{center}
\vfill\cleardoublepage



\end{document}
